\documentclass[JEP,XML,SOM,Unicode,published]{cedram}
\datereceived{2025-10-29}
\dateaccepted{2026-07-20}
\dateepreuves{2026-08-13}

\usepackage[scr=rsfs,cal=euler]{mathalfa}
\multlinegap0pt
\hyphenation{adjust-ing asso-cia-ted assump-tion assump-tions ener-gy esti-mate esti-mates instance obser-va-bi-lity obtain regard-ing result result}

\AtBeginDocument{\renewcommand{\labelitemi}{$\scriptscriptstyle\bullet$}}

\newcommand{\Psfrac}[2]{(\sfrac{#1}{#2})}
\newcommand{\psfrac}[2]{\sfrac{(#1)}{#2}}
\newcommand{\spfrac}[2]{\sfrac{#1}{(#2)}}

\newcommand{\quand}{\quad\text{and}\quad}
\newcommand\mto{\mathchoice{\longmapsto}{\mapsto}{\mapsto}{\mapsto}}
\newcommand\bnor[2]{\bigl\|#1\bigr\|_{#2}}%la norme

\newenvironment{Smallpmatrix}{\Big(\begin{smallmatrix}}{\end{smallmatrix}\Big)}

\newcommand{\lin}{\textup{lin}}
\newcommand{\Nlin}{\textup{Nlin}}

\usepackage{mathtools}

\usepackage[nameinlink,capitalise]{cleveref}
\crefdefaultlabelformat{#2\textup{#1}#3}
\makeatletter
\def\cref@thmoptarg[#1]#2#3#4{%
\ifhmode\unskip\unskip\par\fi%
\normalfont%
\trivlist%
\let\thmheadnl\relax%
\let\thm@swap\@gobble%
\thm@notefont{\fontseries\mddefault\upshape}%
\thm@headpunct{.}% add period after heading
\thm@headsep 5\p@ plus\p@ minus\p@\relax%
\thm@space@setup%
#2% style overrides
\@topsep \thm@preskip % used by thm head
\@topsepadd \thm@postskip % used by \@endparenv
\def\@tempa{#3}\ifx\@empty\@tempa%
\def\@tempa{\@oparg{\@begintheorem{#4}{}}[]}%
\else%
\refstepcounter[#1]{#3}% <<< cleveref modification
\@namedef{cref@#3@alias}{#1}% added
\def\@tempa{\@oparg{\@begintheorem{#4}{\csname the#3\endcsname}}[]}%
\fi%
\@tempa}%
\makeatother

\usepackage{stmaryrd}
\usepackage{bbm}
\usepackage{mathtools}
\DeclarePairedDelimiter\ceil{\lceil}{\rceil}
\DeclarePairedDelimiter\floor{\lfloor}{\rfloor}

\DeclareMathOperator{\supp}{\mathrm{supp}}

\newtheorem{theorem}{Theorem}[section]
\newtheorem{lemma}[theorem]{Lemma}
\newtheorem{proposition}[theorem]{Proposition}
\newtheorem{corollary}[theorem]{Corollary}

\theoremstyle{definition}
\newtheorem{definition}[theorem]{Definition}
\newtheorem{example}[theorem]{Example}
\newtheorem{remark}[theorem]{Remark}
\newtheorem{assum}{Assumption}
\newtheorem*{claim}{Claim}

\newenvironment{assump}[2][]
{\renewcommand\theassum{#2}\begin{assum}[#1]}
{\end{assum}}

\let\mc\mathcal
\let\co\overline
\let\Ld\Lambda
\let\al\alpha
\let\ld\lambda
\let\vp\varphi
\let\veps\varepsilon
\let\wto\rightharpoonup

\newcommand{\A}{\mc{A}}
\newcommand{\B}{\mc{B}}
\newcommand{\E}{\mc{E}}
\newcommand{\J}{\mc{J}}
\newcommand{\M}{\mc{M}}
\newcommand{\T}{\mc{T}}
\renewcommand{\O}{\mc{O}}
\newcommand{\X}{\mc{X}}
\newcommand{\Y}{\mc{Y}}
\newcommand{\C}{\boldsymbol{C}}
\newcommand{\R}{\mathbb{R}}
\newcommand{\N}{\mathbb{N}}

\newcommand{\Q}{\mathcal{Q}}
\renewcommand{\P}{\mathcal{P}}
\newcommand{\bC}{\boldsymbol{C}}
\newcommand{\bA}{\boldsymbol{A}}
\newcommand{\weto}{\rightharpoonup^*}
\providecommand{\norm}[1]{\lVert#1\rVert}
\newcommand{\Norm}[1]{\left\lVert#1\right\rVert}
\providecommand{\inn}[1]{\langle#1\rangle}
\providecommand{\Int}[1]{\textup{Int}(#1)}
\newcommand\nor[2]{\left\|#1\right\|_{#2}}
\newcommand\FL{\mathcal{F}_{n}}

\datepublished{2026-08-27}
\begin{document}

\title[Propagation results for semilinear waves]{Global propagation of analyticity and unique\nobreakspace continuation for semilinear waves}

\author[\initial{C.} \lastname{Laurent}]{\firstname{Camille} \lastname{Laurent}}
\address{CNRS UMR 9008, Université Reims-Champagne-Ardennes, Laboratoire de Mathématiques de Reims (LMR),\\
Moulin de la Housse-BP 1039, 51687 Reims cedex 2, France}
\email{camille.laurent@univ-reims.fr}
\urladdr{https://laurentca.perso.math.cnrs.fr/}

\author[\initial{C.} \lastname{Loyola}]{\firstname{Cristóbal} \lastname{Loyola}}
\address{Sorbonne Université, Université Paris Cité, CNRS, Laboratoire Jacques-Louis Lions, LJLL,\\
F-75005 Paris, France}
\curraddr{\& Université Marie et Louis Pasteur, CNRS, Institut UTINAM, Équipe de physique théorique,\\F-25000 Besançon, France}
\email{cristobal.loyola@sorbonne-universite.fr}
\urladdr{https://cloybar.github.io/}

\thanks{The second author has received funding from the European Union's Horizon 2020 research and innovation programme under the Marie Sk\l{}odowska-Curie grant agreement No\,945332}
\CDRGrant[EU Horizon 2020]{Marie Sk\l{}odowska-Curie grant agreement No 945332}

\begin{abstract}
In this article, we~develop a new method to prove both global propagation of analyticity and unique continuation in finite time for solutions of semilinear wave-type equations with analytic nonlinearity. It~combines control theory techniques and Galerkin approximation, inspired by Hale-Raugel \cite{HR03}, to prove that analyticity in time can be propagated for the nonlinear equation from a zone where linear observability holds towards the full space. For semilinear wave equations on a bounded domain, this implies that analyticity can be propagated to the entire domain from a subset $\omega$ that satisfies the geometric control condition. It~also implies the unique continuation when the solution is assumed to be zero on $\omega$. When the nonlinearity is assumed to be subcritical and defocusing, we~also obtain observability estimates in the optimal time of the geometric control condition. For semilinear plate equations, similar propagation of analyticity is achieved by assuming the controllability of the linear Schrödinger equation.
\end{abstract}

\subjclass[2020]{35A20, 35B60, 93B07, 35L71, 35L75} %

\keywords{Propagation of analyticity, unique continuation, semilinear wave equation, semilinear plate equation}

\altkeywords{Propagation de l'analyticité, prolongement unique, équation des ondes semi-linéaire, équation des plaques semi-linéaire}

\alttitle{Propagation globale de l'analyticité et prolongement unique pour des ondes semi-linéaires}

\begin{altabstract}
Dans cet article, nous développons une nouvelle méthode pour établir à la fois la propagation globale de l'analyticité et le prolongement unique en temps fini pour les solutions d'équations semi-linéaires de type onde à non-linéarité analytique. Combinant des techniques de théorie du contrôle et une approximation de Galerkin, dans l'esprit de Hale-Raugel \cite{HR03}, elle permet de propager l'analyticité en temps de l'équation non linéaire, depuis une zone où l'observabilité linéaire est vérifiée, vers l'espace tout entier. Pour les équations des ondes semi-linéaires sur un domaine borné, cela implique que l'analyticité peut être propagée au domaine tout entier à partir d'un sous-ensemble $\omega$ satisfaisant la condition de contrôle géométrique. Cela implique également le prolongement unique lorsque la solution est supposée nulle sur $\omega$. Lorsque la non-linéarité est supposée sous-critique et défocalisante, nous obtenons aussi des estimations d'observabilité dans le temps optimal de la condition de contrôle géométrique. Pour les équations des plaques semi-linéaires, un résultat analogue de propagation de l'analyticité est obtenu en supposant la contrôlabilité de l'équation de Schrödinger linéaire.
\end{altabstract}

\maketitle
\vspace*{-\baselineskip}
\tableofcontents

\mainmatter

\section{Introduction and main results}
The aim of this article is to study how, for a certain class of evolution PDEs, properties observed from a subset $\omega\subset \M$ through time $(0, T)$ are propagated to the whole solution on $(0, T)\times \M$. Here, $\M$ can be, for instance, a compact Riemannian manifold with boundary. More precisely, we~will study the following three properties:
\begin{enumerate}
\item \emph{Propagation of analyticity}: if the solution is analytic in time on $(0,T)\times \omega$, is the full solution analytic in time on $(0,T)\times\M$?
\item \emph{Unique continuation}: if the solution is zero on $[0, T]\times \omega$, is the solution identically zero in $[0, T]\times\M$?
\item \emph{Observability inequality}: can we quantify the previous unique continuation property?
\end{enumerate}
The method that we develop is quite general for conservative equations and relies on observability estimates. We~provide an abstract result, detailed more precisely in \cref{s:abstrIntro}, that allows to propagate analyticity in time from the observation to the full solution. Afterward, we~will give applications to nonlinear wave (see \cref{s:waveIntro}) and nonlinear plate equations (see \cref{s:plateintro}). We~believe that it could be applied to several other systems and is amenable to generalizations.

Before getting to the details of the results, let us discuss some motivations for the questions we address. Propagation of analyticity and unique continuation for nonlinear equations, in~addition to being of intrinsic interest, arise mainly in various rigidity problems. In~some situations, we~need to identify some objects that satisfy some very specific properties and are often expected to be the asymptotic solution in large time. Here are some examples
\begin{itemize}
\item \emph{Nonlinear stabilization problems:} consider some partially damped equation of the form, for instance, $\partial_t^2u-\Delta u+\mathbbm{1}_{\omega}\partial_t u=f(u)$ or $\partial_t^2u+\Delta^2 u+\mathbbm{1}_{\omega}\partial_t u=f(u)$ with $\omega \subset \M$. A main obstacle to the decay would be the existence of undamped solutions that would satisfy $\partial_t u=0$ on $\omega$. It~is very desirable to obtain such a result with some geometric assumptions on $\omega$ that are as minimal as possible. In~the case that these objects do exist, they are expected to describe the asymptotic behavior of solutions, for instance if they belong to some attractor set. We~refer, for instance, to \cite{Z:91,DLZ03,JL13,L:11} for such problems in the context of wave equations.
\item \emph{The rigidity of stationary black holes:} this problem consists in classifying the set of solutions of the Einstein equation that contain a Killing vector field, are asymptotic to the Minkowski metric, and satisfy suitable additional assumptions. It~is conjectured that there are no solutions of this type beyond the Kerr family. This subject has recently been the subject of intense research (see \cite{IK:15} for a survey). The result is already known to be true for analytic solutions. In~particular, a result of propagation of analyticity for such solutions would certainly allow some great progress on this question. Although our paper is restricted to semilinear equations, we~might expect that our result, along with additional ideas, could extend to the fully nonlinear setting.
\item \emph{Nonlinear scattering theory:} on $\R^d\setminus \mathcal{O}$, where $\mathcal{O}$ is a finite union of compact obstacles, consider, for instance, equations $\partial_t^2u-\Delta u=f(u)$ or $\partial_t^2u+\Delta^2 u=f(u)$. The question of scattering asks whether there exists a linear solution $u_L$ such that $u(t)-u_L(t)\to 0$ as $|t|\to+\infty$ in some suitable norm. There might be several obstacles to this property. One is the existence of solutions with the compactness property, which means that $\{u(t)\ |\ t\in\R\}$ is compact. An even worse scenario is the existence of a solution that remains supported in a fixed compact set in space. This is the typical situation of unique continuation that our analysis might cover. Note that the current methods for this problem use Morawetz-type estimates (see, for instance, \cite{BSS:09} outside of a star-shaped obstacle or \cite{KVZ:16} for NLS outside of a convex obstacle), but it seems that they cannot reach the more general geometries we wish to treat.
\end{itemize}

To motivate the abstract result, we~begin by describing the results for nonlinear wave equations.

\subsection{Main results on semilinear wave equation}\label{s:waveIntro} In this section we consider the semilinear wave equation
\begin{align}\label{eq:nlw-1}
\left\{\begin{aligned}
\partial_t^2 u-\Delta_g u+f(u)&=0 &&(t, x)\in [0, T]\times \text{Int}(\M),\\
u_{|_{\partial\M}}&=0 &&(t, x)\in [0, T]\times\partial\M,\\
(u, \partial_t u)(0)&=(u_0, u_1) && x\in \M,
\end{aligned}\right.
\end{align}
where $(\M, g)$ is a smooth compact connected Riemannian manifold with boundary of dimension $d$ less or equal to $3$ and $\Delta_g$ is the Laplace-Beltrami operator equipped with Dirichlet boundary conditions. The nonlinearity $f: \R\to \R$ is assumed to be \emph{analytic} and to satisfy $f(0)=0$.

We will make two kinds of assumptions concerning the regularity and the nonlinearity:
\begin{enumerate}
\item if there is no more assumption, we~will choose the regularity index $\sigma\in (1/2,1]$ if $d=3$ and $\sigma\in (0,1/2)$ if $d\leq 2$. This ensures $H^{1+\sigma}\hookrightarrow L^{\infty}$.
\item $f$ is energy subcritical: $f$ is assumed to be of polynomial type, in~the sense that there exists $C>0$ such that
\begin{align}\label{hip:nonlinearity-hyp-1}
\ |f(s)|\leq C(1+|s|)^p\ \text{and}\ |f'(s)|\leq C(1+|s|)^{p-1},
\end{align}
with $1\leq p<+\infty$ if $d\leq 2$ and $1\leq p<5$ for $d=3$. In~that case, we~choose $\sigma=0$ in what follows, corresponding to finite energy solutions.
\end{enumerate}

The observation set $\omega\subset \M$ will also always be assumed to be open. A key hypothesis will be that $\omega$ satisfies the Geometric Control Condition at time $T>0$:
\begin{assump}{GCC}\label{assumGCC}
Every generalized geodesic of $\M$ traveling at speed $1$ meets $\omega$ in time $t\in (0, T)$.
\end{assump}
We will always assume that the Hamiltonian vector field of the wave operator does not have contact of infinite order with $\partial \M$ (see \cite{MS:78}). This ensures that the broken bicharacteristic flow is uniquely defined. We~refer to \cref{s:geom} for details and references.

Our first main result concerns the propagation of analyticity in time.

\begin{theorem}\label{thm:analytic-prop}
Let $\sigma\in (1/2, 1]$ if $d=3$ and $\sigma\in (0,1/2)$ if $d\leq 2$. Let $(u,\partial_t u)\in C^0([0, T], H^{1+\sigma}\cap H^1_0\times H^{\sigma}_0(\M))$ be a solution of \eqref{eq:nlw-1}. Assume that the above setting holds and assume moreover that:
\begin{enumerate}
\item $(\omega, T)$ satisfies the \ref{assumGCC}.
\item $t\in (0, T)\mto \chi u(t,\cdot)\in H^{1+\sigma}(\M)\cap H_0^1(\M)$ is analytic for any cutoff function $\chi\in C_c^\infty(\M)$ whose support is contained in $\omega$.
\item $s\mto f( s)$ is real analytic.
\end{enumerate}
Then $t\mto u(t, \cdot)$ is analytic from $(0, T)$ to $H^{2}(\M)\cap H_0^1(\M)$.
\end{theorem}

Note that, once the solution is known to be analytic in time, it~is possible to get more regularity in space. For instance, this occurs when the metric is analytic.

\begin{corollary}
\label{coranalytspacetime}
Under the same assumptions as before, if, moreover, the metric and the boundary are analytic in $x$, then $u$ is analytic in all variables.
\end{corollary}

The second assumption is written as an analyticity property with a Banach valued map, see \cite[Def.\,5.1]{Muj86}. However, it~can be satisfied if we have, for instance, the pointwise estimates $|\partial_t^{\alpha}\partial_x^{\beta}u(t,x)|\leq CR^{|\alpha|}\alpha!$ for $(t,x)\in (0,T)\times \omega$, $\alpha\in \N$, $|\beta|\leq 2$, together with the Dirichlet boundary condition.

These propagation results have applications to unique continuation problems. In~what follows, we~will consider solutions $(u,\partial_t u)\in C^0([0, T], H^1_0\times L^2(\M))$ of \eqref{eq:nlw-1} and say that it has finite Strichartz norm if one of the admissible Strichartz exponent ensuring uniqueness in the energy space is finite (see \cref{thmStrichartz}). The shape of the nonlinearity determines what type of equilibrium will be obtained.

\begin{theorem}[Unique continuation]\label{thm:unique-continuation-nlw}
Assume that $(\omega, T)$ satisfies the \ref{assumGCC} and that~$f$ is energy subcritical and real analytic. If one solution
\[
(u,\partial_t u)\in C^0([0, T], H^1_0\times L^2(\M))
\]
with finite Strichartz norms of \eqref{eq:nlw-1} satisfies $\partial_t u=0$ in $[0, T]\times \omega$, then $\partial_t u=0$ in $[0, T]\times \M$ and $u$ is an equilibrium point of \eqref{eq:nlw-1}, that is, solution of
\begin{align}\label{thm:eq:nlw-equilibrium}
\left\{\begin{aligned}
-\Delta_g u+f(u)&=0 &x&\in\textnormal{Int}(\M),\\
u&=0 &x&\in \partial\M.
\end{aligned}\right.
\end{align}
If, moreover, the nonlinearity satisfies
\begin{align}
\label{defdefocusing}
sf(s)&\geq 0\quad\quad \textnormal{ if }\partial \M\neq \emptyset,\\
\nonumber sf(s)&\geq \gamma s^2\quad \textnormal{ if }\partial \M= \emptyset,
\end{align}
for some $\gamma>0$ and for all $s\in \R$, then $u\equiv 0$.
\end{theorem}

Note that when unique continuation is not possible, we~can still establish a result of finite determining modes, see \cref{propfinitedetwave}. In~the first part of the theorem, $f$~can be focusing. Therefore, \eqref{eq:nlw-1} is not always globally well-posed and can lead to blow-up, that is why we make the a priori assumption of boundedness. In~the case of defocusing nonlinearity, we~are able to prove a quantitative version of the unique continuation, namely, an observability inequality.

\begin{theorem}
\label{thmobserintro}
Assume that $(\omega, T)$ satisfies the \ref{assumGCC}. Assume that $f$ is analytic, energy subcritical and defocusing, that is, satisfying \eqref{hip:nonlinearity-hyp-1}, and \eqref{defdefocusing}. Then for any $R_0\!>\!0$, there exists $C\!>\!0$ such that for any $(u_0,u_1)\!\in\! H^1_0\!\times\! L^2(\M)$, with $\nor{(u_0,u_1)}{H^1_0\times L^2}\leq R_0$, the unique solution of \eqref{eq:nlw-1} with finite Strichartz norms satisfies
\begin{align}\label{obsevNLintro}
C\norm{(u_0, u_1)}_{H_0^1(\M)\times L^2(\M)}^2\leq \int_0^T \norm{\mathbbm{1}_\omega\partial_t u(t)}_{L^2(\M)}^2 dt.
\end{align}
\end{theorem}
For linear equations, this is the typical observability estimate of \cite{BLR92} that is equivalent to the controllability. Here, however, it~is obtained for nonlinear equations. This type of inequality already appeared in the literature, often associated with stabilization results, but under stronger assumptions on $\omega$ and $T$ to ensure unique continuation. Indeed, the stabilization property for the damped equation is often proved thanks to an observability inequality as \eqref{obsevNLintro}. This topic has been studied extensively; for instance, see \cite{H:85}, \cite{Z:91} and \cite{D:01} for $p<3$, and for $p\in[3,5)$, our main reference is the work of Dehman, Lebeau and Zuazua \cite{DLZ03}. This work mainly addressed the stabilization problem previously described, in~the Euclidean space $\R^3$ with flat metric and active damping outside of a ball, or on a bounded domain, but with damping close to the full boundary. It~would be interesting to know how the constant $C$ in~\eqref{obsevNLintro} depends on $R_0$. Under stronger assumptions of multiplier type, it~was made explicit in Duyckaerts-Zhang-Zuazua~\cite{DZZ:08} for linear equations. In~the linear case with time independent potential under \ref{assumGCC}, this question has already been studied in Laurent-Léautaud~\cite{LL16} with an estimate of $C$ depending either on the time or on the size of the potential. Obtaining such dependence in the nonlinear setting might be possible but it would require, at least, to make quantitative every step of the proof of \cref{thmobserintro} and also allow analytic in time potentials in \cite{LL16}. It~would certainly be worthwhile but is out of the scope of this article.

Our main purpose here is to extend observability to more general geometries where multiplier methods cannot be used or do not yield optimal results with respect to the geometry and optimal time. Other stabilization results for the nonlinear wave equation can be found in \cite{AIN:11} or \cite{Per26}, as well as the references therein. Some articles have addressed the more difficult critical case $p=5$; see \cite{L:11} for details. An observability estimate similar to \eqref{obsevNLintro} was also obtained by Joly and the first author in \cite{JL13} under the same Geometric Control Condition, but with a non-uniform time depending on the size of the initial data $R_0$. Their proof relied on some asymptotic regularization result of Hale-Raugel \cite{HR03}. Our scheme of proof here is inspired by their work, with the substantial modification that we always work in finite time, which requires adapting the original method.

We should also note that, although the results of \cite{DLZ03} concern a specific geometry, by~following their proof carefully, we~can extract a rough statement of the type:
\begin{center}
``geometric control condition'' + ``unique continuation'' $\implies$ ``observability''.
\end{center}
Therefore, our proof will follow that line and will mainly focus on the unique continuation property.

Note that we have assumed from the beginning that $\M$ is compact. Yet, several results in unbounded domains can be deduced from our result. We~can, for example, get similar results for a compact perturbation of $\R^3$.

\begin{proposition}
\label{propR3obstacle}
Let $\Omega=\R^3\setminus \mathcal{O}$ where $\mathcal{O}$ is a bounded smooth domain (not necessarily connected). Assume that there exists $R>0$ such that $\R^3\setminus B(0,R)\subset \omega$ and that $(\omega,T)$ satisfies \ref{assumGCC}. Then the same conclusion as \cref{thm:unique-continuation-nlw} holds.
\end{proposition}
\begin{remark}
\label{rk:higherd}
In the previous results, the restriction on the dimension is mainly technical. Similar results would likely hold in any dimension, possibly with additional assumptions. First, one needs to modify the definition of subcriticality. Secondly, in~the case with boundary, we~need to be careful about the fact that the spaces $H^{\sigma}_D$, with $\sigma \in (d/2-1,d/2-1/2)$, impose some specific boundary conditions that do not behave so well with respect to composition and multiplication by cutoff functions. Here, $H^{\sigma}_D$ is the Sobolev space adapted to the Laplace operator with Dirichlet boundary conditions. For the nonlinearity, we~would need to assume that $f^{(k)}(0)=0$ for \hbox{$0\leq k<k_d$}, with $k_d$ chosen such that $f$ sends $H^{1+\sigma}_D$ to $H^{\sigma}_D$ for $1+\sigma>d/2$.
For instance, we~could assume the following condition on $f$, which is supposed to be analytic for simplicity, so that \cref{prop:f-assumptions} holds true:
\begin{equation*}
f^{(k)}(0)=0, \quad \forall\text{ $k$ such that}~0\leq k\leq \lfloor d/2\rfloor-1.
\end{equation*}
This ensures compatibility with the boundary conditions, so that $f$ sends $H^{1+\sigma}_D$ to $H^{\sigma}_D$ for $\sigma\in (d/2-1, (d-1)/2)$. Additionally, the proof of \cref{thm:analytic-prop} involves the action of cutoff functions that require more care for the boundary calculus. That indicates that the second assumption should be replaced by $t\in (0, T)\mto \chi u(t,\cdot)\in H^{1+\sigma}_D(\M)$.
We refer to \cref{rk:proofhigherd} and \cref{rk:comm-reg} below for more details.

Nevertheless, many of these complications are caused by the compatibility conditions. For instance, \cref{thm:analytic-prop} holds for $d\geq 3$ where $\sigma\in (d/2-1, (d-1)/2)$ when $\partial\M=\emptyset$.
\end{remark}

The proof of the previous results are consequences of abstract results described in \cref{s:abstrIntro}. Our methods are inspired by Dynamical Systems techniques from Hale-Raugel \cite{HR03}, originally designed to study the regularity of ``ancient solutions'' for dissipative systems. It~was noted in \cite{JL13} that such results could imply a unique continuation property with Geometric Control Condition on $\omega$, although, in~infinite time or depending on the size of the data. To achieve finite-time results, the method has to be modified. The main idea in this article is to rely on observability properties instead of the decay of semigroup. This approach enables us to obtain some propagation results in finite time for solutions with zero observation. Additionally, the method is flexible enough to allow some source terms, leading to a full propagation result with optimal assumption on $\omega$ and the time $T$, from an observation analytic in time. For more details, we~refer to \cref{s:abstrIntro}.
\subsubsection*{Literature overview on unique continuation and propagation} Let us now review some already known results and why our unique continuation result seems difficult to obtain from the already known methods and results in the literature. We~also refer, for instance, to the survey \cite{LL:23Lecture} and the reference therein that contains an introduction to the unique continuation for wave type operators. For a more complete treatment of Carleman estimates and unique continuation, we~also refer to \cite{LernerBook19}.

One initial approach for proving \cref{thm:unique-continuation-nlw} could involve using the general theory of unique continuation of Hörmander \cite{H:63}. In~that context, this would require considering the nonlinearity $f(u)$ as a potential term $Vu$ where $V$ has, at most, the same regularity as $u$. Achieving global unique continuation would then require iterating some local unique continuation results across a hypersurface $\{\Psi=0\}$. Yet, for a general configuration, the global geometric assumptions resulting from the use of Hörmander theorem for the unique continuation are not very natural and are stronger than \ref{assumGCC}. For instance, for a flat metric, the pseudoconvexity condition of a hypersurface $\{\Psi=0\}$ for the wave operator writes: if $X_{t}^{2}=|X_{x}|^{2}$ and $d\Psi(x_0)(X)=0$, then
\begin{align*}
\text{Hess} \Psi(x_0) (X , X) >0 \quad \text{for all } X=(X_{t},X_{x}) \in \R^{1+d} \setminus \{0\}.
\end{align*}
and imply a kind of convexity of the hypersurface. Typical geometric assumptions are often of ''multiplier type'' (or Morawetz type), meaning $\omega$ is a neighborhood of $\left\{x\in \partial\Omega \left|(x-x_0)\cdot n(x)>0\right.\right\}$ which are known to be stronger than the geometric control condition (see \cite{M:03} for a discussion about the links between these assumptions). Moreover, on curved spaces, this type of condition often needs to be checked by hand in each situation, which is mostly impossible in general. Additionally, concerning the unique continuation with even smooth $V$, the classical counterexamples of Alinhac-Baouendi~\cite{AB:95}, refined by Hörmander \cite{H:00}, are quite striking. They show that Hörmander's pseudoconvexity condition is not far from being optimal for local unique continuation. For any $s>1$ and $d\geq 2$, they construct some $u$ and $V \in \mathcal{G}^s(B_{\R^{1+d}}(0,1), \mathbb{C})$ (Gevrey functions) such that
\begin{align*}
\partial_{t}^{2}u-\Delta u&=V u \textnormal{ on }B_{\R^{1+d}}(0,1),\\
\supp(u)&=\left\{(t,x_{1},\dots, x_{d})\ |\ x_{1}\geq 0\right\}\cap B_{\R^{1+d}}(0,1).
\end{align*}
This suggests that in geometrical situations where the strong pseudoconvexity of the hypersurface is not satisfied, we~cannot expect local unique continuation for potential a $V$ that is not analytic.

Note that quite surprisingly, even in $1$-dimensional linear hyperbolic systems of order $1$, a similar dichotomy seems to exist. It~has been shown in this context by Coron-Nguyen \cite{CN:21} that for time-dependent coupling matrices, it~is possible to construct counterexamples to unique continuation at some natural time while the property is true for coupling depending analytically on time.

Another counterexample that undermines a possible strategy to prove unique continuation in \cref{thm:unique-continuation-nlw} is provided by Métivier in \cite{M:93}. He proved that a nonlinear version of the Holmgren theorem fails in general. The operators for which it applies are not wave operators, but a nonlinear Holmgren theorem, even for a more specific class of operators has, up to our knowledge, never been obtained so far, except for scalar operators of order $1$.

Regarding the propagation of analyticity for nonlinear equations, several results date back to the 1980s and 1990s. Alinhac-Métivier proved in \cite{AM:84,AM:84b} that if $u$ is a regular enough solution of a general nonlinear PDE, the analyticity of $u$ propagates along any hypersurface for which the real characteristics of the linearized operator cross the hypersurface transversally. Subsequently, there has been an intense activity to understand what kind of singularities propagate for nonlinear waves. It~was found that the situation is quite complicated since microlocal analytic singularities do not remain confined to bicharacteristics as in the linear case, but can give rise to nonlinear interactions. For more details, see Godin \cite{G:86} and Gérard \cite{G:88}.

Yet, in~our geometric context, obtaining a global result from local propagation of singularities typically involves propagation from hypersurfaces of the form $S=\{\psi=0\}$ with $\psi(t,x)=\psi(x)$. In~such cases, the operator is never hyperbolic with respect to $S$, and there can be some bicharacteristics transverse to $S$ as soon as $d\geq 2$. In~particular, the results from \cite{AM:84,AM:84b} do not seem to apply.

The problem is better understood in the $C^{\infty}$ or $H^s$ context. For the linear equation, the propagation of $H^s$ regularity along rays for the Dirichlet problem is well understood since the work of Melrose-Sjöstrand \cite{MS:78}; see \cite{H:07} for a complete historical overview on the internal and boundary problem. Concerning the $H^s$ regularity of nonlinear equations, the microlocal propagation has been the object of several studies since the work of Bony \cite{B:81}. The global propagation from a set satisfying \ref{assumGCC} is proved in \cite{DLZ03} using a bootstrap argument and propagation of $H^s$ regularity for smoother source terms. It~is unclear how to adapt such arguments in the analytic context. Even in the linear case $f=0$, local propagation of analytic regularity can be quite complicated, especially for glancing rays; see \cite{RS:81} for details. The propagation of analyticity that we prove may also be of interest in this context. It~seems that, to achieve global propagation of analytic regularity using the propagation of the wavefront set, one would need to take into account analytic rays. They include rays arriving at the boundary in a glancing direction and potentially staying ``stuck'' at the boundary for a certain amount of time, even at diffractive points. It~is uncertain to us whether the global propagation result can be obtained under the same assumptions using this analysis for the linear case.

Beyond Hörmander's general result \cite{H:63} under pseudoconvexity, unique continuation for wave-type operators have been extensively studied. While it is impossible to cover them all here, we~aim to present the variety of results and motivations that appear. A more precise historical overview can be found in \cite[\S 4]{LL:23Lecture}.

We begin with the unique continuation with partial analyticity, which will be crucial for obtaining \cref{thm:unique-continuation-nlw} from \cref{thm:analytic-prop}, see \cref{s:UCPword}. The history of this theory is quite long, with several breakthroughs and improvements that we do not detail here (see \cite[\S 4]{LL:23Lecture}). The proof of local uniqueness results across any non-characteristic hypersurface for $\partial_t^2 - \Delta_g$ was achieved by Tataru in \cite{Tat95}, leading to the global unique continuation result in optimal time. Tataru's result is not restricted to the wave operator;
it holds for operators with coefficients that are analytic in part of the variables, interpolating between Holmgren's theorem and the Hörmander's theorem. Technical assumptions of this article were successively removed by Robbiano-Zuily~\cite{RZ:98}, Hörmander~\cite{Hor:97} and Tataru~\cite{Tat99}, leading to a very general local unique continuation result for operators with partially analytic coefficients (including as particular cases both Holmgren's and Hörmander's theorems). This result (or~some globalized version of it) will be used in our context; see \cref{thm:qucp} below. The key point in applying these results, which are linear, is that the coefficients need to be analytic in time. When applying this to nonlinear solutions such as \eqref{eq:nlw-1}, the nonlinearity $f(u)$ must be seen as a term $Vu$ with $V$ having the same regularity as $u$. This is why the propagation of analyticity in \cref{thm:analytic-prop} is crucial in order to apply this unique continuation result.

In regards to classical Carleman estimates, many authors explored the global assumptions needed to obtain them, as well as their consequences. Global Carleman estimates for waves were proved in~\cite[Ch.\,4]{FI:96} and \cite{BDBE:13} with applications to controllability and inverse problems. Another line of investigations in a geometric context was undertaken in~\cite{DZZ:08,Shao:19,JS:24} aiming to present geometric assumptions that would ensure the usual pseudoconvexity, with applications to observability estimates and null-controllability.

For treating nonlinear problems, admitting lower-order terms in unique continuation results is crucial. For instance, in~the present work, a nonlinearity of the form $f(u)$ is treated as a term $V u$ with potential $V$. Here, a previous analysis allows to obtain that the nonlinear solution is actually very regular. Yet, having nonlinear problems in mind, it~has sometimes been a goal to minimize the regularity of the admissible lower-order terms. A global unique continuation statement was proved in~\cite{Ruiz:92}, with application to energy decay for nonlinear waves. This led to some ``dispersive'' Carleman estimates with Strichartz-type spaces. The literature is vast, and we refer, for instance, to~\cite{KRS:87,DDSF:05, KT:05}, and the references therein.

Unique continuation problems for wave operators also arise from mathematical general relativity. They have recently received a lot of attention, specifically in the context of the rigidity problem for stationary black holes. We~refer to \cite{IK:15} for a precise overview of the problem and recent progress. Note also that the rigidity of stationary black holes is already known under the assumption of analyticity. So, obtaining some propagation of analyticity could be of great interest in this context.

\subsection{Main results on semilinear plate equations}\label{s:plateintro}

As mentioned earlier, the results we obtain for the wave follow from more a general method and abstract framework that could apply to many other systems. Firstly, the result could be extended to systems of waves with the same leading order terms but coupled by lower order terms, provided that the observation is made across all components. Additionally, we~present a first application to nonlinear plate equations, focusing solely on the propagation of analyticity. However, this could represent a first step towards unique continuation in a more general setting.

\begin{theorem}\label{thm:analytic-nlp}
Let $(\M, g)$ be a compact connected manifold with (or without) boundary of dimension $d\leq 3$. Let $\widetilde{T},T$ be such that $0<\widetilde{T}<T$ and let $\omega\Subset \widetilde{\omega}$ be two open subsets of $\M$ such that the Schrödinger equation is observable in $L^2$ from $\omega$ in time~$\widetilde{T}$ (see \cref{sssec:obs-schr} below for more precisions and examples). Assume that $f$ is real analytic with $f(0)=0$. Let $(u, \partial_t u)\in C^0([0, T], H^2\cap H_0^1\times L^2)$ be a solution of
\begin{align}\label{eq:nlp-1intro}
\left\{\begin{aligned}
\partial_t^2 u+\Delta^2 u+f(u)&=0 &(t, x)&\in [0, T]\times \M,\\
u_{|_{\partial\M}}=\Delta u_{|_{\partial\M}}&=0 &(t, x)&\in [0, T]\times\partial\M,
\end{aligned}\right.
\end{align}
such that, for any cutoff function $\chi\in C_c^\infty(\M)$ whose support is contained in $\widetilde{\omega}$,
\[
t\in (0, T)\mto \chi u(t,\cdot)\in H^{3+\veps}(\M)\cap H_0^1(\M)
\]
is analytic for one $\veps>0$. Then
\[
t\in (0, T)\mto \big(u(t, \cdot), \partial_t u(t,\cdot)\big)\in H^4\cap H_0^1\times H^2\cap H_0^1
\]
is analytic.
\end{theorem}

Note that the condition $\Delta u_{|_{\partial\M}}=0$ which does not make sense at this level of regularity is meant as an extension of the related semigroup. The nonlinear equation~\eqref{eq:nlp-1intro} is meant in the sense of the Duhamel formula.

The sharp geometric condition on $\omega$ necessary for the observability of the Schrödinger equation remains an open question. Yet, it~has been the object of many investigations. In~the following situations, the observability of the Schrödinger equation is known, thereby allowing us to apply the previous theorem:
\begin{enumerate}
\item $(\M, g)$ is a compact Riemannian manifold with or without boundary and $\omega$ satisfies the \ref{assumGCC}. See Lebeau \cite{Leb92}.
\item $(\M,g)=((0,1)^d,\text{Euclid})$ with $d\leq 3$, $\omega$ is any nonempty open set. This was first proved by Jaffard for $d=2$ and Komornik \cite{K:92} for other dimensions (actually directly for the beam equation). Other proofs have also been given later by Burq-Zworski \cite{BZ:04}, Anantharaman-Macià \cite{AM:14}. Note here that the proofs are given for the torus $\mathbb{T}^d$, but an easy argument of symmetrization allows to recover the same result for the Dirichlet boundary condition.
\item $(\M, g)=(\mathbb{D}, \text{Euclid})$ is the Euclidean closed disk in $\R^2$ and $\omega\cap \partial \M\neq\emptyset$. See Anantharaman, Léautaud, Macià \cite[Th.\,1.2]{2016:anantharaman-leautaud-macia:wigner-schrodinger-disk},
\item $(\M, g)$ is a compact, boundaryless connected Riemannian surface whose flow has the Anosov property, $\omega$ is any nonempty open set. See Dyatlov-Jin-Nonnenmacher \cite[Th.\,5]{DJN:22}.
\item $(\M, g)$ is the Bunimovich stadium and $\omega$ controls geometrically $\M\setminus R$, $R$ being the rectangular part, see Burq-Zworski \cite[Th.\,9]{BZ:04}.
\item $(\M, g)$ is a compact, boundaryless Riemannian manifold of dimension $d$ and constant curvature $\equiv -1$, and the observation $\mc{C}\psi= a\psi$ is made through a smooth function $a$ on $M$ such that the set
$\{\rho\in S^*\M\ |\ a^2(\phi_t(\rho))=0,\ \forall t\in \R\}$
has Hausdorff dimension $<d$. Here $\phi_t$ is the bicharacteristic flow on $T^*\M$. See Anantharaman-Rivi\`ere \cite[Th.\,2.5]{AR:12}.
\end{enumerate}
We also provide an abstract result following Lebeau \cite{Leb92}, establishing a connection between the observability of the Schrödinger equation and the plate equation.

We expect that the unique continuation for the situation in \cref{thm:analytic-nlp} is true. Yet, we~would need the unique continuation for the plate equation with lower order coefficients analytic in time. It~is very likely to be true, but is not yet proved. We~refer to \cite{Tat99,RZ:98,Hor:97} in analytic regularity and \cite{FLL25} in Gevrey spaces for the closely related Schrödinger equation.

Systems of nonlinear plate equations as \eqref{eq:nlp-1intro} have been considered by Eller-Toundykov \cite{ET:15} for $d=2$ with $\M$ a bounded open subset of $\R^2$, where they have addressed the question of (semi-global) exact controllability for such system. It~is known that establishing a unique continuation property for PDEs is a crucial step to achieve controllability results. They proved using Carleman estimates that unique continuation holds for \eqref{eq:nlp-1}, when $\omega$ is a subset of a collar neighborhood of the boundary $\partial \M$ (namely, a multiplier-type condition) and without an analyticity assumption on the coefficients, see \cite[Lem.\,4.1]{ET:15}. We~also refer to some recent results on nonlinear plates due to Bournissou-Ervedoza-Tucsnak \cite{TBE24}, where they show that the nonlinear system described by the von K\'arm\'an plate equation is locally exactly controllable around any stationary state defined by a real analytic function. Actually, in~\cite{ET:15}, more general models than the one presented here have been considered. However, it~was stressed by the authors that the study of \eqref{eq:nlp-1intro} presents a stepping stone to a further study of control-related questions for systems with even more complicated nonlinearities.
\subsection{An abstract result}\label{s:abstrIntro} Even though we have presented several results related to the semilinear wave and plate equations, at the core of all of them lies an abstract result. This result essentially states that solutions to a nonlinear problem, with a skew-adjoint linear part, and a compact nonlinearity which is also analytic, are analytic in time, provided that the observed solution is zero.

We first need to introduce some notation and assumptions in order to state the aforementioned result. Let $T>0$ and let us consider the nonlinear observability system
\begin{align}\label{intro:eq-obs}
\left\{\begin{aligned}
\partial_t U&=AU+F(U+H_1)+H_2, &t&\in (0, T), \\
\bC U(t)&=0, &t&\in (0, T),
\end{aligned}\right.
\end{align}
on a suitable real Hilbert space $X$, where $A$ is a skew-adjoint operator on $X$, $F$ is a mapping on $X$, $H_1$ and $H_2$ are some parameters, and $\bC$ is a bounded observation operator in $X$. The first assumption dictates the class of PDEs we will be working with.
\begin{assump}{1}
\label{assumAA}$A$ is a skew-adjoint operator with domain $D(A)$ on a real separable Hilbert space $X$, so that $A^*A=-A^2$ has a compact resolvent.
\end{assump}
This assumption allows us to introduce $X^\sigma$ as the interpolation space in between $D(A)$ and $X$ for $\sigma\in [0, 1]$, with $X^0=X$ and $X^1=D(A)$. From now on, we~will work at a fixed level of regularity, for which we fix $\sigma\in [0, 1]$ and instead consider \eqref{intro:eq-obs} in $X^\sigma$. We~will need to exert some control on the linear semigroup $t\in [0, T]\mto e^{At}\in \mc{L}(X^\sigma)$ generated by $A$ and even more, we~will need some extra control at a slightly higher regularity. This will be accomplished through the following assumption.

\begin{assump}{2}
\label{assumCC}Let $\bC\in\mc{L}(X^\sigma, X^\sigma)$ be an observation operator. We~assume that $t\mto e^{tA}$ is observable on $[0,T]$, namely, there exists a constant $\mathfrak{C}_{obs}>0$ such that
\begin{align}\label{observabstractintro}
\norm{W_0}_{X^\sigma}^2\leq \mathfrak{C}_{\textup{obs}}^2\int_0^T \norm{\bC e^{tA}W_0}_{X^\sigma}^2dt,\quad \forall W_0\in X^\sigma.
\end{align}
Furthermore, we~assume that $\bC\in\mc{L}(X^{\sigma+\veps}, X^{\sigma+\veps})$ for some $\veps>0$ and that there exists (another) constant $\mathfrak{C}_{obs}>0$ such that
\begin{align}\label{observabstractvepsintro}
\norm{W_0}_{X^{\sigma+\veps}}^2\leq \mathfrak{C}_{\textup{obs}}^2\int_0^T \norm{\bC e^{tA}W_0}_{X^{\sigma+\veps}}^2dt,\quad \forall W_0\in X^{\sigma+\veps}.
\end{align}
\end{assump}

We now make the last assumption regarding the nonlinearity. For a given real Banach space $Y$, we~denote the ball centered at $0$ of radius $M$ by
\begin{align*}
\mathbb{B}_M(Y):=\{y\in Y\ |\ \norm{y}_Y\leq M\}.
\end{align*}
For a given interval $I\subset \R$, we~denote the ball of radius $M$ of $C^0(I, Y)$ by
\begin{align*}
\B_{M}^I(Y)=\{U\in C^0(I, Y)\ |\ \forall t\in I,\ \norm{U(t)}_Y\leq M\},
\end{align*}
Moreover, we~introduce the canonical complexification $Y_\mathbb{C}$, defined as the set of elements $y_1+iy_2$, $y_j\in Y$, see \cite[\S 2]{BS:71-polynomials} for more details. We~then introduce the notation for the \emph{cylinder} on $Y_\mathbb{C}$
\begin{align*}
\mathbb{B}_{M, \delta}(Y)=\{y\in Y_\mathbb{C}\ |\ \norm{\Re(y)}_Y\leq M\ \text{and}\ \norm{\Im(y)}_Y\leq \delta\},
\end{align*}
and similarly on $C^0(I, Y_\mathbb{C})$
\begin{align*}
\B_{M, \delta}^I(Y)=\{U\in C^0(I, Y_\mathbb{C})\ |\ \forall t\in I,\ \norm{\Re(U(t))}_Y\leq M\ \text{and}\ \norm{\Im(U(t))}_Y\leq \delta\}.
\end{align*}
The latter space is naturally endowed with the $L^\infty(I,Y_\mathbb{C})$-norm. When working on balls not centered at $0$, we~will simply denote by $B_Y(y_0, r)$ the ball of center $y_0$ and radius $r>0$ on $Y$. With the previous notations, we~will make the following assumption on $F$.
\begin{assump}{3}
\label{assumFholom}
$F$ is a nonlinear Lipschitz and bounded operator from $\mathbb{B}_{4R_{0}}(X^{\sigma})$ into $X^{\sigma+\veps}$ for some $R_{0}>0$, $\sigma>0$ and $\veps>0$. Furthermore, there exists $\delta>0$ such that~$F$ has a holomorphic extension from $\mathbb{B}_{4R_{0},2\delta}(X^{\sigma})$ into $X^{\sigma+\veps}_{\mathbb{C}}$, namely, it~is holomorphic on the interior of $\mathbb{B}_{4R_{0},2\delta}(X^{\sigma})$ and continuous up to the boundary. Moreover, there exists $C>0$ such that
\begin{align}
\label{boundFanalyticintro}
\nor{F(U_{0})}{X^{\sigma+\veps}_{\mathbb{C}}}\leq C,\quad \nor{F(U_{0})-F(V_{0})}{X^{\sigma+\veps}_{\mathbb{C}}}\leq C \nor{U_{0}-V_{0}}{X^{\sigma}_{\mathbb{C}}}
\end{align}
for any $U_{0},~V_{0}\in \mathbb{B}_{4R_{0},2\delta}(X^{\sigma})$.
\end{assump}

Unless stated otherwise, we~will understand holomorphic extensions in the sense given in the assumption above. We~will also need a technical assumption on the pair $(A, \bC)$ related to commutator estimates, useful for the regularization of solutions.
\begin{assump}{4}
\label{assumcommu}
There exists $s>0$ such that $[(A^*A)^s,\bC]\in \mc{L}(X^{\sigma+2s}, X^{\sigma+\veps})$.
\end{assump}

We now state the main abstract result.

\begin{theorem}\label{thmabstractanalyticintro}
Let $R_{0}\!>\!0$ and $T\!>\!0$. Suppose that Assumptions \ref{assumAA}, \ref{assumCC} (with~$T$), \ref{assumFholom} (with $R_{0}$) and \ref{assumcommu} hold. Let $T^{*}\!>\!T$. Let $H_{1}\in \B_{R_{0}}^{[0,T^{*}]}(X^{\sigma})$ and $H_{2}\in C^{0}([0,T^{*}],X^{\sigma+\veps})$ that admit some extension in the space $C^{0}([0,T^{*}]+i[-\mu,\mu],X^{\sigma}_{\mathbb{C}})$, respectively in \hbox{$C^{0}([0,T^{*}]+i[-\mu,\mu],X^{\sigma+\veps}_{\mathbb{C}})$}, with $\mu\!>\!0$, so that the map
\begin{align*}
\left\{\begin{aligned}
(0,T^{*})+i(-\mu,\mu) & \to X^{\sigma}_{\mathbb{C}}\\
z&\mto H_{1}(z)
\end{aligned}\right.
\end{align*}
is holomorphic. We~assume the same for $H_{2}$ with value in $X^{\sigma+\veps}_{\mathbb{C}}$. We~assume moreover that $\Re H_{1}(z)\in \mathbb{B}_{R_0}(X^{\sigma})$ for any $z\in [0,T^{*}]+i[-\mu,\mu]$.

Then any $U\in C^{0}([0,T^{*}],X^{\sigma})$ satisfying
\begin{align}\label{UCabstractT*intro}
\left\{\begin{aligned}
\partial_t U&=AU+F(U+H_{1})+H_{2}&& \textnormal{ on } [0,T^{*}], \\
\bC U(t)&=0 &&\textnormal{ for }t\in [0,T^{*}],
\end{aligned}\right.
\end{align}
is real analytic in $t$ in $(0,T^*)$ with value in $X^{\sigma}$.
\end{theorem}

The previous result can still hold with some variants of the aforementioned assumptions (for instance, considering different assumptions on $A$). However, we~have chosen to keep enough generality to showcase that the technique can be applied to other PDEs. Note that the assumption on $H_1$ and $H_2$ is satisfied if $H_1$, \resp $H_2$, are analytic on $(-\tau,T^*+\tau)$ for some $\tau>0$ with value in $X^{\sigma}$, \resp $X^{\sigma+\veps}$. This result can be seen as a sort of finite-time adaptation of an abstract result due to Hale-Raugel \cite[Th.\,2.5]{HR03} with the flexibility of adding analytic source terms.

Hale-Raugel were concerned with the regularity properties (including analyticity) of evolutionary equations whose solutions are defined on $t\in\R$ and lie on a compact invariant set. They prove that such solutions are as smooth (in time) as the nonlinearity, encompassing a wide range of PDEs, including dissipative hyperbolic equations, which, unlike parabolic ones, do not have smoothing properties. Their proof relies upon a generalized Galerkin procedure, already used by dynamicists to study the regularity of attractors in different contexts, see for instance \cite{FT:89,Go:00-torus,Go:18} and the references therein. The approach of Hale-Raugel was to find the high-frequency component as a fixed point of an associate adequate mapping, depending on the low-frequency component of the solution. Then the problem is reduced to study the system associated with the low-frequency component and the corresponding fixed point map parameterized by it. We~can roughly say that there are two key hypotheses for the technique to work: the \emph{exponential decay} of the linear semigroup and some sort of \emph{compactness} on the nonlinearity. To adapt Hale-Raugel's technique to a finite-time setting, we~replace the decay of the linear semigroup by its finite-time counterpart: the \emph{observability} of the linear semigroup. We~are then led to solve a nonlinear observability system to find the high-frequency component. To this end, we~will heavily rely on the observability properties of the linear semigroup generated by $A$ and the compactness of the nonlinearity $F$ to set up an appropriate fixed point. It~is worth mentioning that, to stabilize some semilinear damped wave equations without \ref{assumGCC}, Joly and the first author \cite{2020:joly-laurent:decay-nlw-no-gcc} successfully adapted this technique when a weaker decay of the linear semigroup is assumed.

\subsection{Outline of the article}\cref{sec:abstract-construction} is devoted to the proof of the abstract \cref{thmabstractanalyticintro}. It~also contains several preliminaries as the property of finite determining modes and the abstract propagation of regularity. \cref{SEC:wave-eq} contains the applications to the nonlinear wave equation. It~contains the verification that the abstract \cref{thmabstractanalyticintro} can be applied for a sufficiently high regularity index $\sigma$. It~also contains some propagation of regularity arguments that allow to reach this regularity $\sigma$ starting from the energy space. \cref{s:plate} contains the applications to the plate equation. In~the Appendix, we~gathered some results about complex analysis in Banach spaces and some geometric facts about the generalized geodesic flow that are used in the rest of the article.

\subsubsection*{Acknowledgements} The present article is certainly a consequence of all the earliest discussions between Romain Joly and both authors. We~warmly thank him for everything it brought to this work.

\section{Analytic reconstruction for nonlinear observability systems}\label{sec:abstract-construction}

The purpose of this section is to prove \cref{thmabstractanalyticintro}. For the reader's convenience, we~will briefly outline its proof, following the Galerkin decomposition introduced in Hale-Raugel \cite{HR03}. The key replacement will be the observability estimate that allows to reconstruct the state from the observation, at least for the high-frequency part.

Let $T>0$ and fix $\sigma\in [0, 1]$. From \cref{assumAA}, we~introduce the low and high-frequency projections $\P_n=\mathbbm{1}_{[0, n]}\big((AA^*)^{1/2}\big)$ and $\Q_n=I-\P_n$, respectively. Let $U=U(t)$ be a mild solution of \eqref{UCabstractT*intro} in $C^0([0, T], X^\sigma)$ and suppose $H_1=0$, $H_2=0$ for simplicity. Let us consider the splitting
\begin{align*}
U(t)=\P_n U(t)+\Q_nU(t)=V(t)+W(t),
\end{align*}
where $\big(V(t), W(t)\big)$ solves the following system
\begin{align*}
\left\{\begin{array}{rl}
\partial_t V(t)&=AV(t)+\P_n F(V+W), \\[3pt]
\partial_t W(t)&=AW(t)+\Q_n F(V+W), \\[3pt]
\bC V(t)&=-\bC W(t).
\end{array}\right.
\end{align*}
By Duhamel's formula, the high-frequency component $W$ can be written as
\begin{align*}
W(t)=e^{tA}W(0)+\int_0^t e^{(t-s)A}\Q_n F\big(V(s)+W(s)\big)ds.
\end{align*}
The observation condition $\bC V=-\bC W$ suggests that given $V$, we~can reconstruct $W$ by considering the corresponding nonlinear observability system. Indeed, according to \cref{assumCC}, the observability of the linear semigroup $t\in [0, T]\mto e^{tA}$ enables us to construct an initial condition $W(0)$ solely in terms of an observation. Forgetting first about the source term given by the nonlinearity, that would allow to reconstruct $W(0)$ in terms of the observation of $W$, which is $\bC W=-\bC V$. \cref{lmCauchyobs} below provides a generalization of this reconstruction problem when source terms are present. In~this context, the first part of \cref{assumFholom}, namely, that $F$ is a nonlinear map from bounded sets of $X^\sigma$ into $X^{\sigma+\veps}$, will allow, at frequency sufficiently large, to consider the nonlinearity as a perturbation and to complete this reconstruction procedure. More precisely, this will imply the existence of a nonlinear map $\mc{N}$ such that $W(0)=\mc{N}(V)$. Consequently, we~have the formula $W=\Phi_V(W)$, where $\Phi_V: C^0([0, T], \Q_nX^\sigma)\to C^0([0, T], \Q_nX^\sigma)$ is given by
\begin{align*}
\Phi_V(W)(t)=e^{tA}\mc{N}V(\cdot)+\int_0^t e^{(t-s)A}\Q_n F\big(V(s)+W(s)\big)ds,\ t\in [0, T].
\end{align*}
This suggests that we can find the high-frequency component $W^*$ as a fixed point with the low-frequency component $V$ as an input.

At this point, the solution $U$ can be represented as $U(t)=V(t)+W^*(V)(t)$, where~$V$ solves
\begin{align}\label{intro:eq-low-split}
\partial_t V(t)=AV+\P_n F(V+W^*(V)).
\end{align}
To demonstrate that $t\in (0, T)\mto U(t)\in X^\sigma$ is analytic, the second part of \cref{assumFholom}, namely, that $F$ admits a holomorphic extension, is essential. This will be achieved by establishing that $t\mto V(t)$ and $V\mto W^*(V)$ are both analytic maps. If instead, we~consider \eqref{intro:eq-low-split} as a differential equation on the space Banach space $C^0([0, T], \P_n X^\sigma)$, classical ODEs theory imply that $t\mto V(t)$ is as smooth as $F$, and therefore analytic. The uniform contraction principle further ensures that $W^*$ depends analytically on $V$, from which the result follows.

In what follows, we~develop these ideas towards the proof of the main theorem. Throughout the section, we~will prove several intermediate results under weaker assumptions, some of which are of independent interest. For instance, we~mention that \cref{prop:finite-det-modes} (Finite determining modes) and \cref{prop:abs-prop-reg} (Propagation of regularity) below do not require analyticity. Nevertheless, the hurried reader could consider all the following results under Assumptions~\ref{assumAA}, \ref{assumCC}, \ref{assumFholom} and \ref{assumcommu}, which will always contain the intermediate assumptions made throughout this section.

\Subsection{Linear reconstruction, determining modes and propagation of regularity} In this section, we~will revisit the assumptions outlined in the introduction, organizing them according to the different results we aim to establish in order to prove our main result.

From now on, we~work under \cref{assumAA}. We~will now list some consequences of such an assumption. That $A^*A=-A^2$ is non-negative self-adjoint, allows us to define the Hilbert space $X^{\sigma}=D((A^*A)^{\sigma/2})$ for any $\sigma\in\R$. Note that the assumptions imply $X^{\sigma+\veps}\hookrightarrow X^{\sigma} $ for any $\veps>0$. Unless specifically noticed, we~will often omit the embedding $\iota: X^{\sigma+\veps}\to X^{\sigma}$.

By the spectral theorem, and since $A^*A$ has a compact resolvent, thus, the spectrum of $A^*A$ is real and discrete, allowing us to construct an orthonormal basis of eigenvectors of $A^*A$ in $X$, denoted by $(E_j)_{j\in\N}$ and associated to the nonnegative eigenvalues $(\lambda_j)_{j\in \N}$ (ranged increasingly) with $\lambda_j \to+\infty$ as $j\to +\infty$. We~introduce the high-frequency projectors $\Q_n$ on the space $\overline{\text{Span}\{E_j\}_{j\geq n}}$ and then we set the low-frequency projection $\P_n=I-\Q_n$. Note that $A$ commutes with $\P_n$ and $A\P_n$ is a bounded operator of $X^{\sigma}$ to itself with norm $\inn{\ld_n}$.

The parameter $\sigma$ will be fixed from now on. We~will use the notation $\P_n X^{\sigma}$ or $\Q_n X^{\sigma}$ for that related image of the Hilbert space endowed with the topology of~$X^{\sigma}$. The restriction of the embedding $\iota$ to $\Q_n X^{\sigma+\veps}$(denoted with the same name) $\iota: \Q_n X^{\sigma+\veps}\to \Q_nX^{\sigma}$ has norm $\inn{\lambda_n}^{-\veps}$.

By the spectral theorem, since the spectrum of $A$ is purely imaginary, we~can define $e^{tA}$ for any $t\in \R$, together with the estimate $\norm{e^{tA}U_0}_{X^{\sigma}}=\norm{U_0}_{X^{\sigma}}$. Also, $e^{tA}$~commutes with $\P_n$ and $\Q_n$.

\subsubsection{Linear reconstruction} We will now introduce the following assumption regarding the observability of the semigroup generated by $A$, which corresponds to the first part of Assumption~\ref{assumCC}.
\begin{assump}{2a}
\label{assumC}Let $\bC\in\mc{L}(X^\sigma, X^\sigma)$ be an observation operator. We~assume that $t\mto e^{tA}$ is observable on $[0,T]$, namely, there exists a constant $\mathfrak{C}_{obs}>0$ such that
\begin{align}\label{observabstract}
\norm{W_0}_{X^\sigma}^2\leq \mathfrak{C}_{\textup{obs}}^2\int_0^T \norm{\bC e^{tA}W_0}_{X^\sigma}^2dt,\quad \forall W_0\in X^\sigma.
\end{align}
\end{assump}
Let $\mc{O}\in \mc{L}(X^\sigma, L^2([0, T], X^\sigma))$, defined by $\mc{O}:=\bC e^{\cdot A}$, be the observation operator of linear waves and $\mc{O}_n:=\mc{O}|_{\Q_nX^\sigma}$ the \emph{high-frequency} observation operator. The observability inequality \eqref{observabstract} can be written
\begin{align}\label{observabstractO}
\norm{W_0}_{X^\sigma}&\leq \mathfrak{C}_{\textup{obs}}\nor{\O W_{0}}{L^{2}([0,T],X^\sigma)},\quad\;\; \forall W_0\in X^\sigma,\\
\label{observabstractOn} \norm{W_0}_{X^\sigma}&\leq \mathfrak{C}_{\textup{obs}}\nor{\O_{n} W_{0}}{L^{2}([0,T],X^\sigma)},\quad \forall W_0\in \Q_n X^\sigma,
\end{align}
where the last inequality is uniform in $n\in \N$.
It implies that $\mc{O}$ is injective and it has a closed range. Moreover, since $\Q_nX^\sigma$ is closed in $X^\sigma$, then the \emph{high-frequency} observation operator $\mc{O}_n:=\mc{O}|_{\Q_nX^\sigma}$ has closed range as well, which allows us to define $\Pi_n$ as the orthogonal (according to the natural scalar product in $L^2([0, T], X^\sigma)$) projection onto its image $\text{Im}(\mc{O}_n)\subset L^2([0, T], X^\sigma)$. From now on, we~equip $\Y_n:=\text{Im}(\O_n)$ with the induced topology from $L^2([0, T], X^\sigma)$ which makes it a Banach space. By~\eqref{observabstractOn}, we~know that $\O_n: \Q_n X^\sigma\to \Y_n$ is a bijection, and hence, $\Y_n$ is closed, a bounded reconstruction operator $\O_n^{-1}: \Y_n\to \Q_n X^\sigma$ exists.

By applying the observability inequality \eqref{observabstractOn}, consequence of \cref{assumC}, we~get for any $y\in \Y_n \subset L^2([0, T], X^\sigma)$,
\begin{align}
\label{boundOn-1}
\norm{\O_n^{-1}y}_{X^\sigma}&\leq \mathfrak{C}_{\textup{obs}}\norm{\O_{n}\O_n^{-1}y}_{L^2([0, T], X^\sigma)}=\mathfrak{C}_{\textup{obs}}\norm{y}_{L^2([0, T], X^\sigma)},
\end{align}
where again, the inequality is uniform in $n\in \N$.

To ease notation, we~consider the operator $\mc{I}(t): g\mto \int_0^t e^{(t-s)A}g(s)ds$ and we denote $\mc{I}(\cdot)$ when the operator is seen with value in $C^{0}([0,T],Y)$ for a suitable Banach space. The above construction will enable us to solve an \emph{observability Cauchy problem}, which is the content of the following lemma.

\begin{lemma}
\label{lmCauchyobs}
There exists $C(T,\mathfrak{C}_{\textup{obs}},\nor{\bC}{\mc{L}(X^\sigma)})>0$ such that for any $n\in\N$, $H\in L^{1}([0,T],X^{\sigma})$ and $G\in L^{2}([0,T],X^{\sigma})$, there exists a unique $W\in C^{0}([0,T],\Q_{n}X^{\sigma})$ solution of
\begin{align}\label{solprojectlin}
\left\{\begin{array}{rl}
\partial_t W(t)&=AW(t)+\Q_{n}H, \\[3pt]
\Pi_n\bC W &=\Pi_{n}G.
\end{array}\right.
\end{align}
It satisfies $W(0)= W_0:=\O_n^{-1}\Pi_n\left[G-\bC\mc{I}(\cdot)\Q_n H\right]$ and is given by $W(t)=e^{tA}W_{0}+\mc{I}(t)\Q_nH$. We~denote $\FL$ the associated linear operator defined by $\FL(G,H):=W$. Moreover, we~have the estimate
\begin{align}
\label{estimFL}
\nor{\FL(G,H)}{C^{0}([0,T],\Q_{n}X^{\sigma})}\leq C \nor{\Pi_{n}G}{L^{2}([0,T],X^{\sigma})}+C\nor{\Q_{n} H}{L^{1}([0,T],X^{\sigma})}
\end{align}
\end{lemma}
\begin{proof}
To be solution of the first line of \eqref{solprojectlin}, it~is equivalent to be written as a Duhamel formula $W(t)=e^{tA}W_{0}+\mc{I}(t)\Q_nH$ for some $W_{0}\in \Q_{n}X^{\sigma}$. So, we~only need to compute $W_{0}$. With the previous formula, we~have $W\in C^{0}([0,T],\Q_{n}X^{\sigma})\subset L^{2}([0,T],\Q_{n}X^{\sigma})$ and we can compute
\begin{align*}
\Pi_n\bC W= \Pi_n\bC \left[e^{\cdot A}W_{0}+\mc{I}(\cdot)\Q_nH \right]=\Pi_n\bC \left[e^{\cdot A}W_{0} \right]+\Pi_n\bC \left[\mc{I}(\cdot)\Q_nH \right].
\end{align*}
Note that if $W_{0}\in \Q_{n}X^{\sigma}$, then $\bC\left[e^{\cdot A}W_{0}\right]=\O W_{0}=\O_{n} W_{0}$ and therefore, by~definition of $\Pi_{n}$,
\begin{align*}
\Pi_{n}\bC\left[e^{\cdot A}W_{0}\right]=\Pi_{n}\O_{n} W_{0}=\O_{n} W_{0}.
\end{align*}

In particular, since both belong to $\mathcal{Y}_n$, we~want $ \Pi_n\bC W=\Pi_{n}G$, we~should have
\begin{align*}
\O_n^{-1}\Pi_{n}G= \O_n^{-1} \Pi_n\bC W &=\O_n^{-1}\O_{n} W_{0}+\O_n^{-1}\Pi_n\bC \left[\mc{I}(\cdot)\Q_nH \right]\\ &=W_{0}+\O_n^{-1}\Pi_n\bC[\mc{I}(\cdot)\Q_nH].
\end{align*}
This gives that $W_{0}$ should be defined by the formula given in the lemma. It~indeed belongs to $\Q_{n}X^{\sigma}$ and therefore $W$, as defined, satisfies the second line of \eqref{solprojectlin} by reproducing the same computation backward, that is\vspace*{-3pt}
\begin{align*}
\Pi_n\bC W&=\Pi_n\bC\left[e^{\cdot A}W_{0}+\mc{I}(\cdot)\Q_nH\right] =\Pi_{n} \O_{n}W_{0}+\Pi_n\bC\mc{I}(\cdot)\Q_nH\\
&=\Pi_{n} \O_{n}\O_n^{-1}\Pi_n\left[G-\bC\mc{I}(\cdot)\Q_n H\right]+\Pi_n\bC\mc{I}(\cdot)\Q_nH\\
&=\Pi_n\left[G-\bC\mc{I}(\cdot)\Q_n H\right]+\Pi_n\bC\mc{I}(\cdot)\Q_nH=\Pi_{n}G.
\end{align*}
The uniqueness could actually be obtained from the unique definition of $W_{0}$ that we obtained, but we prefer to give a precise proof. We~consider the difference $R=W_{1}-W_{2}\in C^{0}([0,T],\Q_{n}X^{\sigma})$ between two such solutions $W_{1}$ and $W_{2}$. It~satisfies\vspace*{-3pt}
\begin{align*}
\left\{\begin{array}{rl}
\partial_t R(t)&=AR(t), \\[3pt]
\Pi_n\bC R &=0.
\end{array}\right.
\end{align*}
That is $R(t)=e^{tA}R(0)$ and $\Pi_n\bC R=\Pi_{n}\O R(0)=\Pi_{n}\O_{n}R(0)=\O_{n}R(0)$. In~particular, $R(0)=0$ by injectivity of $\O_{n}$.

Concerning the estimates, since $A$ is skew-adjoint on $X^{\sigma}$, standard semigroup estimates give\vspace*{-3pt}
\begin{align}
\nor{W}{C^{0}([0,T],\Q_{n}X^{\sigma})}\leq \nor{W_{0}}{X^{\sigma}}+\nor{\Q_{n} H}{L^{1}([0,T],X^{\sigma})}.
\end{align}
So, we~need to estimate $W_{0}$. For any $\widetilde{G}\in L^2([0, T], X^\sigma)$, applying \eqref{boundOn-1} to $\Pi_n \widetilde{G}\in \Y_n$, we~have\vspace*{-3pt}
\begin{align*}
\norm{\O_n^{-1}\Pi_n \widetilde{G}}_{X^\sigma}&\leq \mathfrak{C}_{\textup{obs}}\norm{\Pi_n\widetilde{G}}_{L^2([0, T], X^\sigma)}.
\end{align*}
In particular, we~can estimate $W_{0}$ by\vspace*{-3pt}
\begin{align*}
\nor{W_{0}}{X^{\sigma}}&= \nor{\O_n^{-1}\Pi_n\left[G-\bC\mc{I}(\cdot)\Q_n H\right]}{X^{\sigma}}\\ &\leq \mathfrak{C}_{\textup{obs}}\norm{\Pi_nG}_{L^2([0, T], X^\sigma)}+\mathfrak{C}_{\textup{obs}}\norm{\Pi_n\bC\mc{I}(\cdot)\Q_n H}_{L^2([0, T], X^\sigma)}.
\end{align*}
We can finally estimate by the unitarity of $\Pi_{n}$ and Hölder inequality in time\vspace*{-3pt}
\begin{align*}
\norm{\Pi_n\bC\mc{I}(\cdot)\Q_n H}_{L^2([0, T], X^\sigma)}&\leq \norm{\bC\mc{I}(\cdot)\Q_n H}_{L^2([0, T], X^\sigma)}\\ &\leq \nor{\bC}{\mc{L}(X^\sigma)}\norm{\mc{I}(\cdot)\Q_n H}_{L^2([0, T], X^\sigma)}\\
&\leq T^{1/2}\nor{\bC}{\mc{L}(X^\sigma)}\norm{\mc{I}(\cdot)\Q_n H}_{L^{\infty}([0, T], X^\sigma)}\\ &\leq T^{1/2}\nor{\bC}{\mc{L}(X^\sigma)}\norm{\Q_n H}_{L^{1}([0, T], X^\sigma)}.
\end{align*}
Recollecting the previous estimates, we~have finally proved\vspace*{-8pt}
\begin{multline*}
\nor{W}{C^{0}([0,T],\Q_{n}X^{\sigma})}\\\leq \mathfrak{C}_{\textup{obs}} \nor{\Pi_{n}G}{L^{2}([0,T],X^{\sigma})}+\bigl(1+T^{1/2}\mathfrak{C}_{\textup{obs}}\nor{\bC}{\mc{L}(X^\sigma)}\bigr)\nor{\Q_{n} H}{L^{1}([0,T],X^{\sigma})}.\qedhere
\end{multline*}
\end{proof}
\begin{remark}
It will be very important for what follows that the constant $C$ involved in the previous lemma is independent of $n\in\N$.
\end{remark}

\subsubsection{Finite determining modes}
As a first direct consequence of the previous result, we~can get a finite determining mode result: two solutions of a nonlinear equation with the same observation and the same low frequency are the same. This result will not be used directly later, but can be considered as an easier version of what will follow where we will actually construct the reconstruction operator and study its regularity.

We make the following assumption on the nonlinearity $F$, akin to a compactness property. This corresponds to the first part of Assumption~\ref{assumFholom}.

\begin{assump}{3a}
\label{assumF}
$F$ is a nonlinear operator from $\mathbb{B}_{4R_{0}}(X^{\sigma})$ to $X^{\sigma+\veps}$ for some $R_{0}>0$, $\sigma>0$ and $\veps>0$. Moreover, there exists $C>0$ such that
\begin{align}
\label{boundF}
\nor{F(U_{0})}{X^{\sigma+\veps}}\leq C,\quad \nor{F(U_{0})-F(V_{0})}{X^{\sigma+\veps}}\leq C \nor{U_{0}-V_{0}}{X^{\sigma}}
\end{align}
for any $U_{0},~V_{0}\in \mathbb{B}_{4R_{0}}(X^{\sigma})$.
\end{assump}

Note that the second bound implies the first one with another constant. We~show the following property of finite determining modes.

\begin{proposition}\label{prop:finite-det-modes}
Let $R_{0}>0$. Under Assumptions \ref{assumAA}, \ref{assumC} and \ref{assumF} (with $R_{0}$), there exists $n\in \N$ such that the following holds. Let $H\in L^1([0,T],X^{\sigma})$ and $G\in L^2([0, T], X^\sigma)$.

Let $U(t)$ and $\widetilde{U}(t)$ be two solutions on $(0,T)$ of
\begin{align*}
\left\{\begin{aligned}
\partial_t U&=AU+F(U)+H, &&\text{on } (0, T), \\
\bC U(t)&=G(t), && \text{for } t\in (0, T),
\end{aligned}\right.
\end{align*}
such that $\norm{U(t)}_{X^\sigma}\leq R_{0}$ and $\norm{\widetilde{U}(t)}_{X^\sigma}\leq R_{0}$ for all $t\in [0, T]$. If $\P_n U(t)=\P_n \widetilde{U}(t)$ for all times $t\in [0, T]$, then $U(t)\equiv \widetilde{U}(t)$ for all $t\in [0, T]$.
\end{proposition}
\begin{proof}
By assumption, $\P_{n}U=\P_{n}\widetilde{U}$ as maps in $\B_{R_{0}}^{[0,T]}(X^{\sigma})$. Let us consider the difference of solutions $Z=U-\widetilde{U}$. It~satisfies
\begin{align*}
\left\{\begin{aligned}
\partial_t Z&=AZ+F(U)-F(\widetilde{U}), \\
\bC Z&=0.
\end{aligned}\right.
\end{align*}
Moreover, we~have $\P_n Z=0$, that is $Z\in C^{0}([0,T],\Q_{n}X^{\sigma})$ and therefore, applying~$\Q_n$, it~also satisfies
\begin{align*}
\left\{\begin{aligned}
\partial_t Z&=AZ+\Q_n \bigl(F(U)-F(\widetilde{U})\bigr), \\
\Pi_{n}\bC Z&=0.
\end{aligned}\right.
\end{align*}
In particular, we~are in the situation of \cref{lmCauchyobs} and $Z=\FL(0,\Q_n \big(F(U)-F(\widetilde{U})\big))$, so that estimate \eqref{estimFL} gives
\begin{align*}
\nor{Z}{C^{0}([0,T],\Q_{n}X^{\sigma})}&\leq C \bnor{\Q_n \bigl(F(U)-F(\widetilde{U})\bigr)}{L^{1}([0,T],X^{\sigma})}\\
&\leq \dfrac{C}{\inn{\ld_n}^\veps} \bnor{F(U)-F(\widetilde{U})}{L^{1}([0,T],X^{\sigma+\veps})}.
\end{align*}

By hypothesis, $U$, $\widetilde{U}\in \B_{R_0}^{[0, T]}(X^\sigma)$, and \cref{assumF} implies that $F(U)$, $F(\widetilde{U})\in \B_{C}^{[0,T]}(X^{\sigma+\veps})$ with the following Lipschitz estimate
\begin{align*}
\nor{Z}{C^{0}([0,T],X^{\sigma})}\leq \dfrac{C}{\inn{\ld_n}^\veps}\, \bnor{U-\widetilde{U}}{L^{1}([0,T],X^{\sigma})}\leq \dfrac{CT}{\inn{\ld_n}^\veps} \nor{Z}{C^{0}([0,T],X^{\sigma})}.
\end{align*}
Adjusting $n$ if necessary so that $\sfrac{CT}{\inn{\ld_n}^\veps}<1$, we~conclude that $Z=0$, which is $U= \widetilde{U}$.
\end{proof}

\subsubsection{Propagation of regularity} We now turn our attention to the regularity of the nonlinear system, for which we prove a propagation-of-regularity result under suitable assumptions. This result will later be related to a sort of uniformity for the splitting in the Galerkin procedure.

Building upon \cref{assumC}, we~make the following assumption, which will allow us to control the semigroup generated by $A$ in a slightly more regular space:
\begin{align}\label{observabstractveps}
\norm{W_0}_{X^{\sigma+\veps}}^2\leq \mathfrak{C}_{\textup{obs}}^2\int_0^T \norm{\bC e^{tA}W_0}_{X^{\sigma+\veps}}^2dt,\quad \forall W_0\in X^{\sigma+\veps}.
\end{align}
In that context, \cref{assumCC} is the conjunction of \cref{assumC} and \eqref{observabstractveps}, that is, the observability at both levels or regularity $X^{\sigma}$ and $X^{\sigma+\veps}$.

\begin{proposition}\label{prop:abs-prop-reg}
Let $R_{0}>0$, $R_{1}>0$ and $A$, $\bC$ and $F$ satisfying Assumptions~\ref{assumAA},~\ref{assumCC} and~\ref{assumF}, respectively. Moreover, assume that the pair $(A, \bC)$ satisfies \cref{assumcommu}. Then there exists $R_{2}>0$ such that for any
$U\in \B_{R_{0}}^{[0,T]}(X^{\sigma})$, $H_{1}\in \B_{R_{0}}^{[0,T]}(X^{\sigma})$ and $H_{2}\in \B_{R_{1}}^{[0,T]}(X^{\sigma+\veps})$ that satisfy \eqref{UCabstractT*intro} on $[0,T]$, we~have $U\in \B_{R_{2}}^{[0,T]}(X^{\sigma+\veps})$.
\end{proposition}
\begin{proof}

We know that $(AA^*)^{1/2}: X^1\to X$ admits a unique restriction such that $(AA^*)^{1/2}: X^{1+\sigma}\to X^{\sigma}$ is a linear continuous operator. Let us call such extension~$\A_\sigma$. Furthermore:
\begin{itemize}
\item From \cref{assumAA}, $\A_\sigma$ has compact resolvent.
\item $\A_\sigma$ and $(AA^*)^{1/2}$ have the same spectrum, hence the same resolvent set.
\end{itemize}
For simplicity, we~keep the same notation $\A$ for the same operator acting on different spaces. Since $\A$ is non-negative, its resolvent set contains $\R_-$ and for $n\in \N^*$, for any $s>0$, we~have a well-defined smoothing operator $\mc{J}_n=(I+\frac{1}{n}\A^s)^{-1}\in \mc{L}(X^\sigma, X^{\sigma+s})$ with the uniform bound $\norm{\mc{J}_n}_{\mc{L}(X^\sigma, X^{\sigma+s})}\leq n$. We~also have the uniform bound $\norm{\mc{J}_n}_{\mc{L}(X^\sigma)}\leq 1$ and the same estimate holds in $\mc{L}(X^{\sigma+\veps})$. Following \cite[Prop.\,2.3.4]{TW:09}, we~can see that $\J_n \phi\underset{n\to+\infty}{\to} \phi$ for any $\phi\in X^\sigma$.

Let us consider $U^n(t)=\J_nU(t)$ and observe that $U^n$ is uniformly bounded in $X^\sigma$ by some constant $C>0$, which is independent of $n\in \N$. By~Duhamel's formula, let~us split the solution into its linear and nonlinear part as follows
\begin{align*}
U^n(t)&\hphantom{:}=e^{tA}U_0^n+\int_0^t e^{(t-s)A}\big(\J_nF(U(s)+H_1(s))+\J_nH_2(s)\big)ds\\ &:=U_{\lin}^n(t)+U_{\Nlin}^n(t),
\end{align*}
where $U_0^n:=\J_n U_0$. Observe that we have used that $e^{tA}$ and $\J_n$ commute. Since \hbox{$H_2\in \B_{R_{1}}^{[0, T]}(X^{\sigma+\veps})$} and $F$ satisfies \cref{assumF}, we~readily get that $U_{\Nlin}^n$ is uniformly bounded in $L^\infty([0, T], X^{\sigma+\veps})$ by some constants depending on $R_{1}$ and the constant $C$ in \cref{assumF}. The latter property and \cref{assumCC}, allows us to treat the linear part by employing the observability inequality followed by the triangle inequality
\begin{align*}
\norm{U_{\lin}^n(t)}_{X^{\sigma+\veps}}^2&\leq \norm{e^{tA}}_{\mc{L}(X^{\sigma+\veps})}^2 \norm{U_0^n}_{X^{\sigma+\veps}}^2\\
&\leq \mathfrak{C}_{\textup{obs}}^2\int_0^T \norm{\bC U_{\lin}^n(t)}_{X^{\sigma+\veps}}^2dt\\
&\leq 2\mathfrak{C}_{\textup{obs}}^2\int_0^T \norm{\bC U^n(t)}_{X^{\sigma+\veps}}^2dt+2\mathfrak{C}_{\textup{obs}}^2\int_0^T \norm{\bC U_{\Nlin}^n(t)}_{X^{\sigma+\veps}}^2dt.
\end{align*}
Since $\bC U\equiv 0$ on $[0, T]\times \M$, we~get that
\[
\bC U^n=\bC\mc{J}_n U=\mc{J}_n\bC U+[\mc{J}_n,\bC]U=[\mc{J}_n,\bC]U.
\]
Moreover, we~have $[\mc{J}_n,\bC]=\frac{1}{n}\mc{J}_n[\A^s,\bC]\mc{J}_n$. We~have seen that $\norm{\mc{J}_n}_{\mc{L}(X^\sigma)}\leq 1$ and $\norm{\frac{1}{n}\mc{J}_n}_{\mc{L}(X^{\sigma}, X^{\sigma+s})}\leq 1$, so we get, uniformly in $n$,
\begin{align*}
\norm{\bC U^n}_{L^\infty([0, T], X^{\sigma+\veps})}&=\norm{\mc{J}_n[\A^s,\bC]\Psfrac{\mc{J}_n}{n} U}_{L^\infty([0, T], X^{\sigma+\veps})}\\
&\leq R_0\norm{[\A^s,\bC]}_{\mc{L}(X^{\sigma+s}, X^{\sigma+\veps})}.
\end{align*}
Therefore, in~view of \cref{assumcommu}, the term $\norm{[\A^s,\bC]}_{\mc{L}(X^{\sigma+s}, X^{\sigma+\veps})}$ is bounded and so $U_{\lin}^n$ is uniformly bounded in $L^\infty([0, T], X^{\sigma+\veps})$. It~follows that $U^n$ is uniformly bounded in $L^\infty([0, T], X^{\sigma+\veps})$ by some constant $C>0$. Moreover, due to the fact that $\norm{U-U^n}_{L^\infty([0, T], X^\sigma)}\to 0$ and by weak compactness in the Hilbert space $X^{\sigma+\veps}$, we~get that $U(t)\in X^{\sigma+\veps}$ and
\begin{align*}
\norm{U(t)}_{X^{\sigma+\veps}}\leq \liminf \norm{U^n(t)}_{X^{\sigma+\veps}}\leq C,
\end{align*}
for any $t\in [0, T]$, showing that $U$ is uniformly bounded in $L^\infty([0, T], X^{\sigma+\veps})$. We~obtain that $U\in C^0([0, T], X^{\sigma+\veps})$ using Duhamel formula now that we know that, for instance, $U(0)\in X^{\sigma+\veps}$ and $F(U+H_1)\in L^\infty([0, T], X^{\sigma+\veps})$.
\end{proof}
\begin{remark}\label{rk:abs-prop-reg}
Looking at the above proof, in~regards to the nonlinearity, we~only need to ensure that the assignment $t\mto \int_0^t e^{(t-s)A}F(Z(s))ds$ defines a bounded map in $L^\infty([0, T], X^{\sigma+\veps})$. For instance, this was achieved here through the sole hypothesis that $F(Z)$ is bounded in $L^\infty([0, T], X^{\sigma+\veps})$ for $Z$ bounded in $L^\infty([0, T], X^\sigma)$.
\end{remark}

\subsection{An abstract frequency-based reconstruction operator} The objective of this section is to prove that it is possible to reconstruct the high-frequency component for the solutions of our nonlinear system, with the low-frequency component as an input. Furthermore, this reconstruction can be made in a holomorphic way. For most of the facts about complex analysis in Banach spaces, our main reference is \cite{Muj86}.

Building upon \cref{assumF}, we~further assume that there exists $\delta>0$ such that~$F$ has a holomorphic extension from $\mathbb{B}_{4R_{0},2\delta}(X^{\sigma})$ to $X^{\sigma+\veps}_{\mathbb{C}}$ for some $R_{0}>0$, $\sigma\geq 0$ and $\delta>0$. Moreover, there exists $C>0$ such that
\begin{align}
\label{boundFanalytic}
\nor{F(U_{0})}{X^{\sigma+\veps}_{\mathbb{C}}}\leq C,\quad \nor{F(U_{0})-F(V_{0})}{X^{\sigma+\veps}_{\mathbb{C}}}\leq C \nor{U_{0}-V_{0}}{X^{\sigma}_{\mathbb{C}}}
\end{align}
for any $U_{0},~V_{0}\in \mathbb{B}_{4R_{0},2\delta}(X^{\sigma})$. In~this context, \cref{assumFholom} is the conjunction of \cref{assumF} and the previous assumption.

The main result of this section is the following.

\begin{theorem}
\label{thmabstract}Let $R_{0}$ and $R_{1}>0$. Under Assumptions \ref{assumAA}, \ref{assumCC}, \ref{assumF} (with $R_{0}$) and \ref{assumcommu}, there exist $n\in \N$ and a nonlinear Lipschitz (reconstruction) operator $\mathcal{R}$
\begin{align}
\label{Rspace}
\mathcal{R}: \quad \B_{R_{0}}^{[0,T]}(\P_n X^{\sigma}) \times \B_{R_{0}}^{[0,T]}(X^{\sigma}) \times \B_{R_{1}}^{[0,T]}(X^{\sigma+\veps}) \to \B_{R_{0}}^{[0,T]}(\Q_n X^{\sigma})
\end{align}
such that for any $U\in \B_{R_{0}}^{[0,T]}(X^{\sigma})$, $H_{1}\in \B_{R_{0}}^{[0,T]}(X^{\sigma})$ and $H_{2}\in \B_{R_{1}}^{[0,T]}(X^{\sigma+\veps})$ that satisfy
\begin{align}\label{UCabstract}
\left\{\begin{aligned}
\partial_t U&=AU+F(U+H_{1})+H_{2}, &&\textnormal{ on } [0,T], \\[3pt]
\bC U(t)&=0, &&\textnormal{ for }t\in [0,T],
\end{aligned}\right.
\end{align}
then $\Q_n U=\mathcal{R}(\P_{n}U,H_{1},H_{2})$.

Moreover, if~additionally $F$ satisfies \cref{assumFholom} then there exist $\eta,\eta_1>0$ such that, for any $U\in \B_{R_{0}}^{[0,T]}(X^{\sigma})$ that satisfies \eqref{UCabstract}, $\mathcal{R}$ extends holomorphically as
\begin{align*}
\mathcal{R}: \bigl( \P_nU+ \B_{\eta,\eta}^{[0,T]}(\P_n X^{\sigma})\bigr) \times \B_{R_{0},\eta}^{[0,T]}(X^{\sigma}) \times \B_{R_{1},R_{1}}^{[0,T]}(X^{\sigma+\veps}) \to \B_{R_{0},\eta_1}^{[0,T]}(\Q_n X^{\sigma}).
\end{align*}
\end{theorem}

The main idea is to consider the splitting
\begin{align*}
U=\P_n U+\Q_n U:=V+W,
\end{align*}
so the problem \eqref{UCabstract} translates into the following nonlinear observability system
\begin{align*}
\left\{\begin{array}{rl}
\partial_t V(t)&=AV(t)+\P_n F( V+W+H_{1})+\P_{n}H_{2}, \\[3pt]
\partial_t W(t)&=AW(t)+\Q_n F(V+W+H_{1})+\Q_{n}H_{2}, \\[3pt]
\bC V(t)&=-\bC W(t).
\end{array}\right.
\end{align*}
For a given bounded function $V\in C^0([0, T], \P_n X^\sigma)$, we~are interested in solving the nonlinear observability problem
\begin{align*}
\left\{\begin{array}{rl}
\partial_t W(t)&=AW(t)+\Q_n F(V+W+H_{1})+\Q_{n}H_{2}, \\[3pt]
\Pi_n\bC W &=- \Pi_n\bC V.
\end{array}\right.
\end{align*}
Note that the operator $\Pi_{n}$ is nonlocal in time in $[0,T]$.
\begin{proposition}\label{propabstract}
Let $R_{0}$ and $R_{1}>0$. Under Assumptions \ref{assumAA}, \ref{assumC} and \ref{assumF} (with $R_{0}$), there exists $n_{0}\in \N$ and $\eta>0$ such that for any $n\geq n_{0}$, for any $H_{1}\in \B_{5R_{0}/2}^{[0,T]}(X^{\sigma})$, $H_{2}\in \B_{R_{1}}^{[0,T]}(X^{\sigma+\veps})$ and $G\in \mathbb{B}_{\eta}(L^2([0, T], X^{\sigma}))$, there exists a unique solution $W\in \B_{R_{0}}^{[0,T]}(\Q_n X^{\sigma})$ to
\begin{align}\label{eqWG}
\left\{\begin{array}{rl}
\partial_t W&=AW+\Q_n F(W+H_{1})+\Q_{n}H_{2}, \\[3pt]
\Pi_n\bC W(t) &=\Pi_nG.
\end{array}\right.
\end{align}
This defines a nonlinear Lipschitz operator $\mathcal{\widetilde{R}}$
\begin{align}
\left\{\begin{array}{rcl}
\B_{5R_{0}/2}^{[0,T]}(X^{\sigma}) \!\times\! \B_{R_{1}}^{[0,T]}(X^{\sigma+\veps}) \!\times\! \mathbb{B}_{\eta}(L^{2}([0,T],X^{\sigma})) & \dpl\to& C^0([0, T], \Q_{n }X^{\sigma})\\[3pt]
(H_{1},H_{2},G)&\dpl\mto& W:=\mathcal{\widetilde{R}}(H_{1},H_{2},G).
\end{array}\right.
\end{align}
Moreover, if~additionally $F$ satisfies \cref{assumFholom}, then there exists $\eta>0$ such that~$\mathcal{\widetilde{R}}$ extends holomorphically as
\begin{align*}
\mathcal{\widetilde{R}}: \B_{5R_{0}/2,\eta}^{[0,T]}(X^{\sigma}) \times \B_{R_{1}, R_{1}}^{[0,T]}(X^{\sigma+\veps})\times \mathbb{B}_{\eta,\eta}(L^{2}([0,T],X^{\sigma})) \to C^0([0, T], \Q_{n }X^{\sigma}_{\mathbb{C}}).
\end{align*}
\end{proposition}
\begin{proof}[Proof of \cref{thmabstract} from \cref{propabstract}]
Let $\chi\in C^{\infty}_{0}(\R,[0,1])$ be supported in $[-1,1]$ such that $\chi(s)=1$ for $s\in [-1/2,1/2]$. We~define the operator $\mathcal{R}$ by the formula
\begin{align}\label{defR}
\mathcal{R}(V,H_{1},H_{2})=\mathcal{\widetilde{R}}(V+H_{1},H_{2},- \chi(\eta^{-1} \nor{\bC V}{L^{2}([0,T],X^{\sigma})}) \bC V),
\end{align}
where $\eta$ is the one in \cref{propabstract}. If $V$ is complex-valued, this is to be understood as $\chi(\eta^{-1} \nor{\bC V}{L^{2}([0,T],X^{\sigma}_{\mathbb{C}})})$. We~prove that it is well-defined and satisfies the requirements of \cref{thmabstract}.

We first observe that $\chi(\eta^{-1} \nor{\bC V}{L^{2}([0,T],X^{\sigma})}) =0$ if $\nor{ \bC V}{L^{2}([0,T],X^{\sigma})}\geq \eta$, so that we always have $\bnor{ \chi(\eta^{-1} \nor{ \bC V}{L^{2}([0,T],X^{\sigma})}) \bC V}{L^{2}([0,T],X^{\sigma})}\leq \eta$. In~particular, $\mathcal{R}$~is well-defined from the spaces made precise in \eqref{Rspace}.

We now~check that it satisfies the required properties. Let $U$ be as in the theorem and a solution of \eqref{UCabstract}. We~want to check that for $n$ large enough (only depending on the parameters), we~can impose $\Q_n U$ to be small in $C^{0}([0,T],X^{\sigma})$ and $\bC \P_n U$ to be small in $X^{\sigma}$. We~can apply \cref{prop:abs-prop-reg} to $U$ to propagate regularity, resulting in $\norm{U}_{C^0([0, T], X^{\sigma+\veps})}\leq R_{2}$ for some $R_{2}>0$. We~have
\begin{align}
\label{unifboundU}
\nor{ \Q_n U}{C^{0}([0,T],X^{\sigma})} &\leq \frac{1}{\inn{\lambda_n}^{\veps}} \nor{ \Q_n U}{C^{0}([0,T],X^{\sigma+\veps})}\leq \frac{R_{2}}{\inn{\lambda_n}^{\veps}} .
\end{align}
For such a solution satisfying \eqref{UCabstract}, we~have $\bC U=0$ and, in particular, $\bC \Q_n U=-\bC \P_n U$, so that
\begin{align*}
\nor{\bC \P_n U}{L^{2}([0,T],X^{\sigma})}=\nor{\bC \Q_n U}{L^{2}([0,T],X^{\sigma})}&\leq T^{1/2}\nor{\bC}{\mc{L}(X^\sigma)}\nor{ \Q_n U}{C^{0}([0,T],X^{\sigma})}\\
&\leq \frac{T^{1/2}\nor{\bC}{\mc{L}(X^\sigma)}R_{2}}{\inn{\lambda_n}^{\veps}} .
\end{align*}
Since $\lambda_n\to +\infty$ as $n\to +\infty$, we~can find $n\geq n_0$, so that the term above can be made smaller that $\eta/4$. In~particular, defining $W:= \mathcal{R}(\P_{n}U,H_{1},H_{2})$ and by definition of $\mathcal{R}$,
\begin{align*}
W= \mathcal{\widetilde{R}}(\P_{n}U+H_{1},H_{2},- \bC \P_{n}U)
\end{align*}
is the unique solution in $\B_{R_{0}}^{[0,T]}(\Q_n X^{\sigma})$ of
\begin{align*}
\left\{\begin{array}{rl}
\partial_t W&=AW+\Q_n F(W+\P_{n}U+H_{1})+\Q_{n}H_{2}, \\[3pt]
\Pi_n\bC W &=-\Pi_n \bC \P_{n}U.
\end{array}\right.
\end{align*}
Also, since $\bC \P_{n}U=- \bC \Q_{n}U $ we note that $\Q_n U$ is solution of
\begin{align*}
\left\{\begin{array}{rl}
\partial_t \Q_n U&=A\Q_n U+\Q_n F(\Q_n U+\P_n U+H_{1})+\Q_{n}H_{2}, \\[3pt]
\Pi_n\bC \Q_n U &=-\Pi_n \bC \P_{n}U.
\end{array}\right.
\end{align*}
In particular, since $ \nor{\bC \P_n U}{L^{2}([0,T],X^{\sigma})}\leq \eta/4$ and by the various bounds above, we~have
\[
\P_n U+H_{1}\in \B_{2R_{0}}^{[0,T]}(X^{\sigma}),\quad H_{2}\in \B_{R_{1}}^{[0,T]}(X^{\sigma+\veps})\quand-\bC \P_{n}U\in \mathbb{B}_{\eta}(L^2([0, T], X^{\sigma})).
\]
By definition of $\mathcal{\widetilde{R}}$, this implies $\Q_n U= \mathcal{\widetilde{R}}(\P_n U+H_{1},H_{2},- \bC \P_{n}U)$. So we conclude $\Q_{n}U=W=\mathcal{R}(\P_{n}U,H_{1},H_{2})$ as expected.

Concerning the holomorphic extension, since $\mathcal{\widetilde{R}}$ has a holomorphic extension, as proved in \cref{propabstract}, using formula \eqref{defR} and composition of function, it~is enough to prove that $V\mto \chi(\eta^{-1} \nor{\bC V}{L^{2}([0,T],X^{\sigma}_{\mathbb{C}})})$ is constant equal to $1$ and therefore holomorphic in some neighborhood $\P_nU+ \B_{\eta_{1},\eta_{1}}^{[0,T]}(\P_n X^{\sigma})$ that would be included in $\B_{3R_0/2,\eta/2}^{[0,T]}(X^{\sigma})$. This can be obtained for $\eta_{1}$ small enough since $ \nor{\bC \P_n U}{L^{2}([0,T],X^{\sigma})}\leq \eta/4$. This gives the result up to renaming $\eta_{1}$ by $\eta$.
\end{proof}

\begin{proof}[Proof of \cref{propabstract}]
Using \cref{lmCauchyobs} (recall that $\FL$ denotes a linear flow map), we~are looking for $W$ solution of\vspace*{-3pt}
\begin{align*}
W=\FL (G,F(W+H_{1})+H_{2}).
\end{align*}
So, it~is natural to define the nonlinear operator $\Phi_{n,G,H_{1},H_{2}}$ defined by\vspace*{-3pt}
\begin{align}\label{defPhi}
\Phi_{n,G,H_{1},H_{2}}(W):=\FL (G,F(W+H_{1})+H_{2})
\end{align}
which we will later prove to be well-defined from suitable balls of $C^0([0, T], \Q_n X^\sigma)$ to $C^0([0, T], \Q_n X^\sigma)$. To keep notations reasonable, we~will write $\Phi=\Phi_{n,V,G,H_{1},H_{2}}$ keeping in mind all the dependence. The goal will be to find a fixed point of the operator $\Phi$ in a small ball that will also satisfy \eqref{eqWG}.

Following Hale-Raugel \cite[Th.\,2.14]{HR03} and Joly-Laurent \cite[Th.\,10.1]{2020:joly-laurent:decay-nlw-no-gcc}, we~divide the proof in three steps.

\subsubsection*{Step 1. High-frequency fixed point: real case} For this part of the proof, we~consider all the functions involved to be real-valued, in~the sense that we consider the real vector space $X^{\sigma}$. Recall that throughout, we~work with the closed balls $\mathbb{B}$ and $\B$ defined at the beginning of Section~\ref{s:abstrIntro}.

We fix $H_{1}\in \B_{5R_{0}/2}^{[0,T]}(X^{\sigma})$, $H_{2}\in \B_{R_{0}}^{[0,T]}(X^{\sigma+\veps})$ and $G\in \mathbb{B}_{\eta}(L^2([0, T], X^{\sigma}))$, but we will make estimates uniform when these functions are in these sets. We~want to find a fixed point in $W\in \B_{R_{0}}^{[0,T]}(\Q_n X^{\sigma})$ by the fixed point theorem of Banach. We~prove that $\Phi$ is a contraction in this set.

Let $W\in \B_{R_{0}}^{[0,T]}(\Q_n X^{\sigma})$. Estimate \eqref{estimFL} of \cref{lmCauchyobs} gives for some constants uniform in $n\in \N$\vspace*{-5pt}
\begin{multline*}
\nor{\Phi(W)}{C^{0}([0,T],\Q_{n}X^{\sigma})}=\nor{\FL(G,F(W+H_{1})+H_{2})}{C^{0}([0,T],\Q_{n}X^{\sigma})}\\
\leq C \nor{\Pi_{n}G}{L^{2}([0,T],X^{\sigma})}+C\nor{\Q_{n}\left[ F(W+H_{1})+H_{2}\right]}{L^{1}([0,T],X^{\sigma})}.
\end{multline*}
Since $\Pi_{n}$ is a projection on $L^2([0,T],X^{\sigma})$, by~assumption we have\vspace*{-3pt}
\[
\nor{\Pi_{n}G}{L^{2}([0,T],X^{\sigma})}\leq \nor{G}{L^{2}([0,T],X^{\sigma})}\leq \eta.
\]
To estimate the second term, since $\nor{W+H_{1}}{C^{0}([0,T],X^{\sigma})}\leq 3R_{0}$, we~can use \eqref{boundF} and obtain\vspace*{-3pt}
\begin{equation}\label{estimFinPhi}
\begin{aligned}
\nor{\Q_{n}\left[ F(W+H_{1})\!+\!H_{2}\right]}{L^{1}([0,T],X^{\sigma})}&\leq \frac{1}{\inn{\lambda_n}^{\veps}} \nor{F(W+H_{1})\!+\!H_{2}}{L^{1}([0,T],X^{\sigma+\veps})}\\
&\leq \frac{T}{\inn{\lambda_n}^{\veps}}(C+R_{1}).
\end{aligned}
\end{equation}
Summing up, the previous estimates gives
\begin{align*}
\nor{\Phi(W)}{C^{0}([0,T],\Q_{n}X^{\sigma})}\leq C\eta+ \frac{1}{\inn{\lambda_n}^{\veps}} \left(C+R_{1}\right).
\end{align*}
Concerning the difference, similar estimates give for $W, W'\in \B_{R_{0}}^{[0,T]}(\Q_n X^{\sigma})$, using instead the second estimate of \eqref{boundF}
\begin{align}
\nonumber\nor{\Phi(W)-\Phi(W')}{C^{0}([0,T],\Q_{n}X^{\sigma})}&=\nor{\FL(0,F(W+H_{1})-F(W'+H_{1}))}{C^{0}([0,T],\Q_{n}X^{\sigma})}\\
\nonumber&\leq C\nor{\Q_{n}\left[ F(W+H_{1})-F(W'+H_{1}))\right]}{L^{1}([0,T],X^{\sigma})}\\[-7pt]
\label{bounddiffW}\\[-10pt]
\nonumber&\leq \frac{C}{\inn{\lambda_n}^{\veps}} \nor{F(W+H_{1})-F(W'+H_{1}))}{L^{1}([0,T],X^{\sigma+\veps})}\\
\nonumber&\leq \frac{CCT}{\inn{\lambda_n}^{\veps}} \nor{W-W'}{C^{0}([0,T],X^{\sigma})}.
\end{align}
In particular, if~$\eta$ is chosen small enough and $n_{0}$ is large enough, then $\Phi$ reproduces the closed ball $\B_{R_{0}}^{[0,T]}(\Q_n X^{\sigma})$ and is contracting.

\subsubsection*{Step 2. Extension to the complex case} The proof is very similar to the real case. We~define $\Phi$ with the same formula as \eqref{defPhi}, but instead consider $W\in \B_{R_{0},\eta_{1}}^{[0,T]}(\Q_n X^{\sigma})$, $H_{1}\in \B_{5R_{0}/2,\eta}^{[0,T]}(X^{\sigma})$, $H_{2}\in \B_{R_{1}, R_{1}}^{[0,T]}(X^{\sigma+\veps})$ and $G\in \mathbb{B}_{\eta,\eta}(L^{2}([0,T],X^{\sigma}))$, where $\eta_{1}>0$ small is to be determined later.

The flow map $\FL$ is linear, so the extension to complex-valued vectors is clear, see \cite[Th.\,3]{BS:71-polynomials}. So, we~only need to check that the valuation $F(W+H_{1})$ makes sense and the contraction in $\B_{R_{0},\eta_{1}}^{[0,T]}(\Q_n X^{\sigma})$ of $\Phi$ is still true with these parameters, up to making $\eta$ small and $n$ large.

If $W\in \B_{R_{0},\eta_{1}}^{[0,T]}(\Q_n X^{\sigma}), H_{1}\in \B_{5R_{0}/2,\eta}^{[0,T]}(X^{\sigma})$, then
\[
W+H_{1}\in \B_{7R_{0}/2,\eta_{1}+\eta}^{[0,T]}(X^{\sigma})\subset \B_{4R_{0},\delta}^{[0,T]}(X^{\sigma})
\]
if we have $\eta_{1}+\eta\leq \delta$, so that $F(W+H_{1})$ is well-defined and satisfies similar estimates as in the real case, thanks to \eqref{boundFanalytic} in \cref{assumFholom}. It~remains to check $\Phi(W)\in \B_{R_{0},\eta}^{[0,T]}(\Q_n X^{\sigma})$. We~still have
\begin{align*}
\nor{\Pi_{n}G}{L^{2}([0,T],X_\mathbb{C}^{\sigma})}\leq 2\eta,
\end{align*}
when $G\in \mathbb{B}_{\eta,\eta}(L^{2}([0,T],X^{\sigma}))$, while a similar estimate as in \eqref{estimFinPhi} gives
\begin{align*}
\nor{\Q_{n}\left[ F(W+H_{1})+H_{2}\right]}{L^{1}([0,T],X_\mathbb{C}^{\sigma})}&\leq \frac{1}{\inn{\lambda_n}^{\veps}} \left(C+2R_{1}\right).
\end{align*}
Finally, the linear estimate \eqref{estimFL} of \cref{lmCauchyobs} still holds in the complex-valued case, so we~get
\begin{align*}
\nor{\Phi(W)}{C^{0}([0,T],\Q_{n}X_\mathbb{C}^{\sigma})}\leq 2\eta+ \frac{1}{\inn{\lambda_n}^{\veps}} \left(C+2R_{1}\right).
\end{align*}
The bound \eqref{bounddiffW} holds in the complex case without modification. So, we~obtained that if $\eta_{1}+\eta\leq \delta$, $4\eta\leq \eta_{1}\leq R_{0}$ and $n$ is large enough, $\Phi$ is a contraction on $\B_{R_{0},\eta_{1}}^{[0,T]}(\Q_n X^{\sigma})$. This gives a unique fixed point in the complex case.

\subsubsection*{Step 3. Regularity of the fixed point} We prove that for fixed $n$, the map
\[
\widehat{\Phi}: (W, G, H_1, H_2)\mto \Phi_{n, G, H_1, H_2}(W)
\]
is Lispchitz of its arguments under \cref{assumF} and holomorphic when \cref{assumFholom} is made.
Recalling Definition \eqref{defPhi} and that $\FL$ is linear in its arguments, we~see that, by~composition, it~will be enough to verify that $(W,H_1)\mto F(W+H_{1})$ is Lipschitz or holomorphic, depending on the assumption made.

First of all, observe that the same argument to get estimates \eqref{bounddiffW}, shows that $(W,H_1)\mto F(W+H_{1})$ is Lipschitz under \cref{assumF}, and by linearity of the flow map $\mc{F}_n$, so is the map
\begin{multline*}
\widehat{\Phi}: \B_{R_0}^{[0, T]}(\Q_n X^{\sigma})\times \mathbb{B}_{\eta}(L^{2}([0,T],X^{\sigma}))\times \B_{2R_{0}}^{[0,T]}(X^{\sigma}) \times \B_{R_{1}}^{[0,T]}(X^{\sigma+\veps})\\
\to C^0([0, T], \Q_{n }X^{\sigma}).
\end{multline*}
Now, under \cref{assumFholom}, the same argument as before shows that the complex extension, denoted by the same letter $\widehat{\Phi}$,
\begin{multline*}
\widehat{\Phi}: \B_{R_0, \eta}^{[0, T]}(\Q_n X^{\sigma})\times \mathbb{B}_{\eta,\eta}(L^{2}([0,T],X^{\sigma}))\times \B_{2R_{0},\eta}^{[0,T]}(X^{\sigma}) \times \B_{R_{1}, R_{1}}^{[0,T]}(X^{\sigma+\veps})\\ \to C^0([0, T], \Q_{n }X_\mathbb{C}^{\sigma})
\end{multline*}
is Lipschitz, hence continuous. In~view of \cite[Prop.\,8.10]{Muj86} and that $\widehat{\Phi}$ is linear in the variables $G$ and $H_2$ separately, the main task is to show that the map
\begin{align}
\left\{\begin{array}{rcl}
\B_{4R_{0}, \eta}^{[0,T]}(X_{\mathbb{C}}^{\sigma}) \times \B_{2R_{0}, \eta}^{[0,T]}(X^{\sigma}) & \dpl\to& C^0([0, T], X_\mathbb{C}^{\sigma+\veps})\\[3pt]
(W, H_1 )&\dpl\mto& F(W+H_1)
\end{array}\right.
\end{align}
is holomorphic on each variable separately. Since $F$ is holomorphic when defined from some balls to $X^\sigma$ to $X^{\sigma+\veps}$, the only point to check is that it implies the same result as a map on time-dependent functions in $C^0([0,T],X^\sigma)$.

Fix $W\in \text{Int}(\B_{R_0, \eta}^{[0, T]}(\Q_nX^\sigma))$ and $H_1\in \text{Int}(\B_{5R_0/2, \eta}^{[0, T]}(X^\sigma))$, and $\ell>0$ such that $B_{C^0([0,T],X_{\C}^\sigma)}(W+H_1,\ell)\subset \B_{4R_0, 2\eta}^{[0, T]}(X^\sigma))$. Let $C$ be the bound on $F$ in \cref{assumFholom}. For all
\begin{itemize}
\item
$K_W\in C^0([0, T], \Q_n X_\mathbb{C}^\sigma)$,
\item
$K_H\in C^0([0, T], X_\mathbb{C}^\sigma)$ with $\norm{K_W+K_H}_{C^0([0, T], X_\mathbb{C}^\sigma)}\leq \ell/4$,
\item
and $t\in [0,T]$,
\end{itemize}
we~have $(W+H_1+K_W+K_H)(t) \in \mathbb{B}_{4R_{0},2\delta}(X^{\sigma})$ (recall $\eta<\delta$). Being $F$ holomorphic, by~applying a Taylor expansion (see \cite[Th.\,7.1, Cor.\,7.3]{Muj86}) together with Cauchy estimates (see \cite[Cor.\,7.4]{Muj86}), we~have the bound, uniform in $t\in [0,T]$,
\begin{align*}
\lVert F(W(t)+H_1(t)+K_W(t)+K_H(t))&-F(W(t)+H_1(t))\\
&-\delta F(W(t)+H_1(t); K_W(t)+K_H(t)) \rVert_{ X_\mathbb{C}^{\sigma+\veps}}\\
&\leq \dfrac{4C\norm{K_W(t)+K_H(t)}_{X_\mathbb{C}^\sigma}^2}{\ell(\ell-2\norm{K_W(t)+K_H(t)}_{X_\mathbb{C}^\sigma})} \\
&\leq \dfrac{4C\norm{K_W+K_H}_{C^0([0, T], X_\mathbb{C}^\sigma)}^2}{\ell(\ell-2\norm{K_W+K_H}_{C^0([0, T], X_\mathbb{C}^\sigma)})}\\
&\leq \dfrac{8C\norm{K_W+K_H}_{C^0([0, T], X_\mathbb{C}^\sigma)}^2}{\ell^2},
\end{align*}
where
\begin{align*}
\delta F(Z; K):=\lim_{s\to 0}\dfrac{1}{s}(F(Z+sK)-F(Z)).
\end{align*}
Observe that the differential of $F$ can be easily extended from $C^0([0, T], X_\mathbb{C}^\sigma)$ into $C^0([0, T], X_\mathbb{C}^{\sigma+\veps})$. Since, the previous estimate is uniform in $t\in [0,T]$, we~can write
\begin{multline}
\bigl\Vert F(W+H_1+K_W+K_H)-F(W+H_1)\\
\shoveright{-\delta F(W+H_1; K_W+K_H)\bigr\Vert_{C^0([0, T], X_\mathbb{C}^{\sigma+\veps})}}\\
\leq \dfrac{8C\norm{K_W+K_H}_{C^0([0, T], X_\mathbb{C}^\sigma)}^2}{\ell^2}.\label{prop:proof:frechet-1}
\end{multline}

Actually, the previous estimate \eqref{prop:proof:frechet-1} holds uniformly in a ball around $(W, H_1)$ of radius $\ell/4$. Consequently, we~establish that the map
\begin{align*}
\left\{\begin{array}{rcl}
\B_{R_{0}, \eta}^{[0,T]}(X^{\sigma}) \times \B_{5/2R_{0}, \eta}^{[0,T]}(X^{\sigma}) &\dpl\to& C^0([0, T], X_\mathbb{C}^{\sigma+\veps})\\[3pt]
(W, H_1 )&\dpl\mto& F(W+H_1)
\end{array}\right.
\end{align*}
is Fréchet differentiable (in the complex sense), hence holomorphic as an application of Theorem \cite[Th.\,13.16]{Muj86}. The estimates in \cref{lmCauchyobs} show that the map $\FL$ is bilinear continuous with respect to each of its parameters and therefore has an obvious holomorphic extension. Since the embedding $\iota: C^0([0,T],X_{\mathbb{C}}^{\sigma+\veps})\to L^1([0,T],X_{\mathbb{C}}^{\sigma})$ is linear, an application of the chain rule (see \cite[Th.\,13.6]{Muj86}) shows that the map $\widehat{\Phi}$ is holomorphic since it is defined by $\Phi_{n,G,H_{1},H_{2}}(W):=\FL (G,\iota\left(F(W+H_{1})+H_{2}\right)) $. The conclusion that the nonlinear reconstruction operator $\mathcal{\widetilde{R}}$ is holomorphic follows as a direct application of \cite[Th.\,2.2]{1982:chow-hale:bifurcation-theory} of regularity of fixed points with respect to parameters.
\end{proof}

\subsection{Analyticity in time of the observed solution} Here we prove the main theorem of this section, that is the abstract \cref{thmabstractanalyticintro} of analytic regularity stated in the introduction.
\begin{proof}[Proof of \cref{thmabstractanalyticintro}]
First, since $y\mto \max_{t\in [0,T^{*}]}\nor{\Im H_{1}(t+iy)}{X^{\sigma}}$ is a continuous function on $[-\mu,\mu]$ that is equal to zero at $y=0$, there exists $0<\mu'<\mu$ such that $\max_{t\in [0,T]}\left|\Im H_{1}(t+iy)\right| \leq \eta$ for $y\in [-\mu',\mu']$, where $\eta$ is given by \cref{thmabstract}. We~can do the same for $H_{2}$. We~denote $R_{1}=\max_{z\in [0,T^{*}]+i[-\mu,\mu]}\nor{H_{2}(z)}{X^{\sigma+\veps}}$.

Let $n\in \N$, $\eta>0$ and $\mathcal{R}$ be given by \cref{thmabstract}, with the chosen $R_0$ and $R_{1}$. Let~$\nu$ be such that $0<\nu<T^*-T$.

With these two choices, if~we define for $s\in [-\nu, T^*-T-\nu]$, $H_{1}^{s}$ as the map $t\mto H_{1}(t+\nu+s)$, we~have $H_{1}^{s} \in \B_{R_{0},\eta}^{[0,T]}(X^{\sigma})$ for any $s\in [-\nu, T^*-T-\nu]$. Moreover, the map $s\mto H_1^s$ is continuous. Indeed, since $t\in [0, T^*]\mto H_1(t)\in X^\sigma$ is a continuous map defined on a compact set, it~is uniformly continuous. Hence, for every $\veps>0$ and any $t\in [0, T]$, there exists $\delta>0$ such that for any $s$, $s_0\in [-\nu, T^*-T-\nu]$ with $|(s+t)-(s_0+t)|\leq \delta$, then $\norm{H_1(t+\nu+s)-H_1(t+\nu+s_0)}_{X^\sigma}\leq \veps$. Since the later property does not depend on $t$, we~can take $\sup_{t\in [0, T]}$ in the last inequality to obtain that $\norm{H_1^s-H_1^{s_0}}_{C^0([0, T], X^\sigma)}\leq \veps$. We~do the same for $H_2$, so we have $H_2^s\in \B_{R_{1},R_{1}}^{[0,T]}(X^{\sigma+\veps})$ for any $s\in [-\nu, T^*-T-\nu]$ and moreover $s\mto H_2^s$ is continuous.

For $U$, we~can define similarly $U^s\in \B_{R_{0},\eta}^{[0,T]}(X^{\sigma})$ for any $s\in [-\nu, T^*-T-\nu]$. We~can also decompose $U=\P_n U+\Q_n U=:V+W$ and $U^s=V^s+W^s$. For any fixed $s\in [-\nu, T^*-T-\nu]$, the assumption \eqref{UCabstractT*intro} implies
\begin{align}\label{UCabstractwiths}
\left\{\begin{aligned}
\partial_t U^s&=AU^s+F(U^s+H_{1}^s)+H_{2}^s&& \textnormal{on } [0,T],\\
\bC U^s(t)&=0 &&\textnormal{for }t\in [0,T].
\end{aligned}\right.
\end{align}
In particular, for any fixed $s\in [-\nu, T^*-T-\nu]$, \cref{thmabstract} gives
\begin{align}
\label{equalWs}
W^s=\Q_n U^s=\mathcal{R}(\P_{n}U^s,H_{1}^s,H_{2}^s)=\mathcal{R}(V^s,H_{1}^s,H_{2}^s),
\end{align} the equality being meant in $C^0([0,T], \Q_n X^{\sigma})$.

But if we denote
\[
G=\P_n F(U+H_{1})+\P_n H_{2}=\P_n F(V+W+H_{1})+\P_n H_{2}\in C^0([0,T^*],\P_n X^{\sigma}),
\]
then $V\in C^0([0,T^*],\P_n X^{\sigma})$ is solution of
\begin{align*}
\partial_t V=\P_n AV+G& \textnormal{ on } [0,T^{*}].
\end{align*}
With the related notations, \cref{lmtranslat} implies
\begin{align}
\label{eqVG}
\partial_s V^s=\P_n AV^s+G^s& \textnormal{ on } [-\nu, T^*-T-\nu].
\end{align}
But for any fixed $s\in [-\nu, T^*-T-\nu]$, we~have, using \eqref{equalWs},
\begin{align*}
G^s&=\P_n F(V^s+W^s+H_{1}^s)+\P_n H_{2}^s\\ &=\P_n F(V^s+\mathcal{R}(V^s,H_{1}^s,H_{2}^s)+H_{1}^s)+\P_n H_{2}^s
\end{align*}
with equality in $C^0([0,T], \P_n X^{\sigma})$.

It is then natural to consider the map
\begin{align*}
\Tilde{F}: [-\nu, T^*-T-\nu]\times \B_{R_{0}}^{[0,T]}(\P_n X^{\sigma})\to C^0([0, T], \P_n X^\sigma)
\end{align*}
defined by
\begin{align*}
\Tilde{F}(s,\Tilde{V})=\P_n F(\Tilde{V}+\mathcal{R}(\Tilde{V},H_{1}^s,H_{2}^s)+H_{1}^s)+\P_n H_{2}^s.
\end{align*}
This map is well-defined, continuous and locally Lipschitz (with respect to the second variable). We~observe that $H_1^s\in \B_{R_0}^{[0, T]}(X^\sigma)$, $H_2^s\in \B_{R_{1}}^{[0, T]}(X^{\sigma+\veps})$ for all \hbox{$s\in [-\nu, T^*-T-\nu]$} and $\Tilde{V}\in \B_{R_{0}}^{[0,T]}(\P_n X^{\sigma})$ implies, by~construction, that
\[
\Tilde{V}+\mathcal{R}(\Tilde{V},H_{1}^s,H_{2}^s)+H_{1}^s\in \B_{3R_0}^{[0, T]}(X^\sigma).
\]
Therefore $F(\Tilde{V}+\mathcal{R}(\Tilde{V},H_{1}^s,H_{2}^s)+H_{1}^s)$ is well-defined and so is $\Tilde{F}$. Recall that $\mc{R}$ is Lipschitz on its variables, that is, there exists $L>0$ such that
\begin{multline*}
\norm{\mc{R}(\Tilde{V}, H_{1}^{s}, H_2^{s})-\mc{R}(\Tilde{V}^*, H_{1}^{s^*}, H_2^{s^*})}_{C^0([0, T], X^\sigma)}\leq L\big( \norm{\Tilde{V}-\Tilde{V}^*}_{C^0([0, T], X^\sigma)}\\+\norm{H_1^{s}-H_1^{s^*}}_{C^0([0, T], X^\sigma)}+\norm{H_2^{s}-H_2^{s^*}}_{C^0([0, T], X^{\sigma+\veps})} \big),
\end{multline*}
for all $\Tilde{V}$, $\Tilde{V}^*\in \B_{R_{0}}^{[0,T]}(\P_n X^{\sigma})$ and $s$, $s^*\in [-\nu, T^*-T-\nu]$. The continuity of $\Tilde{F}$ then follows by the continuity of $s\mto (H_1^s, H_2^s)$, \cref{assumFholom} and by algebra of continuous maps. The same inequality along with \cref{assumFholom} shows that
\begin{align*}
\norm{\Tilde{F}(s, \Tilde{V})-\Tilde{F}(s, \Tilde{V}^*)}_{C^0([0, T], \P_n X^\sigma)}\leq C\norm{\Tilde{V}-\Tilde{V}^0}_{C^0([0, T], \P_n X^\sigma)},
\end{align*}
for any $\Tilde{V}$, $\Tilde{V}^*\in \B_{R_{0}}^{[0,T]}(\P_n X^{\sigma})$ and $s\in [-\nu, T^*-T-\nu]$.

Moreover, we~claim that for any $s_0\in (-\nu,T^*-T-\nu)$, there exists $\rho>0$ such that it admits a holomorphic extension (in both variables)
\begin{align*}
\Tilde{F}: B_{\mathbb{C}}(s_0, \rho)\times \bigl[V^{s_0}+\B_{\eta,\eta}^{[0,T]}(\P_n X^{\sigma})\bigr]\to C^0([0, T], \P_n X_\mathbb{C}^\sigma)
\end{align*}
given by
\begin{align*}
\Tilde{F}(z,\Tilde{V})=\P_n F(\Tilde{V}+\mathcal{R}(\Tilde{V},H_{1}^z,H_{2}^z)+H_{1}^z)+\P_n H_{2}^z.
\end{align*}
By \cref{thmabstract}, around $V^{s_0}=\P_n U^{s_0}$, the map $\mc{R}$ extends holomorphically as
\begin{align*}
\mathcal{R}: \bigl[ V^{s_0}+ \B_{\eta,\eta}^{[0,T]}(\P_n X^{\sigma})\bigr] \times \B_{R_{0},\eta}^{[0,T]}(X^{\sigma}) \times \B_{R_{1},R_{1}}^{[0,T]}(X^{\sigma+\veps}) \to C^0([0, T], \Q_{n }X_\mathbb{C}^{\sigma}).
\end{align*}
We need to argue that, the map $z\mto H_1^z$ is holomorphic from $B_{\mathbb{C}}(s_0, \rho)$ to $C^0([0, T], X_\mathbb{C}^\sigma)$ for $\rho>0$ small enough, and that the same holds for $z\mto H_2^z$ from $B_{\mathbb{C}}(s_0, \rho)$ with values in $C^0([0, T], X_\mathbb{C}^{\sigma+\veps})$. Indeed, since
\begin{align*}
z\in (0, T^*)+i(-\mu', \mu')\mto H_1(z)\in X_\mathbb{C}^\sigma
\end{align*}
is holomorphic, for each $s_0\in [-\nu, T^*-T-\nu]$ we can find $\rho>0$ such that, for any $t\in [0, T]$, the map
\begin{align*}
z\in B_\mathbb{C}(s_0, \rho)\mto H_1(z+t+\nu)\in X_\mathbb{C}^\sigma
\end{align*}
is holomorphic. As we did in the proof of \cref{propabstract}, by~using Cauchy estimates and the bound on $H_1(z)$, for each $t\in [0, T]$ and any $z_0\in B_\mathbb{C}(s_0, \rho)$, we~can find $\ell>0$ such that
\begin{multline*}
\norm{H_1(z+h+t+\nu)-H_1(z+t+\nu)-\delta H_1(z+t+\nu, h)}_{X_\mathbb{C}^\sigma}\\
\leq \dfrac{4(R_0+\eta)|h|^2}{\ell(\ell-2|h|)}\leq \dfrac{8(R_0+\eta)|h|^2}{\ell^2}
\end{multline*}
holds uniformly for all $z\in B_\mathbb{C}(z_0, \ell/4)$ and $B_\mathbb{C}(0, \ell/4)$. Moreover, the last estimate is independent of $t\in [0, T]$, so we actually have
\begin{align*}
\norm{H_1^{z+h}-H_1^z-\delta H_1(z+\cdot+\nu, h)}_{C^0([0, T], X_\mathbb{C}^\sigma)}\leq \dfrac{8(R_0+\eta)|h|^2}{\ell^2},
\end{align*}
uniformly for all $z\in B_\mathbb{C}(z_0, \ell/4)$ and $|h|\leq \ell/4$. The previous estimate shows that the map
\begin{align*}
z\in B_\mathbb{C}(s_0, \rho)\mto H_1^z\in \B_{R_0, \eta}^{[0, T]}(X^\sigma)
\end{align*}
is holomorphic. The argument works similarly to show that $z\mto H_2^z$ is holomorphic.

By composition of holomorphic functions and \cite[Prop.\,8.10]{Muj86} we get that
\begin{align*}
\left\{\begin{array}{rcl}
B_{\mathbb{C}}(s_0, \rho)\times \bigl[V^{s_0}+\B_{\eta,\eta}^{[0,T]}(\P_n X^{\sigma})\bigr]&\dpl\to& C^0([0, T], \Q_n X_\mathbb{C}^\sigma)\\[3pt]
(z,\Tilde{V})&\dpl\mto& \mathcal{R}(\Tilde{V},H_{1}^z,H_{2}^z)
\end{array}\right.
\end{align*}
is a well-defined and holomorphic map. That the extension $\Tilde{F}$ is holomorphic, follows from \cref{assumFholom}, algebra of holomorphic maps, and that $\P_n$ is a linear bounded operator.

Since we have proved that for any $s\in [-\nu, T^*-T-\nu]$, $G^s=\Tilde{F}(s,V^s)$, we~obtain that the equation \eqref{eqVG} verified by $V^s$ can be written as
\begin{align*}
\partial_s V^s=\P_n AV^s+\Tilde{F}(s,V^s)& \textnormal{ on } [-\nu, T^*-T-\nu].
\end{align*}
We then consider the following ODE, defined in the Banach space $C^0([0, T], \P_n X_\mathbb{C}^\sigma)$,
\begin{align}\label{thm:proof:eq:low-freq-evo-nlw}
\left\{\begin{array}{rl}
\partial_s \xi(s)&=\P_nA\xi(s)+\Tilde{F}(s, \xi(s)), \\[3pt]
\xi(s_0)&=V^{s_0}.
\end{array}\right.
\end{align}
Note that $C^0([0,T], \P_n X^\sigma)$ is of infinite dimension, but we can check that $\P_n A$ is a linear bounded operator on it. \cref{lmDuhamelCauchy} shows that $V^s$, which is a mild solution of \eqref{thm:proof:eq:low-freq-evo-nlw} (\ie satisfying \eqref{appendix:ode:ode-banach-1-duhamel}) is also a classical solution in the usual sense of ODE in Banach space.

Note that $\P_n A+\Tilde{F}$ is holomorphic from $ B_{\mathbb{C}}(s_0, \rho)\times \left[V^{s_0}+ \B_{\eta,\eta}^{[0,T]}(\P_n X^{\sigma})\right]$ into $C^0([0, T], \P_n X_{\mathbb{C}}^\sigma)$. Therefore, by~the classical theory of ODEs in Banach spaces \cite[Th.\,10.4.5]{1969:dieudonne:modern-analysis}, we~obtain a unique classical solution of \eqref{thm:proof:eq:low-freq-evo-nlw}
\begin{align*}
\xi:B_{\mathbb{C}}(s_0, \rho')\mto \Tilde{V}(z)\in C^0([0,T], \P_n X_{\mathbb{C}}^\sigma)
\end{align*}
for some $\rho'\in (0, \rho]$. Moreover, such a solution map inherits the regularity of the right-hand side of the ODE and so it is a holomorphic map. By~uniqueness of the solutions, we~have $\xi(s)=V^s$ for $s\in (s_0-\rho',s_0+\rho')$.

Moreover, the maps $z\mto \xi(z)$ and
\begin{align*}
\left\{\begin{array}{rcl}
B_{\mathbb{C}}(s_0, \rho)\times \bigl[V^{s_0}+ \B_{\eta,\eta}^{[0,T]}(\P_n X^{\sigma})\bigr]&\dpl\to& C^0([0, T], \P_n X_{\mathbb{C}}^\sigma)\\[3pt]
(z,\Tilde{V})&\dpl\mto&\Tilde{V}+ \mathcal{R}(\Tilde{V},H_{1}^z,H_{2}^z)
\end{array}\right.
\end{align*}
are both holomorphic in their respective spaces, and so is its composition\vspace*{-3pt}
\[
z\mto \xi(z)+\mathcal{R}(\xi(z),H_{1}^z,H_{2}^z).
\]

But for $s\in (s_0-\rho',s_0+\rho')$, we~have $\xi(s)=V^s$, so that
\[
\xi(s)+\mathcal{R}(\xi(s),H_{1}^s,H_{2}^s)=V^s+\mathcal{R}(V^s,H_{1}^s,H_{2}^s)=V^s+W^s=U^s,
\]
where we have used \eqref{equalWs} and the inclusion $(s_0-\rho',s_0+\rho')\subset (-\nu,T^*-T-\nu)$ for~$\rho'$ small enough. In~particular, the map $s\in (s_0-\rho',s_0+\rho')\mto U^s\in C^0([0, T], X^\sigma)$ is the restriction to a real interval of a holomorphic map, and hence is real analytic.

Since for any $t_0\in [0,T]$, the trace map $C^0([0, T], X^\sigma)\to X^\sigma$ is linear continuous, we~obtain by composition that the map\vspace*{-3pt}
\begin{align*}
\left\{\begin{array}{rcl}
(s_0-\rho',s_0+\rho')&\dpl\to& X^\sigma\\[3pt]
s&\dpl\mto&U^s(t_0)=U(t_0+\nu+s)
\end{array}\right.
\end{align*}
is real analytic. Recall that $\rho'$ depends on all the other parameters, but $\nu$ is an arbitrary number with $0<\nu<T^*-T$, $s_0$ is any number with $s_0\in (-\nu,T^*-T-\nu)$, while $t_0$ is arbitrary in $[0,T]$. It~means that $t\mto U(t)$, (which is well-defined for $t\in [0,T^*]$), is real analytic in a neighborhood of any $t_1$ of the form $t_1=t_0+\nu+s_0$ with $t_0$, $\nu$ and $s_0$ as before. It~is not hard to see that this implies that $t\mto U(t)$ is real analytic on $(0,T^*)$, as expected.
\end{proof}

\section{Application to the wave equation}\label{SEC:wave-eq}

The aim of this section is to prove the main results of propagation of analyticity and unique continuation related to the semilinear wave equation.

\subsection{Notation and preliminaries}\label{s:not_wave} Let $\beta\geq 0$. Set $X=H_0^1(\M)\times L^2(\M)$ and introduce the operator\vspace*{-3pt}
\begin{align*}\arraycolsep4.5pt
A=\begin{pmatrix}
0 & I \\
\Delta_g-\beta & 0
\end{pmatrix}\ \text{ with }\ D(A)=\big(H^2(\M)\cap H_0^1(\M)\big)\times H_0^1(\M).
\end{align*}
For the standard scalar product on $X$ (with $\norm{u}_{H^1_0}^2=\norm{\nabla u}_{L^2}^2+\beta \norm{u}_{L^2}^2$), we~can compute\vspace*{-3pt}
\begin{align*}
A^*=-A\ \text{ and }\ A^*A=-A^2=\begin{pmatrix}
-\Delta_g+\beta& 0 \\
0 & -\Delta_g+\beta
\end{pmatrix}\!,
\end{align*}
where we used the same notation for the Dirichlet Laplacian defined on $H^1_0(\M)$ and $L^2(\M)$ with their natural respective domains.
Observe that $A$ satisfies \cref{assumAA} on the real Hilbert space $X$. For $\sigma\geq 0$, $H^\sigma_D$ and $X^\sigma$ denote the spaces\vspace*{-3pt}
\begin{align}
\label{defSob}
H^\sigma_D&=D((-\Delta_g)^{\sigma/2}),\\
X^\sigma&= D((A^*A)^{\sigma/2})= D((-\Delta_g)^{(1+\sigma)/2})\times D((-\Delta_g)^{\sigma/2})=H^{1+\sigma}_D\times H^{\sigma}_D.
\end{align}
Here, $\Delta_g$ denotes the Dirichlet Laplacian defined on $L^2$.
\begin{remark}
For $\sigma \in [0,1/2)$, we~have $ X^\sigma=(H^{1+\sigma}(\M)\cap H_0^1(\M))\times H^\sigma(\M)$, while for $\sigma \in (1/2,1]$, we~have $ X^\sigma=(H^{1+\sigma}(\M)\cap H_0^1(\M))\times H_0^\sigma(\M)$, see \cite{G:67}. Note that $X^0=X$ and $X^1=D(A)$. When $\partial\M=\emptyset$, we~still write $H^1_D=H_0^1=H^1$ and assume $\beta>0$ so that Poincaré-like inequalities hold.
\end{remark}
By the spectral theorem, $-\Delta_g$ has a compact resolvent and thus we~can construct an orthonormal basis of eigenfunctions of $-\Delta_g$ in $L^2(\M)$, denoted by $(e_j)_{j\in \N}$ and associated to the eigenvalues $(\lambda_j)_{j\in \N}$. We~introduce the high-frequency projectors $Q_n$ on the space $\overline{\text{Span}\{e_j\}_{j\geq n}}$ and then we set the low-frequency projection $P_n=I-Q_n$. In~the same direction, we~then set $\Q_n=(Q_n, Q_n)$ and $\P_n=(P_n, P_n)$ on $X$, fitting with the abstract framework given in \cref{sec:abstract-construction}.

In regards to the nonlinearity, for some $\widetilde{\chi}\in C^{\infty}(\M)$ to be chosen later, we~use the notation\vspace*{-3pt}
\begin{align}
\label{defFwave}
F: (u, v)\mto (0, -\widetilde{\chi} f(u)+\beta u),
\end{align}
where $f: \R \to \R$ will be either $C^4$ or analytic. By~setting $U(t):=(u(t), \partial_t u(t))$, if~$\widetilde{\chi}=1$, we~can write the associated Cauchy problem to \eqref{eq:nlw-1} as\vspace*{-3pt}
\begin{align}\label{eq:nlw-1-cauchy}
\left\{\begin{array}{rl}
\partial_t U&=AU+F(U), \\[3pt]
U(0)&=U_0.
\end{array}\right.
\end{align}
We will later consider some slightly different problems, which is why we introduce the cutoff $\widetilde{\chi}$. Since $A$ is skew-adjoint, by~Stone's theorem it generates a unitary group $t\mto e^{tA}$ on $X$ and on $D(A)$. In~particular,\vspace*{-3pt}
\begin{align*}
\forall t\in [0, T],\quad \norm{e^{tA}}_{\mc{L}(X)}= 1\quand\norm{e^{tA}}_{\mc{L}(D(A))}=1.
\end{align*}
By linear interpolation, the same property holds on $X^\sigma$ for any $\sigma\in (0, 1)$.

We now establish some regularity properties for $F$. First of all, we~recall a version of a result that can be found in Alinhac-Gérard \cite[Prop.\,2.2]{AG91}, in~relation to the regularity of a composition.
\begin{lemma}\label{lem:comp-reg}
Let $g: \R\to \R$ be a $C^3(\R, \R)$ function, with $g(0)=0$. If $u\in L^\infty(\M)\cap H^s(\M)$, with $s\in (0, 2)$, then $g(u)\in L^\infty(\M)\cap H^s(\M)$ and $\norm{g(u)}_{H^s}\leq C\norm{u}_{H^s}$, where~$C$ depends only on $g$ and $\norm{u}_{L^{\infty}}$.
\end{lemma}

\begin{remark}
The previous lemma is actually written in \cite{AG91} for function in $H^s(\R^d)$ and $f\in C^{\infty}$. Yet, the same result holds for functions in $H^s(\M)$ when $\M$ is a compact manifold with boundary, using the definition of the norm of $H^s(\M)$ by partition of unity and sum of the norm in $H^s(\R^d)$ of the functions in local coordinates and with extension. Concerning the requirement $g\in C^3$, we~only note that the Meyer multiplier lemma \cite[Lem.\,2.2.]{AG91} requires estimates of the derivatives of the multiplier up to $\floor{s}+1\leq 2$. Since it is applied to $g'$ in \cite[Prop.\,2.2]{AG91}, it~requires $g'\in C^2$. Note that, as proved in \cite{G:67}, for functions satisfying the appropriate boundary condition, (that is $u=0$ on $\partial \M$ for $1/2<s<5/2$ and no condition if $0\leq s<1/2$), the norms $H^s(\M)$ (defined as restrictions of functions in $H^s(\R^d)$ in local coordinates) and the norms $H^s_D(\M)$ defined by spectral theory are equivalent.
\end{remark}

We now prove that $F$ satisfies \cref{assumF} and \ref{assumFholom} under suitable hypothesis on the function $f$.

\begin{proposition}\label{prop:f-assumptions}
Let $\sigma\in (1/2, 1]$ if $d=3$ and $\sigma\in (0,1]$ if $d\leq 2$. Assume that $f\in C^4(\R,\R)$ and $\widetilde{\chi}\in C^{\infty}(\M)$. Then $F$, defined in \eqref{defFwave}, satisfies \cref{assumF} for any $\veps\in (0, 1]$ and $R_0>0$. If, in addition, $f$ is real analytic, then $F$ satisfies \cref{assumFholom}.
\end{proposition}
\begin{proof} We will use the notation introduced in \eqref{defSob} for the Sobolev spaces. In~what follows, $C$ will denote a generic constant, that may change from line to line, but we will specify its dependency on the parameters involved in the estimates. We~consider $\beta=0$ for simplifying the statements, since this term is easier to treat. Thanks to the choice of $\sigma$, we~have the embedding $H_D^{1+\sigma}\hookrightarrow L^\infty$ with constant denoted by $\kappa$. Since $f\in C^4$, we~can apply \cref{lem:comp-reg}, so that for any $v\in H_D^{1+\sigma}$, both $f(v)$ and $Df(v)-Df(0)$ are well-defined in $H_D^{1+\sigma}$ and, moreover, $\norm{Df(v)-Df(0)}_{H_D^{1+\sigma}}\leq C\norm{v}_{H_D^{1+\sigma}}$, where~$C$ depends on $Df$ and the $L^{\infty}$-norm of $v$. Therefore, using that $H^{1+\sigma}$ is an algebra, we~get, for any $v$, $v'\in H_D^{1+\sigma}$,
\begin{equation}\label{fineq1}
\begin{aligned}
\norm{f(v)&-f(v')}_{H_D^{1+\sigma}}=\Norm{\int_0^1 Df(v'+\tau(v-v'))(v-v')d\tau}_{H^{1+\sigma}}\\
&\leq C\biggl(1+\biggl\|\int_0^1 \big(Df(v'+\tau(v-v'))-Df(0)\big)d\tau\biggr\|_{H^{1+\sigma}}\biggr)\norm{v-v'}_{H_D^{1+\sigma}}\\
&\leq C\bigl(1+\norm{v}_{H_D^{1+\sigma}}+\norm{v'}_{H_D^{1+\sigma}}\bigr)\norm{v-v'}_{H_D^{1+\sigma}},
\end{aligned}
\end{equation}
where $C$ is a constant depending only on $Df$ and the $L^\infty$-norm of both $v$ and $v'$. Note that we have used $f(0)=0$ to get $\norm{f(v)-f(v')}_{H_D^{1+\sigma}}=\norm{f(v)-f(v')}_{H^{1+\sigma}}$, where we consider one norm of $H^{1+\sigma}$ defined in local coordinates so that \cref{lem:comp-reg} can be applied.

We now establish that $F$ satisfies \cref{assumF}. Let $V$, $V'\in \mathbb{B}_{4R_0}(X^\sigma)$. Observe that the first component of $V$ and $V'$, denoted by $v$ and $v'$, respectively, always stay smaller than $4\kappa R_0$ in $L^{\infty}$. This implies that $F$ is a well-defined map on $\mathbb{B}_{4R_0}(X^\sigma)$. For $\veps\in (0, 1]$, we~have $\sigma+\veps\leq 1+\sigma$ and thus the continuous embedding $H_D^{1+\sigma}\hookrightarrow H_D^{\sigma+\veps}$. Therefore,
\begin{align}\label{fineq2}
\norm{F(V)-F(V')}_{X^{\sigma+\veps}}=\norm{\widetilde{\chi}(f(v)-f(v'))}_{H_D^{\sigma+\veps}}&\leq C\norm{f(v)-f(v')}_{H_D^{1+\sigma}},
\end{align}
where $C$ is a constant depending on $\chi$. Combining inequalities \eqref{fineq1} and \eqref{fineq2}, we~get
\begin{align}\label{ineq:Flip1}
\norm{F(V)-F(V')}_{X^{\sigma+\veps}}&\leq C(1+8R_0)\norm{V-V'}_{X^\sigma},
\end{align}
where $C$ depends only $\widetilde{\chi}$, $f$ and $R_0$. Inequality \eqref{ineq:Flip1} proves the claim.

Let us show that $F$ satisfies \cref{assumFholom} if we assume that $f$ is real analytic. By~compactness, there exists $\delta>0$ small such that $f$ extends holomorphically into the interior of the complex strip
\begin{align*}
\mathbb{S}_{R_0, \delta}:=\{z_1+iz_2\in \mathbb{C}\ |\ |z_1|\leq 4\kappa R_0\ \text{and}\ |z_2|\leq 2\kappa \delta\}.
\end{align*}
Moreover, this extension is continuous up to the boundary and there exists a constant $M>0$ such that $|f(z)|\leq M$ for all $z\in \mathbb{S}_{R_0, \delta}$. We~still denote by $f$ such an extension. A slight variant of \cref{lem:comp-reg} allows to consider the composition of smooth functions defined on domains of $\mathbb{C}$ by functions in $H^{1+\sigma}$, assuming that the composition makes sense. Since $\kappa$ is the constant of the embedding $H_D^{1+\sigma}\hookrightarrow L^\infty$, we~see that $f(v)$ is well-defined in $\mathbb{B}_{4R_0, 2\delta}(H^{1+\sigma}_D)$. In~particular, $F$ is well-defined in $\mathbb{B}_{4R_0, 2\delta}(X^\sigma)$ and satisfies the same estimate as \eqref{fineq1}.

Next, let us prove that this map $F$ is $\mathbb{C}$-differentiable. We~write for $z\in \Int{\mathbb{S}_{R_0, \delta}}$ and $h\in\mathbb{C}$ small, $f(z+h)=f(z)+f'(z)h+h^2\int_0^1f^{(2)}(z+th)(1-t)dt$. So, for $v\in \Int{\mathbb{B}_{4R_0, 2\delta}(H^{1+\sigma}_D)} $, we~can write
\begin{align}\label{eq:f-exp}
f(v(x)+r(x))=f(v(x))+f'(v(x))r(x)+r(x)^2\hspace*{-1mm}\int_0^1\hspace*{-1mm}f^{(2)}(v(x)+tr(x))(1-t)dt,
\end{align}
for $r\in \mathbb{B}_{\eta, \eta}(H^{1+\sigma}_D)$ (for some $\eta$ depending on $v$). Since $f$, $f'$ and $f^{(2)}$ are smooth functions on the range of $v+tr$ for any $t\in [0,1]$, as before, all terms in \eqref{eq:f-exp} are well-defined and bounded in $H^{1+\sigma}$. The Dirichlet boundary condition being satisfied for each of the three terms, we~conclude that $f$ is Fréchet differentiable (in the complex sense) at~$v$, with derivative $r\mto f'(v)r$ which is continuous $\mathbb{C}$-linear from \hbox{$H^{1+\sigma}_D+iH^{1+\sigma}_D$} into $H^{1+\sigma}_D+iH^{1+\sigma}_D$ and therefore into $H^{\sigma+\veps}_D+iH^{\sigma+\veps}_D$. Therefore, $f$~can be extended to admit a bounded holomorphic extension from $\mathbb{B}_{4R_0, 2\delta}(H_D^{1+\sigma})$ into $H^{\sigma+\veps}_D+iH^{\sigma+\veps}_D$ and the required extension holds for $F$ after composition by linear bounded functions.
\end{proof}

\skpt
\begin{remark}\label{rk:proofhigherd}
Regarding \cref{rk:higherd}, one point in extending \cref{thm:analytic-prop} to any dimension $d\geq 3$ is to verify that the previous proposition holds for $\sigma\in (d/2-1, (d-1)/2)$ and assuming
\begin{equation*}
f^{(k)}(0)=0, \quad \forall~0\leq k\leq \lfloor d/2\rfloor-1,
\end{equation*}
so that $f: H^{1+\sigma}_D\to H^{\sigma}_D$ admits a holomorphic extension into the corresponding complexified spaces. Recall, see \cite{G:67}, that if $\sigma\notin 2\N+1/2$ and $\sigma>1/2$, $u\in H^{\sigma}_D$ if and only if $u\in H^{\sigma}$ and $\Delta_g^ku|_{\partial \M}=0$ for all $k\in \N$ with $2k+1/2< \sigma $. The choice of $\sigma$ ensures that $H^{1+\sigma}$ is an algebra embedded in $L^{\infty}$, so that the same estimates performed in \cref{prop:f-assumptions} hold. It~remains to verify that the boundary conditions required to have $f(u)\in H^{\sigma}_D$ are satisfied when $u\in H_D^{1+\sigma}$. Under these assumptions, $f$~can be written $f(z)=z^{p}h(z)$ for $h$ holomorphic and $p=\lfloor d/2\rfloor$. Since each operator~$\Delta_g^k$ is an operator of order $2k$ with $2k<\sigma-1/2$, in~suitable local coordinates near~$\partial\M$, the Leibniz formula gives $\Delta_g^k (f(u))=\sum_{|\alpha|+|\beta|\leq 2k} c_{\alpha,\beta}\partial^{\alpha}(u^{p})\partial^{\beta}(h(u))$ for some smooth coefficients $c_{\alpha,\beta}$. Since $|\alpha|\leq 2k<\sigma-1/2<d/2-1$, we~have $|\alpha|\leq \lfloor d/2\rfloor-1=p-1$. As $u\in H_D^{1+\sigma}$, we~have $u|_{\partial\M}=0$, and therefore $\partial^{\alpha}(u^{p})|_{\partial\M}=0$ for every $|\al|\leq p-1$. It~follows that $\Delta_g^k(f(u))|_{\partial\M}=0$ whenever $2k+1/2<\sigma$, and hence $f(u)\in H_D^\sigma$.
\end{remark}

\subsubsection{Well-posedness theory} For $d\leq 2$, and with assumption \eqref{hip:nonlinearity-hyp-1} when $1\leq p<+\infty$, the well-posedness theory in $H^1\times L^2$ can be performed with Sobolev embedding, so we omit the details. Assume that $\M$ is of dimension $d=3$. Later on, we~will need to use some results related to the global existence and uniqueness of solutions of the semilinear wave equation in the subcritical case. We briefly recall them.

A central argument to handle the subcritical case in dimension $d=3$ is the use of Strichartz estimates. They have a long history. We~only quote the results that we use and refer to the references therein for historical background. For general domains with boundary, Strichartz estimates were proved by Burq, Lebeau and Planchon \cite{BLP:08} and later, the range of admissible exponents was extended by Blair, Smith and Sogge \cite[Cor.\,1.2]{BSS:09} leading to the following theorem.

\begin{theorem}[Strichartz estimates]\label{thmStrichartz}
Let $T>0$ and $(q, r)$ satisfying
\begin{align}\label{thm:strichartz-exponents}
\dfrac{1}{q}+\dfrac{3}{r}=\dfrac{1}{2},\ \ \ q\in [7/2, +\infty].
\end{align}
There exists $C=C(T, q)>0$ such that, for every $G\in L^1([0, T], L^2(\M))$ and every $(u_0, u_1)\in X$, the mild solution $u$ of
\begin{align*}
\left\{\begin{array}{rl}
\partial_t^2 u-\Delta_g u&=G(t), \\[3pt]
u_{|_{\partial\M}}&=0, \\[3pt]
(u, \partial_t u)(0)&=(u_0, u_1),
\end{array}\right.
\end{align*}
satisfies the estimate
\begin{align*}
\norm{u}_{L^q([0, T], L^r(\M))}\leq C\big(\norm{(u_0, u_1)}_{H^1(\M)\times L^2(\M)}+\norm{G}_{L^1([0, T], L^2(\M))}\big).
\end{align*}
\end{theorem}

Without loss of generality, it~can be assumed that $p\in (3, 5)$ for the bound on $f$ in~\eqref{hip:nonlinearity-hyp-1}. The exponent $p=3$ is the exponent where Strichartz are no longer necessary and so the cases $p\in [1, 3]$ can be managed using appropriate Sobolev embeddings. Observe that it is enough to consider the pair of exponents $(q, r)\!=\!(\spfrac{2p}{p-3}, 2p)$, since they give $u^p\!\in\! L^{\spfrac{2}{p-3}}([0, T], L^2(\M))\!\subset\! L^1([0, T], L^2(\M))$ because \hbox{$1\!<\!\spfrac{2}{p-3}\!<\!+\infty$}. Once the Strichartz estimates are obtained, the global well-posedness is classical for subcritical nonlinearities. It~first appeared in \cite{BLP:08} for general bounded domains. We~will need the following well-posedness result, which can be found in \cite[Th.\,2.2]{JL13}.

\begin{theorem}[Cauchy problem]\label{thm:cauchy-problem}
Let $f$ satisfies \eqref{hip:nonlinearity-hyp-1} and \eqref{defdefocusing}. For any $T>0$ and for any $(u_0, u_1)\in X=H_0^1(\M)\times L^2(\M)$, there exists a unique solution $U(t)=(u(t), \partial_t u(t))\in C^0([0, T], X)$ with finite Strichartz norm of \eqref{eq:nlw-1-cauchy}. Moreover, for any $E_0\geq 0$ and $(q, r)$ satisfying \eqref{thm:strichartz-exponents}, there exists a constant $C>0$ such that, if~$U=(u, \partial_t u)$ is solution of \eqref{eq:nlw-1-cauchy} with $E(u(0))\leq E_0$, then
\begin{align*}
\norm{u}_{L^q([0, T], L^r(\M))}\leq C\norm{(u_0, u_1)}_{H^1(\M)\times L^2(\M)}.
\end{align*}
In addition, there exists a constant $C>0$ such that, if~$U$ and $\Tilde{U}$ are two solutions of \eqref{eq:nlw-1} with $E(u(0))\leq E_0$ and $E(v(0))\leq E_0$, then
\begin{align*}
\sup_{[-T, T]}\norm{(u, \partial_t u)(t)-(\Tilde{u}, \partial_t \Tilde{u})(t)}_X\leq C\norm{(u, \partial_t u)(0)-(\Tilde{u}, \partial_t \Tilde{u})(0)}_X.
\end{align*}
\end{theorem}

\subsubsection{Observability of linear waves at higher regularity} A crucial tool in the present section will be an observability inequality of linear waves, but at an appropriate regularity. We~start by recalling the following classical observability result from Bardos-Lebeau-Rauch \cite{BLR:88} in the case $\partial\M\neq\emptyset$.

\begin{theorem}[\cite{BLR:88}]\label{thm:blr-wave}
Assume that $(\omega, T)$ satisfies \ref{assumGCC}. Then there exists $C>0$ such that, for any $(v_0, v_1)\in H_0^1(\M)\times L^2(\M)$ and associated solution $v$ of
\begin{align}\label{eq:linear-waves}
\left\{\begin{aligned}
\partial_t^2 v-\Delta_g v&=0 &&(t, x)\in [0, T]\times\M,\\
v_{|_{\partial\M}}&=0 &&(t, x)\in [0, T]\times\partial\M,\\
(v, \partial_t v)(0)&=(v_0, v_1) &&x\in \M,
\end{aligned}\right.
\end{align}
we have
\begin{align}\label{thm:ineq:linear-waves-obs}
C\norm{(v_0, v_1)}_{H_0^1(\M)\times L^2(\M)}^2\leq \int_0^T \norm{\mathbbm{1}_\omega\partial_t v(t)}_{L^2(\M)}^2 dt.
\end{align}
\end{theorem}
From now on, we~assume that $(\omega, T)$ satisfies the \ref{assumGCC} and $b_\omega$ is a smooth function defined on $\M$ such that
\begin{align}
\label{lowerb}
b_{\omega}(x)\geq 1 \quad\textnormal{for }x\in \omega.
\end{align}
Let us consider the observation operator $\bC\in \mc{L}(X^\sigma, X^\sigma)$ given by
\begin{align}\label{eq:observation-operator-wave}
\bC(\phi, \psi)=(0, b_\omega\psi).
\end{align}

\begin{proposition}\label{prop:higher-reg-obs}
Assume that $(\omega, T)$ satisfies \ref{assumGCC} and $b_{\omega}$ satisfies \eqref{lowerb}. Then for $\bC$ defined by \eqref{eq:observation-operator-wave}, there exists $\mathfrak{C}_{\textup{obs}}>0$ such that for any $\sigma\in [0,1]$ and $Z_0\in X^\sigma$, we~have
\begin{align}\label{thm:ineq:LWobshigh}
\norm{Z_0}_{X^\sigma}^2\leq \mathfrak{C}_{\textup{obs}}^2\int_0^T \norm{\bC e^{tA}Z_0}_{X^\sigma}^2 dt.
\end{align}
\end{proposition}
\begin{proof}
Let $\beta\geq 0$ and let us consider $z$ solution of
\begin{align}\label{eq:LWbeta}
\left\{\begin{aligned}
\partial_t^2 z-\Delta_g z+\beta z&=0 &&(t, x)\in [0, T]\times\M,\\
z_{|_{\partial\M}}&=0 &&(t, x)\in [0, T]\times\partial\M\ \text{ if } \partial\M\neq\emptyset,\\
(z, \partial_t z)(0)&=(z_0, z_1) &&x\in \M,
\end{aligned}\right.
\end{align}
with $(z_0, z_1)\in H_0^1(\M)\times L^2(\M)$. When $\partial\M\neq\emptyset$, from \cref{thm:blr-wave} and energy estimates, we~readily see that
\begin{align}\label{ineq:LWobsL2}
\norm{(z_0, z_1)}_{H_0^1(\M)\times L^2(\M)}^2\leq C_1\int_0^T \norm{b_\omega\partial_t z(t)}_{L^2(\M)}^2 dt.
\end{align}
If $\partial\M=\emptyset$ and $\beta>0$, considering the component $\partial_t z$, with a slight modification of \cite[Prop.\,2.2]{LL16} along with G{\aa}rding's inequality \cite[Th.\,A.9]{LL16} (see also \cite[Th.\,1.3]{CLW:20} for an equivalence with ODE along bicharacteristics, also valid for systems), we~obtain a weak observability inequality, namely, \eqref{ineq:LWobsL2} with an additional $\norm{(z_0, z_1)}_{H^{1/2}\times H^{-1/2}}$ term on the right-hand side. Such a term can easily be removed by employing a nowadays classical compactness-uniqueness argument (as in \cref{prop:obs-schr-plates} below), which leads us to \eqref{ineq:LWobsL2} in this case.

Furthermore, from \eqref{ineq:LWobsL2} along with energy estimates (commutator estimates to be more precise) allow us to deduce that there exists $C_0>0$ such that\vspace*{-3pt}
\begin{align}\label{thm:ineq:LWobsH1}
\norm{(z_0, z_1)}_{H_0^1(\M)\times L^2(\M)}^2\leq C_0\int_0^T \norm{b_\omega z(t)}_{H_0^1(\M)}^2 dt.
\end{align}
Therefore, if~$(z_0, z_1)\in (H^2(\M)\cap H_0^1(\M))\times H_0^1(\M)$, we~look at the equation satisfied by $y=\partial_t z$ and we then apply the observability inequality \eqref{thm:ineq:LWobsH1} to $y$, to obtain\vspace*{-3pt}
\begin{align}\label{ineq:LWobsH2}
\norm{(z_0, z_1)}_{(H^2(\M)\cap H_0^1(\M))\times H_0^1(\M)}^2\leq C_2\int_0^T \norm{b_\omega \partial_t z(t)}_{H_0^1(\M)}^2dt.
\end{align}
With the operator's notation of \cref{s:not_wave}, by~linear interpolation in between inequalities \eqref{ineq:LWobsL2} and \eqref{ineq:LWobsH2}, we~obtain the desired inequality \eqref{thm:ineq:LWobshigh}.
\end{proof}

\begin{remark}
We also refer to \cite{Per25} for the link between observability at different levels of regularity.
\end{remark}

The previous result shows that, if~$(\omega, T)$ satisfies the \ref{assumGCC}, then $t\mto e^{tA}$ satisfies \cref{assumC} with $\sigma\in [0,1]$, $T>0$ and observation operator $\bC$ given by \eqref{eq:observation-operator-wave}.

We also show that the pair $(A, \bC)$ satisfies \cref{assumcommu}.
\begin{proposition}\label{propcommutwave}
Let $\sigma\in [0, 1]\setminus \{1/2\}$. Then if~$\bC$ is given by \eqref{eq:observation-operator-wave} with $b_{\omega}$ smooth satisfying $\partial_{\vec n}b_{\omega}=0$, then \cref{assumcommu} is fulfilled with $s=1$ as long as $\veps\leq 1$.
\end{proposition}
\begin{proof}
We compute $[A^*A,\bC]=\begin{psmallmatrix} 0 & 0 \\ 0 & [b_{\omega},\Delta_g]\end{psmallmatrix}$ so that the result is true as long as $[b_{\omega},\Delta_g]=-2\nabla_g b_{\omega}\cdot \nabla_g -\Delta_g b_{\omega}$ sends $H^{\sigma+2}_D$ into $H^{\sigma+\veps}_D$. We~claim that $[b_{\omega},\Delta_g]$ sends $H^{\sigma+2}_D$ into $H^{\sigma+1}_D$ when $\sigma\in [0, 1]\setminus \{1/2\}$, which will give the result since $\veps\leq 1$.

Since $b_{\omega}$ is smooth and $\sigma\in [0, 1]\setminus \{1/2\}$, we~only need to verify that the term $\nabla b_{\omega}\cdot \nabla u$ is indeed zero on $\partial\M$. Decomposing $\nabla_g =\partial_{\vec n}+\nabla_{T}$, where $\nabla_{T}$ is the tangential derivative, we~can write $\nabla b_{\omega}\cdot \nabla u=\partial_{\vec n}b_{\omega} \partial_{\vec n}u+\nabla_T b_{\omega}\cdot \nabla_T u$. The first term is zero on the boundary since we assumed $\partial_{\vec n}b_{\omega}=0$, while the second term $\nabla_T u$ cancels since $u=0$ on $\partial \M$. This proves the claimed result.
\end{proof}

\subsection{Propagation of analyticity}
Here we prove \cref{thm:analytic-prop} and then Corollary~\ref{coranalytspacetime}.

\begin{proof}[Proof of \cref{thm:analytic-prop}] According to \cref{lemma:gcc-smaller-subset}, there exists $\chi\in C^{\infty}_c(\omega)$ with non negative values such that
\begin{itemize}
\item there exists a nonempty open set $\Tilde{\omega}$ which is compactly contained in $\omega$ and a time $\Tilde T\in(0,T)$ such that $(\Tilde{\omega}, \Tilde{T})$ satisfies the \ref{assumGCC} and $\chi =1$ on $\Tilde{\omega}$,
\item $\partial_{\vec n} \chi=0$ on $\partial \M$ where $\partial_{\vec n}$ is the normal derivative to the boundary.
\end{itemize}

Let $\chi_2 \in C^{\infty}_c(\omega)$ be such that $\chi_2 =1$ on $\supp(\chi)$. Since $\chi_2$ is a cutoff function whose support is contained on $\omega$, by~hypothesis, we~get that $t\mto \chi_2u(t)$ is analytic with value in $ H^{1+\sigma}\cap H_0^1(\M)$. Note that, up to exchanging $(0,T)$ with a compact subinterval such that the geometric control condition is still satisfied, we~can assume without loss of generality that the analyticity holds in a neighborhood of $(0,T)$. In~the same way as done in the proof of \cref{propcommutwave}, we~verify that the map $[\Delta_g, \chi]$ maps $ H^{1+\sigma}\cap H_0^1(\M)$ into $H^{\sigma}_0(\M)$ when $\sigma\in [0, 1]\setminus \{1/2\}$. By~composition of analytic maps, and noticing that $[\Delta_g, \chi]=\chi_2[\Delta_g, \chi]=[\Delta_g, \chi]\chi_2$, we~get that the map\vspace*{-3pt}
\begin{align*}
t\in (0, T)\mto [\Delta_g, \chi] u\in H_0^{\sigma}(\M)
\end{align*}
is analytic.

By writing $u=\chi u+(1-\chi)u$, it~remains to show that $t\mto (1-\chi)u$ is analytic. Set $\Tilde{\chi}=(1-\chi)$ and consider the new variable $z=\Tilde{\chi} u$. We~then have\vspace*{-3pt}
\begin{align*}
\partial_t^2 z-\Delta_g z+\beta z&=\Tilde{\chi}(\partial_t^2 u-\Delta_g u)+\beta z+[\Delta_g, \Tilde{\chi}]u\\
&=-\Tilde{\chi} f(u)+\beta z-[\Delta_g, \chi]u=-\Tilde{\chi} f(z+\chi u)+\beta z-[\Delta_g, \chi]u\\
&=-\Tilde{\chi} f(z+h_1)+\beta (z+h_1)+h_2-\beta h_1,
\end{align*}
where we have defined the functions $h_1=\chi u$ and $h_2=-[\Delta_g, \chi]u$.

Let $\sigma_*\in (1/2, \sigma)$ if $d=3$ and $\sigma_*\in (0,\sigma)$ if $d\leq 2$. We now~define\vspace*{-3pt}
\begin{align*}
\left\{\begin{array}{ccl}
t\in [0, T] & \dpl\mto & H_1(t)=(\chi u(t), 0)\in X^{\sigma_*}, \\[3pt]
t\in [0, T] & \dpl\mto & H_2(t)=(0, -[\Delta_g, \chi]u(t)-\beta \chi u(t))\in X^\sigma.
\end{array}\right.
\end{align*}
We observe that $H_1\in C^0([0, T], X^{\sigma_*})$ and $H_2\in C^0([0, T], X^\sigma)$. Since $U$ is bounded in $C^0([0, T], X^\sigma)$ by some $M>0$, it~follows that\vspace*{-3pt}
\begin{align*}
H_1\in \B_{CM}^{[0, T]}(X^{\sigma_*})\ \text{ and }\ H_2\in \B_{CM}^{[0, T]}(X^\sigma),
\end{align*}
for some $C=C(\chi,\chi_2)=C(\omega)>0$. Since $t\in (0, T)\mto \chi u\in H^{1+\sigma_*}(\M)\cap H_0^1(\M)$ and $t\in (0, T)\mto [\Delta, \chi] u\in H^{\sigma}_0(\M)$ are analytic, an application of \cite[Th.\,7.2]{BS:71-analytic} and compactness, gives the existence of $\mu>0$ such that $H_1$ and $H_2$ can be extended holomorphically as
\begin{align*}
\left\{\begin{array}{ccl}
z\in (0, T)+i(-\mu, \mu) & \dpl\mto & H_1(z)=(\chi u(z), 0)\in X_\mathbb{C}^{\sigma_*}, \\[3pt]
z\in (0, T)+i(-\mu, \mu) & \dpl\mto & H_2(z)=(0, -[\Delta_g, \chi]u(z)-\beta \chi u(z))\in X_\mathbb{C}^\sigma.
\end{array}\right.
\end{align*}
Moreover, since $\Re(H_1(z))\in \mathbb{B}_{CM}(X^{\sigma_*})$ for $z\in[0,T]+i\{0\}$, shrinking $\mu>0$ if necessary, by~continuity and compactness, we~can assume that $\Re(H_1(z))\in \mathbb{B}_{2CM}(X^{\sigma_*})$ for every $z\in [0, T]+i[-\mu, \mu]$.

Applying \cref{lemma:gcc-smaller-subset} again to $(\Tilde{\omega},\Tilde{T})$, we~can find $(\widetilde{\omega}_1,\Tilde{T}_1)$ satisfying \ref{assumGCC} and $b\in C^{\infty}_c(\Tilde{\omega})$ such that $b=1$ on $\widetilde{\omega}_1$ and $0<\widetilde{T}_1<\Tilde{T}<T$. Defining $\bC$ by $\bC(\phi, \psi)=(0, b\psi)$ as in \eqref{eq:observation-operator-wave}, \cref{prop:higher-reg-obs} gives the following observability estimate\vspace*{-3pt}
\begin{align}\label{thm:ineq:LWobshighTstar}
\norm{Z_0}_{X^{\sigma_*}}^2\leq \mathfrak{C}_{\textup{obs}}^2\int_0^{\widetilde{T}_1} \norm{\bC e^{tA}Z_0}_{X^{\sigma_*}}^2 dt
\end{align}
and the same observability at the level of regularity $X^{\sigma}$.

Observe that $\chi=1$ on $\Tilde{\omega}$, then $z=(1-\chi)u=0$ on $\widetilde{\omega}$. Then since $b\in C^{\infty}_c(\Tilde{\omega})$, it~implies $b\partial_t z=0$ on $[0,T]\times \M$.
Finally, we~see that $Z=(z, \partial_t z)$ satisfies the system
\begin{equation}\label{thm:proof:eq:}
\left\{\begin{aligned}
\partial_t Z&=AZ+F(Z+H_{1})+H_{2}&&\textnormal{ on } [0,T], \\
\bC Z(t)&=0 &&\textnormal{ for }t\in [0,T],
\end{aligned}\right.
\end{equation}
where $F$ is as in \eqref{defFwave}. We~are then in the framework of \cref{thmabstractanalyticintro} with $(T^*, T)$, $(\sigma, \veps)$ and $R_0$ of the theorem replaced by $(T, \widetilde{T}_1)$, $(\sigma_*, \sigma-\sigma_*)$ and $2CM$, respectively. \cref{assumFholom} is fulfilled thanks to \cref{prop:f-assumptions}. The observability estimate \eqref{thm:ineq:LWobshighTstar} for $\sigma$ and $\sigma_*$ ensures \cref{assumCC}. We~deduce that $t\in (0, T)\mto Z(t)\in X^{\sigma_*}$ is real analytic, hence $t\in (0, T)\mto U(t, \cdot)\in X^{\sigma_*}$ is real analytic as well.

Observe that $t\in (0, T)\mto \partial_t U(t, \cdot)\in X^{\sigma_*}$ is an analytic map (see \cite[Prop.\,6.4]{BS:71-analytic}), and so is $t\in (0, T)\mto \partial_t U(t, \cdot)-F(U(t,\cdot))\in X^{\sigma_*}$, where $F(u,v)=(0,-f(u))$ (that is, as in \eqref{defFwave} with $\widetilde{\chi}=1$ and $\beta=0$). We~readily get that the map $t\in (0, T)\mto A U(t, \cdot)\in X^{\sigma_*}$ is analytic. This implies that $t\in (0, T)\mto U(t, \cdot)\in X^{1+\sigma_*}$ is analytic as well.

For the last statement, we~can use local holomorphic extension and prove that it is weakly holomorphic, which is sufficient (see \cite[Th.\,8.12]{Muj86}). Indeed, a continuous linear form $L$ on $X^{1+\sigma_*}$ can be written, for some $V\in X^{1+\sigma_*}$,\vspace*{-3pt}
\begin{align*}
L(U)=\left<V,U\right>_{X^{1+\sigma_*}}&=\bigl<(A^*A)^{1/2}V,(A^*A)^{1/2}U\bigr>_{X^{\sigma_*}}+\left<V,U\right>_{X^{\sigma_*}}\\
&=\left<AV,AU\right>_{X^{\sigma_*}}+\left<V,U\right>_{X^{\sigma_*}}.
\end{align*}
In~particular, since $AV\in X^{\sigma_*}$, the extension $z\mto L(U(z))$ is well-defined and holomorphic.
\end{proof}
\begin{remark}\label{rk:comm-reg}
In the previous proof, the restriction $d\leq 3$ has been used mainly in the fact that the space $H^{\sigma}_D$ for $\sigma\in (d/2-1, d/2-1/2)$ only involves the trace of~$u$ on $\partial \M$, and not on that of $\Delta_g u$. Otherwise, the action of the maps $u \mto \chi u$ and $u \mto [\Delta_g, \chi ]u$ is more involved. This subtlety already appeared in \cite{DL:09} concerning the action of the HUM operator on regular data. This problems can be avoided for instance if $\mathbbm{1}_{\omega}$ is constant near any connected component of the boundary. This allows to select some cutoff functions that are constant near the boundary for which the compatibility conditions disappear.
\end{remark}

\begin{proof}[Proof of Corollary \ref{coranalytspacetime}]
We know that $u$ admits a holomorphic extension to $[0,T]+i[-\delta,\delta]$ with value in $H^{1+\sigma}(\M)\cap H^1_0(\M)$ with $\sigma\in (1/2,1]$. Let $\varphi\in C^{\infty}_c(\M)$. For any $z=(t,s)\in [0,T]+i[-\delta,\delta]$, we~consider the well-defined quantity $m(z)=\left<\partial_z^2u(z)-f(u(z)),\varphi\right>_{L^2(\M)}-\left<u(z),\Delta_g \varphi\right>_{L^2(\M)}$. Note that $m$ is holomorphic by composition. Since $u$ is holomorphic, it~satisfies $\partial_{\bar z}u(z)=0$ for $z\in (0,T)+i(-\delta,\delta)$ where $\partial_{\bar z}=\frac{1}{2}\left(\partial_t+i\partial_s\right)$. In~particular, $\partial_{z}u(z)=\partial_t u(z)=\Psfrac{1}{i}\partial_s u(z)$ with equality meant in $H^{1+\sigma}\cap H^1_0$. Moreover, if~we restrict $m$ to the real interval $(0,T)$, we~know that it is zero since $u_{\left|(0,T)\right.}$ is solution of \eqref{eq:nlw-1}. By~analytic continuation, we~conclude that $m(z)=0$ for $z\in (0,T)+i(-\delta,\delta)$. That means that $\partial_z^2 \left<u(z),\varphi\right>_{L^2(\M)}=\left<f(u(z)),\varphi\right>_{L^2(\M)}+\left<u(z),\Delta_g \varphi\right>_{L^2(\M)}$ for any $\varphi\in C^{\infty}_c(\M)$ and $z\in (0,T)+i(-\delta,\delta)$.

For any $t_0\in (0,T)$, we~consider the function $v\in C^{\infty}((-\delta,\delta),H^{1+\sigma}\cap H^1_0(\M))$ defined by $s\in (-\delta,\delta)\mto v(s):=u(t_0+is)\in H^{1+\sigma}\cap H^1_0(\M)$. For any $\varphi\in C^{\infty}_c(\M)$, we~compute\vspace*{-3pt}
\begin{align*}
\partial_s^2 \left<v(s),\varphi\right>_{L^2(\M)}&=\left<\partial_s^2 v(s),\varphi\right>_{L^2(\M)}=-\left<\partial_z^2 u(t_0+is),\varphi\right>_{L^2(\M)}\\
&=-\left<f(u(t_0+is)),\varphi\right>_{L^2(\M)}-\left<u(t_0+is),\Delta_g \varphi\right>_{L^2(\M)}\\
&=-\left<f(v(s)),\varphi\right>_{L^2(\M)}-\langle v(s),\Delta_g \varphi\rangle_{L^2(\M)}.
\end{align*}
In particular, $v$ is solution, in~the distributional sense of $\partial_s^2v+\Delta_g v=-f(v)$ on $(-\delta,\delta)\times \M$ with Dirichlet boundary condition on $\partial \M$. Since
\[
v\in C^{\infty}((-\delta, \delta),H^{1+\sigma}\cap H^1_0(\M))\subset C^0((-\delta,\delta)\times \M),
\]
standard elliptic regularity states that $v$ is smooth and therefore analytic, see for instance Friedman \cite[Th.\,5]{Friedman:58}. In~particular, it~gives that for any $x_0\in \M$, in~some charts around $x_0$, there exists $R>0$ and $C>0$ such that
\begin{align*}
\left| \partial_s^{\alpha}\partial_x^{\beta}v(0,x_0)\right|\leq C R^{\alpha+|\beta|}\alpha!\beta!.
\end{align*}
By definition of $u$ and holomorphy with value in $H^{1+\sigma}\cap H^1_0$, we~have $\partial_s^{\alpha}v(0)=(i\partial_t)^{\alpha}u(t_0)$ with equality in $H^{1+\sigma}\cap H^1_0$. Since the valuation at $x_0\in \M$ is continuous on $H^{1+\sigma}\cap H^1_0$, we~get $\partial_s^{\alpha}v(0,x_0)=(i\partial_t)^{\alpha}u(t_0,x_0)$ for any $x_0\in \M$. Taking now derivative in $x_0$, we~have $\partial_s^{\alpha}\partial_x^{\beta}v(0,x_0)=(i\partial_t)^{\alpha}\partial_x^{\beta}u(t_0,x_0)$. So, we~obtain
\begin{align*}
\left| \partial_t^{\alpha}\partial_x^{\beta}u(t_0,x_0)\right|\leq C R^{\alpha+|\beta|}\alpha!\beta!.
\end{align*}
This means analyticity close to $(t_0,x_0)$ and gives the result since $(t_0,x_0)\in (0,T)\times \M$ are arbitrary.
\end{proof}

\subsection{Finite determining modes} We now show that the property of finite determining modes holds the observed problem \eqref{eq:nlw-1}.
\begin{proposition}
\label{propfinitedetwave}
Let $\sigma\in (1/2, 1]$ if $d=3$ and $\sigma\in (0,1/2)$ if $d\leq 2$. With the notations of \cref{s:not_wave}, assume $(\omega, T)$ satisfies \ref{assumGCC} and that $f\in C^4(\R, \R)$. For any $R_0>0$, there exists $n\in \N$ such that the following holds. Let $h\in C^0([0,T], H_0^{\sigma}(\M))$ and $g\in L^2([0, T], H_0^\sigma(\M))$. Let $U(t)=(u(t),\partial_t u(t))$ and $\widetilde{U}(t)=(\tilde{u}(t), \partial_t\tilde{u}(t))$ be two solutions on $(0,T)$ of\vspace*{-3pt}
\begin{align}\label{eq:nlw-uc_source}
\left\{\begin{aligned}
\partial_t^2 u-\Delta_g u+f(u)&=h(t, x) &&(t, x)\in [0, T]\times \text{Int}(\M),\\
u_{|_{\partial\M}}&=0 &&(t, x)\in [0, T]\times\partial\M,\\
\partial_t u&=g &&(t, x)\in [0, T]\times \omega,
\end{aligned}\right.
\end{align}
such that $\norm{U(t)}_{X^\sigma}\leq R_{0}$ and $\norm{\widetilde{U}(t)}_{X^\sigma}\leq R_{0}$ for all $t\in [0, T]$. If $\P_n U(t)=\P_n \widetilde{U}(t)$ for all times $t\in [0, T]$. Then $U(t)\equiv \widetilde{U}(t)$ for all $t\in [0, T]$.
\end{proposition}
\begin{proof}
We have already established in \cref{s:not_wave} that $A$ satisfies \cref{assumAA}. Under the hypothesis that $(\omega, T)$ satisfies the \ref{assumGCC}, \cref{prop:higher-reg-obs} implies \cref{assumC}. Finally, from \cref{prop:f-assumptions}, \cref{assumF} is satisfied for some $\veps\in (0, 1]$, where $F$ is defined by \eqref{defFwave} with $\widetilde{\chi}=1$. Then the result follows, as a direct application of the abstract \cref{prop:finite-det-modes}.
\end{proof}

\subsection{Unique continuation and equilibrium points} Here we will consider $U=(u, \partial_t u)\in C^0([0, T], X)$ solution of the system\vspace*{-3pt}
\begin{align}\label{eq:nlw-uc}
\left\{\begin{aligned}
\partial_t^2 u-\Delta_g u+f(u)&=0 &&(t, x)\in [0, T]\times \text{Int}(\M),\\
u_{|_{\partial\M}}&=0 &&(t, x)\in [0, T]\times\partial\M,\\
\partial_t u&=0 &&(t, x)\in [0, T]\times \omega.
\end{aligned}\right.
\end{align}
The purpose of this section is to prove \cref{thm:unique-continuation-nlw}.

\subsubsection{Propagation of regularity} Here we aim to prove a regularization effect for functions satisfying \eqref{eq:nlw-uc} in the subcritical case. When $\M$ is of dimension $d=2$ and~$f$ has polynomial growth with $p\in [1, +\infty)$, it~is possible to prove a gain of regularity in the nonlinearity by means of appropriate Sobolev embedding. The same property holds when $d=3$ and $p\in [1, 3]$. However, when $d=3$ and $p\in (3, 5)$, this gain of regularity needs to be handled by means of Strichartz estimates. Dehman-Lebeau-Zuazua \cite[Th.\,8]{DLZ03} proved that the nonlinearity term is more regular than it seems to be and they use this fact in a key fashion to stabilize the semilinear wave equation in an unbounded domain. We~recall such a regularity result from \cite[Cor.\,4.2]{JL13}, which in particular fits with our geometrical setting.

\begin{proposition}\label{prop:gain-reg-nonlinearity}
Assume $d=3$. Let $R>0$ and $T>0$. Let $s\in [0, 1)$ and $\veps=\min\{1-s, (5-p)/2, (17-3p)/14\}>0$ with $p$ as in \eqref{hip:nonlinearity-hyp-1}. There exist $(q, r)$ satisfying~\eqref{thm:strichartz-exponents} and $C>0$ such that the following property holds. If
\[
v\in L^\infty([0, T], H^{1+s}(\M)\cap H_0^1(\M))
\]
is a function with finite Strichartz norms $\norm{v}_{L^q([0, T], L^r(\M))}\leq R$, then we have $f(v)\in L^1([0, T], H_0^{s+\veps}(\M))$ and, moreover,
\begin{align*}
\norm{f(v)}_{L^1([0, T], H_0^{s+\veps}(\M))}\leq C\norm{v}_{L^\infty([0, T], H^{1+s}(\M)\cap H_0^1(\M))}.
\end{align*}
The constant $C$ depends only on $\M$, $(q, r)$, $R$ and the constant in estimate \eqref{hip:nonlinearity-hyp-1}.
\end{proposition}
\begin{remark}\label{rk:gain-reg-cpct-pert}
According to the geometric framework considered in \cite{JL13}, the previous result remains true when $\M$ is a compact perturbation of $\R^3$, that is, $\R^3\setminus\O$ where $\O$ is a bounded smooth domain, endowed with a smooth metric equal to the euclidean one outside of a ball.
\end{remark}

With this gain of regularity at hand, along with the assumption that $(\omega, T)$ satisfies the \ref{assumGCC}, we~are able to propagate the regularity through the observability estimate \eqref{thm:ineq:LWobshigh}.

\begin{proposition}\label{prop:propagation-regularity}
Let $U\in C^0([0, T], X)$ be a mild solution of the system \eqref{eq:nlw-uc} with finite Strichartz norm. Then
\begin{align*}
U\in C^0([0, T], X^{\upsilon})
\end{align*}
for all $\upsilon\in [0, 1)$. In~particular $u\in L^\infty([0, T]\times \M)$.
\end{proposition}

\begin{proof}
The proof is the same as that of \cref{prop:abs-prop-reg}, except that the nonlinear term is not bounded. Yet, the Duhamel term is well-defined thanks to Sobolev embedding or Strichartz estimates depending on the case; see Remark \ref{rk:abs-prop-reg}. Therefore, we~only need to check that $\T: t\mto \int_0^t e^{(t-s)A}F(U(s))ds$ defines a bounded map from $L^\infty([0, T], X)$ into $L^\infty([0, T], X^\veps)$ for some $\veps>0$.

From \cite[Prop.\,3.2]{2020:joly-laurent:decay-nlw-no-gcc}, when $d=2$, we~know that $F$ maps bounded sets of $X$ into bounded sets of $X^\veps$ for any $\veps\in [0, 1)$. When $d=3$ and $p\in [1, 3]$, the same property holds with $\veps\in [0, (3-p)/2)$. When $d=3$ and $p\in (3, 5)$, since $u\in H_0^1(\M)$ and $f$ is subcritical, \cref{prop:gain-reg-nonlinearity} implies that, for a given $\veps>0$ depending on $p$, $f(u)$~is globally bounded in $L^1([0, T], H_0^{\veps}(\M))$. In~any case, $\T$ is bounded in $C^0([0, T], X^\veps)$ for an appropriate choice of $\veps>0$.

Without loss of generality, let us now fix $\veps>0$ small enough so we can encompass all the aforementioned cases simultaneously. For $d=2$ and $d=3$, we~deduce from \cref{prop:higher-reg-obs} that \cref{assumCC} is satisfied for $\sigma=0$ and $\veps>0$. Let $b_\omega$ be given by \cref{lemma:gcc-smaller-subset} and let $\bC$ be as in \eqref{eq:observation-operator-wave}. Then \cref{propcommutwave} ensures that the pair $(A, \bC)$ satisfies \cref{assumcommu} with $s=1$ if $\veps\leq 1$. We~then reproduce the proof of \cref{prop:abs-prop-reg} to obtain that $U$ is bounded in $C^0([0, T], X^\veps)$.

We can iterate the previous process to obtain that $U$ is bounded in $C^0([0, T], X^{k\veps})$ for $k\in \N$ as long as \cref{prop:gain-reg-nonlinearity}, or simply \cref{prop:f-assumptions}, applies. We~finally obtain that $U$ is bounded in $C^0([0, T], X^\upsilon)$ for any $\upsilon\in [0, 1)$.
\end{proof}

\begin{remark}
\label{rk:propagsource}
\cref{prop:propagation-regularity} can easily be extended with the same proof by replacing \eqref{eq:nlw-uc} with equations of the form $\partial_t^2 u-\Delta_g u+\chi(x)f(u+h_1)=h_2$ with $h_1\in C^0([0,T],H^2_D)$ and $h_1\in C^0([0,T],H^1_D)$.
\end{remark}

\subsubsection{On unique continuation for linear waves}\label{s:UCPword} Based on the works of Tataru, Robbiano-Zuily and Hörmander, it~is established that global unique continuation holds under the framework of partial analyticity and very general geometric assumptions, provided sufficient time has passed, so that we do not contradict the finite speed of propagation. To obtain the nonlinear unique continuation property, we~will follow~\cite{JL13} and aim to treat the nonlinearity as a potential term. Therefore, we~must ensure that our framework allows for the application of the result for linear waves. In~what follows, we~will recall a version of the Robbiano-Zuily-Hörmander-Tataru result, which is well suited to our specific context.

We will first introduce some geometric quantities needed to state the result. For $E\subset \M$, we~can define the largest distance from $E$ to a point in $\M$ by\vspace*{-3pt}
\begin{align*}
\mc{L}(\M, E)=\sup_{x\in\M}\text{dist}(x, E).
\end{align*}
If $E$ is open, the quantity\vspace*{-3pt}
\begin{align*}
T_{\textup{UC}}(E):=2\mc{L}(\M, E)
\end{align*}
is the minimal time of unique continuation for the (linear) wave equation from an open set $E$, see Tataru \cite{Tat95}. We~say that an open set $\omega$ satisfies the \ref{assumGCC} if there exists $T>0$ such that $(\omega, T)$ satisfies the \ref{assumGCC}. One can then define the minimal control time associated with $\omega$ by\vspace*{-3pt}
\begin{align*}
T_{\textup{GCC}}(\omega)=\inf\{T>0\ |\ (\omega, T)\ \text{satisfies \ref{assumGCC}}\}.
\end{align*}
It can be proved that $T_{\textup{GCC}}(\omega)\geq T_{\textup{UC}}(\omega)$, see Laurent-Léautaud \cite[Lem.\,B.4]{LL16}.

\begin{remark}
For the case where $\partial\M=\emptyset$, the critical time $T_{\textup{GCC}}(\omega)$ is not allowed, since, as shown in \cite[Th.\,1.1]{LL16}, the observability estimate always fails for such time.
\end{remark}

We now state the unique continuation property for linear waves with coefficients analytic in time, due to Tataru \cite{Tat95,Tat99}, Robbiano-Zuily \cite{RZ:98} and Hörmander \cite{Hor:97}. We~refer to \cite[Th.\,6.1]{LL19} for a quantitative statement that implies the unique continuation and contains the construction of hypersurfaces that allows to obtain the global result.

\begin{theorem}[Tataru-Robbiano-Zuily-Hörmander]\label{thm:qucp}
Let $\M$ be a compact Riemannian manifold with (or without) boundary, $\Delta_g$ the Laplace-Beltrami operator on~$\M$, and
\begin{align*}
P=\partial_t^2-\Delta_g+W_0\partial_t+W_1\cdot\nabla+V
\end{align*}
with $V$, $W_0$, $W_1$, $\text{div}(W_1)$ bounded and depending analytically on the variable $t\in (0, T)$. Let $\omega$ be a nonempty open subset of $\M$ and $T>2\mc{L}(\M, \omega)$. Let $(u_0, u_1)\in H_0^1(\M)\times L^2(\M)$ and associated solution $u$ of
\begin{align*}
\left\{\begin{aligned}
Pu&=0&&\text{ in }\ (0, T)\times\Int{\M},\\
u_{|_{\partial\M}}&=0 &&\text{ in }\ (0, T)\times \partial\M,\\
(u, \partial_t u)(0)&=(u_0, u_1). &
\end{aligned} \right.
\end{align*}
Then if~$u$ satisfies $u=0$ on $[0,T]\times \omega$, then $u=0$ on $[0,T]\times \M$.
\end{theorem}

\subsubsection{Unique continuation for semilinear waves} We now prove \cref{thm:unique-continuation-nlw} and then \cref{propR3obstacle}.

\begin{proof}[Proof of \cref{thm:unique-continuation-nlw}] Let $U(t)=(u(t), \partial_t u(t))$ be a solution of \eqref{eq:nlw-uc}, which, by~\cref{prop:propagation-regularity}, belongs to $L^\infty([0, T], X^\sigma)$ for any $\sigma\in (1/2, 1)$. For any $\chi\in C^{\infty}_c(\omega) $, the map
\[
t\in (0, T)\mto \chi u(t,\cdot)\in H^{1+\sigma}(\M)\cap H_0^1(\M)
\]
does not depend on $t$ and is therefore analytic. In~particular, \cref{thm:analytic-prop} applies and we get that $t\in (0, T)\mto U(t)\in X^\sigma$ is analytic.

Since $u$ is smooth with respect to $t$ and $f$ is smooth, by~writing
\begin{align*}
\Delta_g u=\partial_t^2 u+f(u),
\end{align*}
we get that $\Delta_g u\in L^2(\M)$ and so $u\in H^2(\M)$. We~differentiate the above equation to obtain
\begin{align*}
\Delta_g^2 u= \Delta_g (\partial_t^2 u+f(u))&=\partial_t^2 \Delta_g u+\Delta_g f(u)\\
&=\partial_t^4 u+f'(u)\partial_t^2 u+f''(u)(\partial_t^2 u)^2+\Delta_g f(u),
\end{align*}
which shows that $u$ belongs to $H^4(\M)$. This process can be repeated as many times as wanted, so by classical Sobolev embedding we get that $u=u(t, x)$ is smooth with respect to $x$. In~particular, $(t, x)\in [0, T]\times \M\mto u(t, x)\in C^\infty([0, T]\times \M)$ is bounded, together with all its derivatives.

Set $z=\partial_t u$ and observe that $z$ solves
\begin{align}\label{thm:proof:eq:nlw-unique-cont}
\left\{\begin{aligned}
\partial_t^2 z-\Delta_g z+f'(u)z&=0 &&(t, x)\in [0, T]\times \text{Int}(\M),\\
z_{\left|\partial\M\right.}&=0 &&(t, x)\in [0, T]\times\partial\M,\\
z&=0 &&(t, x)\in [0, T]\times \omega.
\end{aligned}\right.
\end{align}
By the previous discussion $(t, x)\mto f'(u(t, x))$ is bounded, analytic in $t$ and smooth in~$x$. Since $T>T_{\textup{GCC}}\geq T_{\textup{UC}}=2\mc{L}(\M, \omega)$, we~can apply \cref{thm:qucp} to get that~$z\equiv 0$ everywhere. This in turn means that $u(t, x)=u(x)$ is constant in time and henceforth it solves
\begin{align}
\left\{\begin{aligned}
-\Delta_g u+f(u)&=0 &&x\in\text{Int}(\M),\\
u&=0 &&x\in \partial\M.
\end{aligned}\right.
\end{align}
Moreover, multiplying the latter equation by $u$ and integrating by parts leads us to the identity
\begin{align*}
0\leq \int_\M |\nabla u(x)|^2 dx=-\int_\M u(x)f(u(x)) dx.
\end{align*}
Under the assumption that $sf(s)\geq 0$ for all $s$, the above identity, the connectedness of $\M$ and the boundary condition imply that $u\equiv 0$ everywhere. In~the case $\partial\M=\emptyset$, we~get
\begin{align*}
0\leq \int_\M |\nabla u(x)|^2 dx=-\int_\M u(x)f(u(x)) dx\leq -\gamma\int_\M |u(x)|^2dx,
\end{align*}
which directly implies $u\equiv 0$ everywhere.
\end{proof}

\begin{proof}[Proof of \cref{propR3obstacle}] Up to increasing $R$, we~can assume without loss of generality that $\O\Subset B(0,R)$. Since $\partial_t u$ vanishes on $(0, T)\times\big(\R^3\setminus B(0,R)\big)\subset (0, T)\times \omega $, we~have $-\Delta u+f(u)=0$ with $u\in H^1(\R^3\setminus B(0, R))$. Since $f$ is subcritical, we~can use elliptic regularity and bootstrap (see, for instance, \cite[Th.\,9.19]{GT01}) to show that $u=u(x)$ belongs to $C^{4, \al}$ for some $\al\in (0, 1)$ in the set $\R^3\setminus B(0, R)$.

Let us take $R_1>R$. We~can easily construct a radial cutoff function $\chi\in C_c^\infty(\overline{\Omega})$ satisfying:
\begin{itemize}
\item $\chi=1$ in $B(0, R)\setminus \O$;
\item $\chi=0$ in $\Omega\setminus B(0, R_1)$;
\item $\partial_{\vec n}\chi=0$ in $\partial\Omega=\partial\O$;
\item $\supp\nabla\chi\subset B(0, R_1)\setminus B(0, R)$.
\end{itemize}
Observe that $(1-\chi)$ is supported in $\R^3\setminus B(0, R)$ and that $u$ does not depend on time in such a set and is regular in space. Henceforth, $(1-\chi) u$ is analytic as a map from $(0, T)$ into $H^{1+\sigma}(\Omega)\cap H_0^1(\Omega)$, for some $\sigma\in (1/2, 1)$, which we fix from now on. Additionally, pick $\veps>0$ such that $\sigma+\veps<1$. It~remains to check that $\chi u$ is analytic on the time variable $t$.

Let us take $\widetilde{R}>R_1$ and consider $\M=B(0,\widetilde{R})\setminus \O$, such that $\M\cap\omega\neq \emptyset$ and $\partial \M=\partial \O\cup S(0,\widetilde{R})$. Since $\chi$ is $0$ outside the ball $B(0, R_1)$, we~can consider it as a cutoff function in $C_c^\infty(\overline{\M})$ with $\partial_{\vec n}\chi=0$ on $\partial\M$. Let $\widetilde{\chi}\in C_c^\infty(\overline{\Omega})$ be another cutoff with the same properties as $\chi$ and $\widetilde{\chi}=1$ on $\textnormal{Supp}(\chi)$. Since the operator $f$ is local, we~have $\chi f(u)=\chi f(\widetilde{\chi}u)$. The equation satisfied by $z:=\chi u$ is then
\begin{align}
\label{eq:zobstacle}
\left\{\begin{aligned}
\partial_t^2 z-\Delta_g z+\chi f(z+(1-\chi)\widetilde{\chi}u)-[\Delta_g, \chi]u&=0 &&\text{on }(0, T)\times \M, \\
z_{|_{\partial \M}}&=0 &&\text{on }(0, T)\times \partial\M, \\
\partial_t z&=0 &&\text{on }(0, T)\times \widetilde{\omega},\\
\end{aligned}\right.
\end{align}
where $\widetilde{\omega}=\M\cap\omega$. Observe that $[\Delta_g, \chi]$ is supported in the annulus
\[
B(0, R_1)\setminus B(0, R)\subset \widetilde{\omega}.
\]
Moreover, due to the regularity of $u$ in such a set (recall that it does not depend on time there) and the condition $\partial_{\vec n} \chi=0$ in $\partial\M$, $[\Delta_g, \chi]u$ is an analytic map from $(0, T)$ into $H_D^{2}(\M)\subset H_D^{\sigma+\veps}(\M)$. The same holds for $(1-\chi)\widetilde{\chi}u$ that defines an analytic map from $(0, T)$ into $H_D^{2}(\M)\subset H_D^{1+\sigma}$, where the cutoff $\widetilde{\chi}$ ensures the correct boundary condition on $S(0,\widetilde{R})$. In~particular, since Proposition \eqref{prop:propagation-regularity} still holds for the equation \eqref{eq:zobstacle} (see Remark \ref{rk:propagsource}), we~obtain that $(z, \partial_t z)\in C^0([0, T], X^\sigma)$.

We are then in the configuration of \cref{thmabstractanalyticintro}. As we did in the proof of \cref{thm:analytic-prop}, we~get that $t\mto z(t)\in H^{1+\sigma}(\M)\cap H_0^1(\M)$ is real analytic.

Summarizing, we~have proved that $t\in (0, t)\mto u(t)\in H^{1+\sigma}(\Omega)\cap H_0^1(\Omega)$ is analytic. A version of \cref{thm:qucp} of unique continuation for linear waves for unbounded domains can be applied (see \cite[Cor.\,3.12]{JL13}), from which we get that $\partial_t u=0$ in the whole cylinder $(0, T)\times \Omega$. We~then conclude as we did for \cref{thm:unique-continuation-nlw}.
\end{proof}

\subsection{Observability inequality for the nonlinear equation}
Once the unique continuation property has been proved, the proof of the observability estimate follows some ideas from earlier articles. We~follow in particular, the scheme introduced in~\cite{DLZ03} with the further simplification of \cite{JL13} that replaced the use of microlocal defect measure by the decay of the semigroup. We~follow a similar path, except that we want to have the observability in finite time, the one of \ref{assumGCC}. Therefore, we~have to use instead the observability estimate as a black box.
\begin{proof}[Proof of \cref{thmobserintro}]We make the proof for $d=3$ and, without loss of generality, we~assume that $p\in (3,5)$. For $d\leq 2$, the proof is the same using different Strichartz norms and the polynomial bound of the nonlinearity. To simplify the notation, we~denote $\mc{Z}$ the Banach space of vectors $W=(u,v)$ such that $W\in\nobreak C^0([0, T],X)$ and $u\in\nobreak L^{4}([0, T],L^{12}(\M))$, endowed with the natural norm. We~argue by contradiction. Assume that \eqref{obsevNLintro} is not satisfied. Then there exists a sequence \hbox{$(u_{0,n},u_{1,n})\in H^1_0\times L^2(\M)$}, with $\nor{(u_{0,n},u_{1,n})}{H^1_0\times L^2}\leq R_0$, such that the unique solution $(u_{n}, \partial_t u_{n})\in \mc{Z}$ of \eqref{eq:nlw-1} satisfies
\begin{align}\label{obsevNLintronegation}
\int_0^T \norm{\mathbbm{1}_\omega\partial_t u_{n}(t)}_{L^2(\M)}^2 dt
\leq \frac{1}{n}\norm{(u_{0,n},u_{1,n})}_{H_0^1(\M)\times L^2(\M)}^2.
\end{align}
Up to taking a subsequence, we~can assume $(u_{0,n},u_{1,n}) $ converges weakly to $(u_{0},u_{1})\in H^1_0\times L^2$. The global well-posedness theory for \eqref{eq:nlw-1} and the defocusing assumption states that the sequence $u_{n}$ is globally bounded in the space $C^{0}([0, T], H_0^1(\M))$ with uniformly bounded Strichartz norms $L^{4}([0, T], L^{12}(\M))\leq R$; see \cref{thm:cauchy-problem}.

In particular, we~can extract a subsequence (still denoted $u_{n}$) such that $u_{n}$ converges weakly to some $u\in L^{4}([0, T],L^{12}(\M))$.

We will prove that $u$ is solution of the nonlinear equation with initial datum $(u_{0},u_{1})$. More precisely, denoting $U_0=(u_{0},u_{1})$, the well-posedness theory allows to define the nonlinear solution
\begin{align*}
V(t)=e^{tA}U_0+\int_0^t e^{(t-s)A}F(V(s))ds=V_{\lin}+V_{\Nlin},
\end{align*}
and we want to prove that $U=V$ where $U=(u,\partial_{t }u)$. The operators $A$ and $F$ are defined as in \cref{s:not_wave} with the appropriate choice of $\beta$.

Observe that the sequence $(u_{n}, \partial_{t }u_{n})$ is bounded in $C^{0}([0, T],H_0^1(\M)\times L^2(\M))$. Using the Aubin-Lions lemma (see, for instance, \cite[Cor.\,4]{S:87}), we~obtain that, for any $\eta>0$, still up to a subsequence, we~can assume that $u_{n}$ converges strongly to $u$ in $C^{0}([0, T],H^{1-\eta}(\M))$. In~particular, choosing $\eta=\psfrac{5-p}{4}>0$, by~Sobolev embedding, we~conclude that $u_{n}$ converges strongly to $u$ in $L^{\infty}([0, T],L^{r}(\M))$ for $r=\spfrac{12}{7-p}<6$. Using \eqref{hip:nonlinearity-hyp-1}, Hölder estimates for $\sfrac{1}{2}=\sfrac{1}{r}+\psfrac{p-1}{12} $ and $p<5$, we~get
\begin{align*}
\|f(u_{n}&)-f(u)\|_{L^{1}([0, T], L^{2}(\M))}\\
&\leq \norm{u_{n}-u}_{L^{\infty}([0, T], L^{r}(\M))}\norm{1+ u_{n}^{p-1}+u^{p-1}}_{L^{1}([0, T],L^{\frac{12}{p-1}}(\M))}\\
&\leq C \norm{u_{n}-u}_{L^{\infty}([0,T], L^{r}(\M))}\Bigl(1+\norm{u_{n}}_{L^{p-1}([0, T],L^{12}(\M))}^{p-1}+\norm{u}_{L^{p-1}([0, T],L^{12}(\M))}^{p-1}\Bigr)\\
&\leq C(T) \norm{u_{n}-u}_{L^{\infty}([0,T], L^{r}(\M))}\Bigl(1+\norm{u_{n}}_{L^{4}([0, T],L^{12}(\M))}^{p-1}+\norm{u}_{L^{4}([0, T],L^{12}(\M))}^{p-1}\Bigr).
\end{align*}
We obtain that $f(u_{n})$ converges strongly to $f(u)$ in $L^{1}([0, T], L^{2}(\M))$. The Duhamel formulation gives
\begin{align*}
U^n(t)=e^{tA}U_0^n+\int_0^t e^{(t-s)A}F(U^n(s))ds=U_{\lin}^n+U_{\Nlin}^n,
\end{align*}
with
\begin{align}
\label{e:cvgNL}
\norm{U_{\Nlin}^n-V_{\Nlin}}_{\mc{Z}}\underset{n\to +\infty}{\to}0.
\end{align}
Since the map $U_0 \mto e^{\cdot A}U_0$ is linear continuous from $X$ to $\mc{Z}$, it~is also continuous for the weak topology on each space. In~particular, since $U_0^n$ converges weakly to $U_0$ in $X$, we~obtain that $e^{\cdot A}U_0^n$ converges weakly to $e^{\cdot A}U_0=V_{\lin}$ in $L^{4}([0, T],L^{12}(\M))$. In~particular, $U^n$ converges weakly to $V$ in $\mc{Z}$ and $U=V$ as expected.

We have obtained that $u\in C^{0}([0, T],H_0^1(\M))\cap L^{4}([0, T],L^{12}(\M))$ is a mild solution~of
\begin{align*}
\left\{\begin{aligned}
\partial_t^2 u-\Delta_g u+f(u)&=0 &&(t, x)\in [0, T]\times \text{Int}(\M),\\
u_{|_{\partial\M}}&=0 &&(t, x)\in [0, T]\times\partial\M,\\
(u, \partial_t u)(0)&=(u_0, u_1)&&x\in \M .
\end{aligned}\right.
\end{align*}
Moreover, using \eqref{obsevNLintronegation} and taking weak limit, we~get $\partial_t u=0$ in $[0, T]\times\omega$.
In particular, we~are in a position to apply \cref{thm:unique-continuation-nlw} and we obtain $u=0$. In~particular, \eqref{e:cvgNL} can be written
\begin{align}
\label{e:cvgNLbis}
\norm{U_{\Nlin}^n}_{\mc{Z}}\underset{n\to +\infty}{\to}0.
\end{align}
The observability inequality allows to write
\begin{align*}
\norm{U_0^n}_{X}^2 &\leq \mathfrak{C}_{\textup{obs}}^2\int_0^T \norm{\mathbbm{1}_\omega\partial_t u_{\lin}^n}_{L^2(\M)}^2dt\\
&\leq 2\mathfrak{C}_{\textup{obs}}^2\int_0^T \norm{\mathbbm{1}_\omega\partial_t u^n}_{L^2(\M)}^2dt+2\mathfrak{C}_{\textup{obs}}^2\int_0^T \norm{\mathbbm{1}_\omega\partial_t u_{\Nlin}^n}_{L^2(\M)}^2dt.
\end{align*}
When combined with \eqref{obsevNLintronegation} and \eqref{e:cvgNLbis}, we~obtain $\alpha_n:=\norm{U_0^n}_{X}\underset{n\to +\infty}{\to}0$.

Now that we know that the initial datum converges to zero strongly, we~can ``linearize'' and consider the nonlinear solution as close to the linear one for which the observability is known. More precisely, denote $w_n=u_{n}/\alpha_n$ mild solution of
\begin{align*}
\left\{\begin{aligned}
\partial_t^2 w_n-\Delta_g w_n+\alpha_n^{-1}f(\alpha_n w_n)&=0 &&(t, x)\in [0, T]\times \text{Int}(\M),\\
{w_n}_{|_{\partial\M}}&=0 &&(t, x)\in [0, T]\times\partial\M,\\
(w_n, \partial_t w_n)(0)&=(w_{n,0}, w_{n,1})&&x\in \M,
\end{aligned}\right.
\end{align*}
with $\norm{(w_{n,0}, w_{n,1})}_{X}=1$. If we write $f(s)=f'(0)s+h(s)$ (with $f'(0)\geq 0$ thanks to the assumption on $f$), we~have the estimate
\begin{align*}
\ |h(s)|\leq C(|s|^2+|s|^p)\ \text{and}\ |h'(s)|\leq C(|s|+|s|^{p-1}).
\end{align*}
The nonlinearity $f_n(s):=\alpha_n^{-1}f(\alpha_n s)$ can be written as $f_n(s)=f'(0)s+\alpha_n^{-1}h(\alpha_n s)$. In~particular, $h_n(s):=\alpha_n^{-1}h(\alpha_n s)$ satisfies, uniformly in $n \in \N$,
\begin{align*}
\ |h_n(s)|\leq C\alpha_n(|s|^2+|s|^p)\ \text{and}\ |h'_n(s)|\leq C\alpha_n(|s|+|s|^{p-1}).
\end{align*}
Now, denoting
\begin{align*}\arraycolsep4.5pt
\widetilde A=\begin{pmatrix}
0 & I \\
\Delta_g -f'(0)& 0
\end{pmatrix}\ \text{ and }
H_n\begin{pmatrix}
u \\
v
\end{pmatrix}=\begin{pmatrix}
0 \\
-h_n(u)
\end{pmatrix}\!,
\end{align*}
we have
\begin{align*}
W^n(t)=e^{t\widetilde A}W_0^n+\int_0^t e^{\widetilde A(t-s)}H_n(W^n(s))ds=W_{\lin}^n+W_{\Nlin}^n,
\end{align*}
with
\begin{align}
\norm{W_{\lin}^n}_{\mc{Z}}&\!\leq\! C\\
\label{estimWNlin}\norm{W_{\Nlin}^n}_{\mc{Z}}&\!\leq\! C\alpha_n \norm{W^n}_{C^0([0, T],X)}\big(\norm{w^n}_{L^{4}([0, T],L^{12}(\M))}\!+\!\norm{w^n}_{L^{4}([0, T],L^{12}(\M))}^{p-1}\big).
\end{align}
This gives
\begin{align*}
\norm{W^n}_{\mc{Z}}&\leq C+ C\alpha_n\bigl( \norm{W^n}_{\mc{Z}}^2+ \norm{W^n}_{\mc{Z}}^p\bigr).
\end{align*}
In particular, a bootstrap argument (see for instance \cite[Lem.\,2.2.]{BG:99} for a slightly different case) allows to prove, for $n$ large enough,
\begin{align*}
\norm{W^n}_{\mc{Z}}&\leq 2C,
\end{align*}
which, after getting back to \eqref{estimWNlin} gives
\begin{align*}
\norm{W_{\Nlin}^n}_{\mc{Z}}&\leq C\alpha_n .
\end{align*}
We are now in position to apply the observability estimate of \cref{prop:higher-reg-obs}, taking into account that the assumptions imply $f'(0)>0$ when $\partial \M=\emptyset$. We~then write it for $W^n_{\lin}=(w^n_{\lin},\partial_tw_{\lin}^n)$,\vspace*{-10pt}
\begin{multline*}
\norm{(w_{0,n},w_{1,n})}_{H_0^1(\M)\times L^2(\M)}^2\leq C \int_0^T \norm{\mathbbm{1}_\omega\partial_t w^{n}_{\lin}(t)}_{L^2(\M)}^2 dt\\
\leq C \int_0^T \norm{\mathbbm{1}_\omega\partial_t w^{n}_{\Nlin}(t)}_{L^2(\M)}^2 dt+ C \int_0^T \norm{\mathbbm{1}_\omega\partial_t w^{n}(t)}_{L^2(\M)}^2 dt.
\end{multline*}
Concerning the first term, we~use energy estimates and get\vspace*{-3pt}
\begin{align*}
\int_0^T \norm{\mathbbm{1}_\omega\partial_t w^{n}_{\Nlin}(t)}_{L^2(\M)}^2 dt\leq C\norm{W_{\Nlin}^n}_{\mc{Z}}^2\leq C\alpha_n^2.
\end{align*}
For the second term, estimate \eqref{obsevNLintronegation} with $\alpha_n=\norm{(u_{0,n},u_{1,n})}_{H_0^1(\M)\times L^2(\M)}$ and the scaling $w_n=u_n/\alpha_n$ can be written as\vspace*{-3pt}
\begin{align*}
\int_0^T \norm{\mathbbm{1}_\omega\partial_t w_{n}(t)}_{L^2(\M)}^2 dt
\leq \frac{1}{n}. \end{align*}
The combination of the previous estimates give\vspace*{-3pt}
\begin{align*}
\norm{(w_{0,n},w_{1,n})}_{H_0^1(\M)\times L^2(\M)}^2\leq C\alpha_n^2+\frac{C}{n}.
\end{align*}
Yet, $\norm{(w_{0,n},w_{1,n})}_{H_0^1(\M)\times L^2(\M)}=1$, which is a contradiction.
\end{proof}

\section{Applications to the plate equation}\label{s:plate}
In this section, we~will give another application of the abstract \cref{thmabstractanalyticintro} concerning the plate equation. The purpose will be to obtain \cref{thm:analytic-nlp}. We~begin by presenting the equation and notations.
\subsection{Semilinear plate equation} Let $T>0$. Let us consider $\M$ to be a compact connected Riemannian manifold with smooth boundary $\partial\M$ and the hinged semilinear plate equation\vspace*{-3pt}
\begin{align}\label{eq:nlp-1}
\left\{\begin{aligned}
\partial_t^2 u+\Delta_g^2 u+f(u)&=0 &&(t, x)\in [0, T]\times \M,\\
u_{|_{\partial\M}}=\Delta u_{|_{\partial\M}}&=0 &&(t, x)\in [0, T]\times\partial\M,\\
(u, \partial_t u)(0)&=(u_0, u_1) &&x\in \M,
\end{aligned}\right.
\end{align}
where $(u_0, u_1)\in \big(H^2(\M)\cap H_0^1(\M)\big)\cap L^2(\M)$ and $f: \R\to \R$ is assumed to be analytic and to satisfy $f(0)=0$.

Given the abstract result described in \cref{s:abstrIntro} and the strategy we have used to obtain unique continuation property for the wave equation, it~is then natural to ask if it also holds for system \eqref{eq:nlp-1} when the nonlinearity is assumed to be analytic.

\subsubsection{Notation}\label{notplate} Let $A_0=-\Delta_g$ be the Laplace-Beltrami operator, equipped with Dirichlet boundary conditions if $\partial\M\neq\emptyset$. Recall that $A_0: D(A_0)\to L^2(\M)$ is a self-adjoint and nonnegative operator. In~this case, its domain is given by $D(A_0)=H^2(\M)\cap H_0^1(\M)$ on $L^2(\M)$. Set $X=D(A_0)\times L^2(\M)$ and introduce the densely defined operator $\bA: D(\bA)\to X$ given by
\begin{align*}\arraycolsep4.5pt
\bA=\begin{pmatrix}
0 & I \\
-A_0^2 & 0
\end{pmatrix}\quad \text{with } D(\bA)=D(A_0^2)\times D(A_0).
\end{align*}
Since $A_0$ is self-adjoint, then $A_0^2$ is a strictly positive operator and so $\bA$ is skew-adjoint. As for the wave operator, a simple computation then shows that $\bA\bA^*=-\bA^2$ has compact resolvent, henceforth $\bA$ satisfies \cref{assumAA} on $X$.

For $\sigma\in [0, 1]$, in~this section, $X^\sigma$ denotes the space
\begin{align*}
X^\sigma=D(A_0^{1+\sigma})\times D(A_0^{\sigma})=H_D^{2+2\sigma}\times H_D^{2\sigma},
\end{align*}
where we recall the notation introduced in \eqref{defSob}. By~Stone's theorem, $\bA$ generates a unitary $C_0$-group on $X$ and $D(\bA)$. By~linear interpolation, so it does on $X^\sigma$ for any $\sigma\in [0, 1]$.

\begin{remark}
We have chosen to consider the \emph{hinged} boundary conditions for simplicity. Other boundary conditions could have been considered, but they would require a more careful analysis; see \cite{ET:15}.
\end{remark}

For some $\widetilde{\chi}\in C^{\infty}(\M)$ to be chosen later, we~set $F\in C^0(X)$ to be the map
\begin{align*}
F: (u, v)\in X\mto (0, -\widetilde{\chi}f(u))\in X.
\end{align*}
Following the exact same steps of \cref{prop:f-assumptions} and adapting the spaces to this setup, we~have an analogous result for $F$. We~only need to note that for $d\leq 3$, $H^{2}_D\subset L^{\infty}$.

\begin{proposition}\label{prop:f-assumptions-pl}
Set $\sigma=0$ and $d\leq 3$. If $f$ is real analytic and $f(0)=0$, then $F$ satisfies \cref{assumFholom} for some $\veps>0$.
\end{proposition}

\subsection{From Schrödinger's to plate's observability}
The objective of this section is to discuss abstractly Lebeau's strategy \cite{Leb92} for deriving an observability inequality for the plate equation from that of Schrödinger. This analysis is based on the observation that, under hinged boundary conditions, the bi-Laplace operator is precisely the square root of the Dirichlet-Laplace operator.

\subsubsection{Abstract framework for transferring observability estimates} Let $H_0$ and $Y$ be Hilbert spaces with norms $\norm{\cdot}_0$ and $\norm{\cdot}_Y$, with $H_0$ separable. Let $A$ be self-adjoint, positive and boundedly invertible unbounded operator on $H_0$ with domain $D(A)$. Furthermore, we~assume that $A$ has compact resolvents.

\begin{remark}
In the application, the operator $A$ will sometimes only be nonnegative because of the eigenvalue $0$. Yet, we~will easily treat this subspace and this assumption makes the proof easier to write.
\end{remark}

We introduce the Sobolev scale of spaces based on $A$. For any $s>0$, let $H_s$ denote the Hilbert space $D(A^{s/2})$ with the norm $\norm{x}_s=\norm{A^{s/2}x}_{0}$ (which is equivalent to the graph norm $\norm{x}_{0}+\norm{x}_s=\norm{A^{s/2}x}_{0}$ since $0\in \rho(A)$). We~identify $H_0$ with its dual with respect to its inner product (\ie we will use it as a pivot space). Let $H_{-s}$ denote the dual of $H_s$. Since $H_s$ is densely continuously embedded in $H_0$, the pivot space $H_0$ is densely continuously embedded in $H_{-s}$, and $H_{-s}$ is the completion of $H_0$ with respect to the norm $\norm{x}_{-s}=\norm{A^{-s/2}x}_0$. We~will still denote by $A$ the restriction of $A$ to $H_s$ with domain $H_{s+2}$. It~is self-adjoint with respect to the $H_s$ scalar product. By~the spectral theorem, there exists an orthonormal basis of $H_0$ consisting of eigenfunctions of $A$, denoted by $\{(e_k, \ld_k)\}_k$, where $Ae_k=\ld_k e_k$ for each $k$. In~such a case, we~can characterize the $H_s$-norm as follows,
\begin{align*}
\norm{x}_s^2=\sum_{k\in \N} \ld_k^{s}|x_k|^2,\ \forall s\in \R
\end{align*}
where $x=\sum_{k\in \N} x_ke_k$. We~will also make the technical assumption that there exists $N\in \N$ such that
\begin{align}
\label{hypWeyl} \sum_{k\in \N}\ld_k^{-N}<+\infty.
\end{align}

Let $\mc{C}\in \mc{L}(H_2, Y)$. Let us consider the first order system for $A$,
\begin{equation}
\label{appendix:first-order-system}
\dot{\psi}(t)-iA\psi(t)=0,\quad \psi(0)=\psi_0\in H_0,\quad y(t)=\mc{C}\psi(t).
\end{equation}
We assume that $\mc{C}$ is an admissible observation operator for $e^{i\cdot A}$, which means that, for some $\tau>0$ (and thus for any times by the group property), there exists $K_\tau\geq 0$ such that
\begin{align*}
\int_0^\tau \norm{\mc{C}e^{itA}\psi_0}_Y^2dt\leq K_\tau \norm{\psi_0}_{H_0}^2,\ \ \forall \psi_0\in H_2.
\end{align*}
\begin{remark}
Under the admissibility assumption, the output map $\psi_0\mto \mc{C}\psi_0$ from $H_2$ to $L_{\textup{loc}}^2(\R, Y)$ has a continuous extension to $H_0$.
\end{remark}
\begin{definition}
Let $-\infty< \tau_1<\tau_2<\infty$. We~say that system \eqref{appendix:first-order-system} is weakly observable by $\mc{C}$ in $(\tau_1, \tau_2)$ if there exists $C>0$ such that
\begin{align*}
\norm{\psi_0}_{H_0}^2\leq C\biggl(\int_{\tau_1}^{\tau_2} \norm{\mc{C}e^{itA}\psi_0}_{Y}^2dt+\norm{\psi_0}_{H_{-2}}^2\biggr),
\end{align*}
for any $\psi_0\in H_0$.
\end{definition}

Note that the conservation of $H_0$ and $H_{-2}$-norms and translation invariance in time show that it is equivalent to state it for $\tau_1=0$, the only relevant quantity being $\tau_2-\tau_1$. Moreover, changing $\psi(t)$ to $\psi(-t)$ allows to get the same result for $e^{-itA}$.

Consider the second-order observability system
\begin{align}\label{appendix:second-order-system}
\left\{\begin{array}{rl}
\ddot{z}(t)+A^2z(t)&=0, \\[3pt]
z(0)&=z_0\in H_2,\\[3pt]
\dot{z}(0)&=z_1\in H_0,
\end{array}\right.
\end{align}
with the observation being
\begin{align*}
\text{either }\ y(t)=\mc{C}Az(t)\ \text{ or }\ y(t)=\mc{C}\dot{z}(t).
\end{align*}
System \eqref{appendix:second-order-system} can be recast as a first order system, with $x(t)=(z(t), \dot{z}(t))\in X=H_2\times H_0$,
\begin{align*}
\A(z_0, z_1):=(z_1, -A^2z_0)\ \text{ with }\ D(\A)=H_4\times H_2,
\end{align*}
considering as possible observations
\begin{equation}
\begin{aligned}
\label{defobservplate}
\bC(z_0, z_1)&=\mc{C}Az_0\\
\bC_1(z_0, z_1)&=\mc{C}z_1.
\end{aligned}
\end{equation}
Note that $\bC\in \mc{L}(H_4\times H_2, Y)$ is admissible for $e^{\cdot \A}$and the same holds for $\bC_1$.

\begin{proposition}\label{prop:w-plates-schr}
Let $\widetilde{T},T$ be such that $0 < \widetilde{T} < T < \infty$. If the Schrödinger-like equation
\eqref{appendix:first-order-system} is weakly observable by $\mc{C}\in \mc{L}(H_2, Y)$ on $(0, \Tilde{T})$, then the plate-like equation \eqref{appendix:second-order-system} is weakly observable by $\bC$ or $\bC_1$ on $(0, T)$ defined in \eqref{defobservplate}. That means, there exists $C_1>0$ such that
\begin{align}\label{ineq:wobs-plates}
\norm{(z_0, z_1)}_{H_2\times H_0}^2\leq C_1\biggl(\int_0^T \norm{\bC e^{t\A}(z_0, z_1)}_{Y}^2dt+\norm{(z_0, z_1)}_{H_{0}\times H_{-2}}^2\biggr),\\
\norm{(z_0, z_1)}_{H_2\times H_0}^2\leq C_1\biggl(\int_0^T \norm{\bC_1 e^{t\A}(z_0, z_1)}_{Y}^2dt+\norm{(z_0, z_1)}_{H_{0}\times H_{-2}}^2\biggr),
\end{align}
for all $(z_0, z_1)\in H_2\times H_0$.
\end{proposition}
\begin{proof}
Based on the factorization $\partial_t^2 z+A^2z=(\partial_t+iA)(\partial_t-iA)z$, we~introduce the splitting
\begin{align*}
z_+=\dfrac{1}{2}(z_0-iA^{-1}z_1),\quad z_-=\dfrac{1}{2}(z_0+iA^{-1}z_1),
\end{align*}
so that
\begin{align*}
z_0=z_++z_-,\quad z_1=iA(z_+-z_-).
\end{align*}
Consequently, $(\partial_t\mp iA)e^{\pm itA}z_{\pm}=0$ implies $(\partial_t^2+A^2)e^{\pm itA}z_{\pm}=0$. Hence, $(z(t), \dot{z}(t))=e^{t\A}(z_0, z_1)$, solution of \eqref{appendix:second-order-system}, can be rewritten as
\begin{align*}
z(t)=e^{itA}z_++e^{-itA}z_-.
\end{align*}
For any $s\in \R$, denote by $\Lambda: H_s\times H_{s-2}\to H_s\times H_s$ the isomorphism corresponding to the previous splitting $\Lambda(z_0, z_1)=(z_+, z_-)$. Observe that it is almost an isometry
\begin{align*}
\norm{(z_0, z_1)}_{H_s\times H_{s-2}}^2&=\norm{z_++z_-}_{H_s}^2+\norm{A(z_+-z_-)}_{H_{s-2}}^2\\&=2\big(\norm{z_+}_{H_s}^2+\norm{z_-}_{H_s}^2\big).
\end{align*}
By considering the above splitting, we~can write for the $H_s\times H_{s-2}$-energy of the system \eqref{appendix:second-order-system} as follows
\begin{align*}
\E_s(z_0, z_1)=\dfrac{1}{2}\norm{(z_0, z_1)}_{H_{s}\times H_{s-2}}^2=\E_s(\Ld^{-1}(z_+, z_-)).
\end{align*}

Let $T',T'' \in (0,T)$ and $\rho \in C^\infty_c(\mathbb{R},\mathbb{R}^+)$
be such that $T''-T'>\widetilde{T}$, $\rho \equiv 1$ on $(T',T'')$ and $\text{supp}(\rho) \subset (0,T)$. We~have the identity
\begin{multline*}
\int_0^T \norm{\rho(t)\bC e^{t\A}(z_0, z_1)}_{Y}^2=\int_0^T \norm{\rho(t)\mc{C}Ae^{itA}z_+}_{Y}^2dt\\+\int_0^T \norm{\rho(t)\mc{C}Ae^{-itA}z_-}_{Y}^2dt+2\Re\int_0^T \rho^2(t)\inn{\mc{C}Ae^{itA}z_+, \mc{C}Ae^{-itA}z_-}_{Y}dt.
\end{multline*}
while
\begin{multline*}
\int_0^T \norm{\rho(t)\bC_1 e^{t\A}(z_0, z_1)}_{Y}^2=\int_0^T \norm{\rho(t)\mc{C}Ae^{itA}z_+}_{Y}^2dt\\+\int_0^T \norm{\rho(t)\mc{C}Ae^{-itA}z_-}_{Y}^2dt-2\Re\int_0^T \rho^2(t)\inn{\mc{C}Ae^{itA}z_+, \mc{C}Ae^{-itA}z_-}_{Y}dt.
\end{multline*}
On the one hand, since $A$ and $e^{\pm i\cdot A}$ commute, the weak observability inequality for \eqref{appendix:first-order-system} implies
\begin{align*}
\int_0^T \norm{\rho(t)\mc{C}Ae^{\pm itA}z_{\pm}}_{Y}^2dt\geq \dfrac{1}{C}\big(\norm{Az_\pm}_{H_0}^2-\norm{Az_\pm}_{H_{-2}}^2\big).
\end{align*}
Putting the two inequalities together, we~get
\begin{multline*}
\int_0^T \norm{\rho(t)\mc{C}Ae^{itA}z_+}_{Y}^2dt+\int_0^T \norm{\rho(t)\mc{C}Ae^{-itA}z_-}_{Y}^2dt\\
\geq \dfrac{1}{2C}\big(\E_2(\Lambda^{-1}(z_0, z_1))-\E_{0}(\Lambda^{-1}(z_0, z_1))=\dfrac{1}{2C}\big(\E_2(z_0, z_1)-\E_{0}(z_0, z_1)\big).
\end{multline*}
On the other hand, to treat the interaction term, we~first look at its spectral expansion. By~denoting $z_{\pm, k}=\inn{z_{\pm}, e_k}_{H_0}$ the Fourier coefficients of $z_{\pm}$, we~have
\begin{multline*}
\int_0^T \rho^2(t)\inn{\mc{C}Ae^{itA}z_+, \mc{C}Ae^{-itA}z_-}_{Y}dt\\[-5pt]
=\sum_{j, k} z_{+, j}z_{-, k} \inn{\mc{C}Ae_j, \mc{C}Ae_k}_{Y} \biggl(\int_\R e^{it(\ld_j+\ld_k)}\rho^2(t)dt\biggr).
\end{multline*}
First of all, we~use that $\mc{C}$ is a bounded operator from $H_2$ to $Y$ to estimate $\inn{\mc{C}e_j, \mc{C}e_k}_{Y}$ as follows:
\begin{align*}
|\inn{\mc{C}e_j, \mc{C}e_k}_{Y}|\leq \norm{\mc{C}e_k}_{Y}\norm{\mc{C}e_j}_{Y}\leq C\norm{Ae_k}_{H_0}\norm{Ae_j}_{H_0}=C\ld_k\ld_j.
\end{align*}
Since $\rho^2\in C_c^\infty(\R)$ is smooth compactly supported function, for any $N\in \N$, we~can find $C_N>0$ such that $|\widehat{\rho^2}(\xi)|\leq C_N \left<\xi\right>^{-N}$ and
\begin{align*}
\left|\int_\R e^{it(\ld_j+\ld_k)}\rho^2(t)dt\right|\leq \dfrac{C_N}{(\ld_j+\ld_k)^{N}} \leq\dfrac{C_N}{(\ld_j\ld_k)^{N/2}},\quad \forall j, k\in \N_0.
\end{align*}
We then have
\begin{align*}
\biggl|\int_0^T \rho^2(t)\inn{A\mc{C}e^{itA}z_+, A\mc{C}&ee^{-itA}z_-}_{Y}dt\biggr|\\
&\leq \sum_{j, k\in \N_0} \ld_j\ld_k |z_{+, j}||z_{-, k}|\frac{C_N}{(\ld_j\ld_k)^{N/2}}\\
&\leq C_N\sum_{j, k\in \N_0} \frac{1}{(\ld_j\ld_k)^{N/2-1}}\left(|z_{+, j}|^2+|z_{-, k}|^2\right)\\
&\leq C_N\biggl(\sum_{j\in \N_0}\frac{1}{\ld_j^{N/2-1}}\biggr)
\biggl(\sum_{k\in \N_0} \frac{1}{\ld_k^{N/2-1}}\big(|z_{+, k}|^2+|z_{-, k}|^2\big)\biggr),
\end{align*}
where we have chosen $N$ large enough so that $\ld_J^{1-N/2}$ is summable by \eqref{hypWeyl}. By~adjusting $N$ if necessary, we~conclude the proof by observing that
\begin{align*}
\sum_{k\in \N_0} \dfrac{1}{\ld_k^{N/2-1}} \big(|z_{+, k}|^2+|z_{-, k}|^2\big)\leq \norm{(z_+, z_-)}_{H_{-2}, H_{-2}}^2=2\norm{(z_0, z_1)}_{H_{0}, H_{-2}}^2.\qedhere
\end{align*}
\end{proof}

Let us make the following unique continuation assumption on the pair $(A, \mc{C})$.

\begin{assump}{UCP}\label{assumAeigen}
For any eigenvector $\psi$ of $A$ such that $\mc{C}\psi=0$, then $\psi\equiv 0$.
\end{assump}

\begin{proposition}\label{prop:nt-plates}
Let $\widetilde{T},T$ be such that $0 < \widetilde{T} < T < \infty$. Assume that the pair $(A, \mc{C})$ satisfies \cref{assumAeigen}. If the Schrödinger-like equation
\eqref{appendix:first-order-system} is weakly observable by $\mc{C}$ on $(0, \Tilde{T})$, then for any non zero $(z_0, z_1)\in H_2\times H_0$,
\begin{align*}
\bC e^{t\A}(z_0, z_1)\neq 0\ \text{ in }\ L^2([0, T], Y),\\
\bC_1 e^{t\A}(z_0, z_1)\neq 0\ \text{ in }\ L^2([0, T], Y).
\end{align*}
\end{proposition}
\begin{proof}
We consider first the observation by $\bC$. Let us consider
\begin{align*}
N_T:=\{(z_0, z_1)\in H_2\times H_0\ |\ \bC e^{t\A}(z_0, z_1)=0,\ \text{ in } L^2([0,T], Y)\}.
\end{align*}
Since the equation is linear, it~is a linear subspace of $H_2\times H_0$. \cref{prop:w-plates-schr} implies that for all $(z_0, z_1)\in N_T$,
\begin{align}\label{prop:ineq:plates_nt}
\norm{(z_0, z_1)}_{H_2\times H_0}\leq C\norm{(z_0, z_1)}_{H_0\times H_{-2}}.
\end{align}

We will argue by contradiction to prove that $N_T=\{0\}$. Let $0<\veps<T-\widetilde{T}$ and $(z_0, z_1)\in N_T$. Let us introduce the sequence
\begin{align*}
(z_0^\veps, z_1^\veps)=\dfrac{1}{\veps}\big(e^{\veps\A}(z_0, z_1)-(z_0, z_1)\big),
\end{align*}
and note that it belongs to $N_{\widetilde{T}}$. Observe that $(z_0, z_1)\in H_2\times H_0$ implies that $(A^{-1}z_0, A^{-1}z_1)\in H_4\times H_2=D(\A)$. By~classical semigroup theory, we~have
\[
(A^{-1}z_0^\veps, A^{-1}z_1^\veps)=\dfrac{1}{\veps}\big(e^{\veps\A}(A^{-1}z_0^\veps, A^{-1}z_1^\veps)-(A^{-1}z_0^\veps, A^{-1}z_1^\veps)\big)
\xrightarrow[\veps\to 0]{} \A (A^{-1}z_0, A^{-1}z_1),
\]
where the convergence holds in $H_2\times H_0$. It~follows that $\big((z_0^\veps, z_1^\veps)\big)_{\veps}$ is a Cauchy sequence for the norm $\norm{\cdot}_{H_0\times H_{-2}}$ and so it is for the norm $\norm{\cdot}_{H_2\times H_{0}}$, due to inequality~\eqref{prop:ineq:plates_nt}. Therefore, the limit $\A(z_0, z_1)$ belongs to $H_2\times H_0$, namely, $(z_0, z_1)\in D(\A)$ and so~$N_T\subset\nobreak D(\A)$. The regularity in time of $e^{t\A}$ allows us to take $\partial_t$ in the observation condition, obtaining
\begin{align*}
\partial_t\bC e^{t\A}(z_0, z_1)=\bC \partial_t e^{t\A}(z_0, z_1)=\bC e^{t\A}\A(z_0, z_1).
\end{align*}
This means that $N_T$ is stable under $\A$ and it only contains elements of $D(A^k)\times D(A^k)$ for any $k\in \N$.

Applying inequality \eqref{prop:ineq:plates_nt} to $\A(z_0, z_1)$ yields $\norm{(z_0, z_1)}_{H_4\times H_2}\leq C\norm{(z_0, z_1)}_{H_2\times H_0}$. We~deduce that the unit ball of $N_T$ in the $H_2\times H_0$ topology, is bounded in $H_4\times H_2$ and thus it is compact by compact embedding (recall, $A$ has compact resolvent). By~Riesz's theorem, $N_T$ is a finite-dimensional subspace of $H_2\times H_0$. Since $-\A^2$ is self-adjoint positive and sends $N_T$ into itself, it~admits an eigenvalue $\ld\geq 0$ associated to the eigenvector $(v_\ld, w_\ld)$. We~arrive at the system
\begin{align*}
\left\{\begin{array}{cc}
A^2 v_\ld=\ld v_\ld, & \\[3pt]
A^2w_\ld=\ld w_\ld, &
\end{array}\right.
\end{align*}
Since $A^2$ is a strictly positive operator, we~can write $\lambda=\alpha^2$ for $\alpha> 0$ and get $(A+\alpha I)(A-\alpha I)v_{\lambda}=0$. Using that $A+\alpha I$ is a strictly positive operator, we~get $Av_{\lambda}=\alpha v_{\lambda}$ and $Aw_{\lambda}=\alpha w_{\lambda}$. In~particular,
\begin{align*}
e^{t\A}(v_\ld, w_\ld)=(\cos(\alpha t)v_\ld+\alpha^{-1}\sin (\alpha t) w_\ld,-\alpha\sin(\alpha t)v_\ld+\cos (\alpha t) w_\ld).
\end{align*}
Therefore, such eigenvector must satisfy $\bC e^{t\A}(v_\ld, w_\ld)=0$ for $t\in [0,T]$ and hence $\mc{C}A v_\ld=\mc{C}A w_\ld=0$. Using \cref{assumAeigen} and $\alpha\neq 0$, we~deduce that $(v_\ld,w_{\ld})=(0,0)$. This is a contradiction to the fact that $N_T\neq \{0\}$.

For $\bC_1$, the proof is the same except that it leads to the unique continuation problem $-\alpha\sin(\alpha t)\mc{C}v_\ld+\alpha \cos (\alpha t) \mc{C}w_\ld=0$, $t\in [0,T]$, for which \cref{assumAeigen} still applies.
\end{proof}

\begin{proposition}\label{prop:obs-schr-plates}
Let $\widetilde{T},T$ be such that $0 < \widetilde{T} < T < \infty$. If the Schrödinger-like equation
\eqref{appendix:first-order-system} is weakly observable on $(0, \Tilde{T})$ and \cref{assumAeigen} is satisfied, then there exists $\mathfrak{C}_{obs}>0$ such that
\begin{align}\label{ineq:obs-plates}
\norm{(z_0, z_1)}_{H_2\times H_0}^2&\leq \mathfrak{C}_{obs}^2\int_0^T \norm{\bC e^{t\A}(z_0, z_1)}_{Y}^2dt,\\
\label{ineq:obs-platesC1} \norm{(z_0, z_1)}_{H_2\times H_0}^2&\leq \mathfrak{C}_{obs}^2\int_0^T \norm{\bC_1 e^{t\A}(z_0, z_1)}_{Y}^2dt,
\end{align}
for all $(z_0, z_1)\in H_2\times H_0$.
\end{proposition}
\begin{proof}
The result follows from \cref{prop:w-plates-schr} and \cref{prop:nt-plates}. We~write it for $\bC$, but it is exactly the same for $\bC_1$.

Assume that the inequality \eqref{ineq:obs-plates} does not hold. Then we can find a sequence $(z_0^n, z_1^n)$ of norm $1$ in $H_2\times H_0$ such that
\begin{align}\label{plates-obs-conv}
\int_0^T \norm{\bC e^{t\A}(z_0^n, z_1^n)}_{Y}^2 dt\xrightarrow[n\to\infty]{} 0.
\end{align}
Let $(z_0, z_1)$ be the weak limit of $(z_0^n, z_1^n)$. The above inequality implies that $(z_0, z_1)\in N_T$ and so $(z_0, z_1)=(0, 0)$ by \cref{prop:nt-plates}. We~then have that $(z_0^n, z_1^n)\wto (0, 0)$ weakly in $H_2\times H_0$ and by compact embedding (recall $A$ has compact resolvent), we~get $\norm{(z_0^n, z_1^n)}_{H_0\times H_{-2}}\to 0$ as $n\to \infty$. Applying the weak observability estimate~\eqref{ineq:wobs-plates} along with \eqref{plates-obs-conv} we get that $\norm{(z_0^n, z_1^n)}_{H_2\times H_0}\to 0$ as $n\to \infty$, which is a contradiction.
\end{proof}

\begin{remark}
Lebeau's strategy gives a loss on time, harmless for our ends. Following Miller \cite{2012:miller:resolvent-control}, it~is possible to avoid such loss in time but it would imply to inherit the geometric setup of the wave equation, which is known to be more restrictive than the one of Schrödinger.
\end{remark}

\subsection{Proof of \cref{thm:analytic-nlp}}
We follow the notations of \cref{notplate}. We~will first need the following observability estimates in order to apply the abstract results.
\begin{proposition}\label{prop:obs-plate}
Let $\Tilde{T}>0$ and $\omega$ be an open subset of $\M$ such that the Schrödinger equation
\eqref{appendix:first-order-system} is weakly observable by $\mc{C}=\mathbbm{1}_{\omega}\in \mc{L}( L^2(\M))$ on $(0, \Tilde{T})$. Then for any $T$ such that $0<\Tilde{T}<T<+\infty$ and any $b_{\omega}\in C^{\infty}(\M)$ such that $b_{\omega}=1$ on $\omega$, there exists $C>0$ such that for any $s\in [0,2]$, any $(z_0, z_1)\in H_D^{2+s}\times H^s_D(\M)$ and associated solution $z$ of
\begin{align*}
\left\{\begin{aligned}
\partial_t^2 z+\Delta_g^2 z&=0 &&(t, x)\in [0, T]\times \M,\\
z_{|_{\partial\M}}=\Delta_g z_{|_{\partial\M}}&=0 &&(t, x)\in [0, T]\times\partial\M,\\
(z, \partial_t z)(0)&=(z_0, z_1) &&x\in \M,
\end{aligned}\right.
\end{align*}
we have, if~$\partial M\neq \emptyset$
\begin{align}\label{thm:ineq:obs-plate-1}
C\norm{(z_0, z_1)}_{H^{2+s}_D\times H^s_D(\M)}^2\leq \int_0^T \norm{b_\omega \Delta_g z(t)}_{H^s_D(\M)}^2 dt,
\end{align}
and, if $\partial M= \emptyset$,
\begin{align}\label{thm:ineq:obs-plate-1bis}
C\norm{(z_0, z_1)}_{H^{2+s}\times H^s(\M)}^2\leq \int_0^T \norm{b_\omega z(t)}_{H^{2+s}(\M)}^2 dt.
\end{align}
\end{proposition}
\begin{proof}
In the case $\partial M= \emptyset$, the operator $-\Delta_g$ is not strictly positive because of constants. So, we~first prove \eqref{thm:ineq:obs-plate-1} in the case $\partial M\neq \emptyset$ or $\partial M= \emptyset$ but $(z_0, z_1)$ are orthogonal to the set of constants. By~interpolation, it~is enough to consider the cases $s=0$ and $s=2$.

First of all, the condition on the eigenvalues \eqref{hypWeyl} is satisfied, for instance, due to Weyl's law. The unique continuation of eigenfunctions for the laplacian $\Delta_g$ is known to hold in our framework, see \cite[Prop.\,5.2.]{LLR:book1}, hence \cref{assumAeigen} is satisfied.

For $s=0$, this is a consequence of the \cref{prop:obs-schr-plates} for the case \eqref{ineq:obs-plates} and the fact that $b_{\omega}=1$ on $\omega$. For the second one, assume that $(z_0, z_1)\in H^4_D\times H^2_D$. Take~$w=\nobreak\partial_t z$ and observe that it is a mild solution of
\begin{align*}
\left\{\begin{aligned}
\partial_t^2 w+\Delta_g^2 w&=0, &\\
w_{|_{\partial\M}}=\Delta w_{|_{\partial\M}}&=0 &&\text{if } \partial\M\neq\emptyset,\\
(w, \partial_t w)(0)&=(z_1, \Delta z_0),
\end{aligned}\right.
\end{align*}
with $(z_1, \Delta_g z_0)\in H_D^2(\M)\times L^2(\M)$. Applying the observability inequality \eqref{ineq:obs-platesC1} to $w$ and then going back to the $z$ variable, we~get
\begin{align*}
\norm{(z_1, z_0)}_{H_D^4\times H_D^2}^2\leq C_2\int_0^T \norm{b_\omega \Delta_g^2 z(t)}_{L^2(\M)}^2dt.
\end{align*}
Note that
\begin{align*}
\int_0^T \norm{b_\omega \Delta_g^2 z(t)}_{L^2(\M)}^2dt\leq \int_0^T \norm{\Delta_g(b_\omega\Delta_g z)}_{L^2(\M)}^2dt+\int_0^T \norm{[b_\omega, \Delta_g]\Delta_g z}_{L^2(\M)}^2dt,
\end{align*}
so we now need to estimate the commutator term appearing on the right-hand side of the inequality above. Recall $[b_\omega, \Delta_g]\Delta_g z=2\nabla b_\omega\cdot\nabla(\Delta_g z)+\Delta_g b_\omega\Delta_g z$. In~the following estimates, the constant $C>0$ may change from line to line,
\begin{align*}
\int_0^T \norm{[b_\omega, \Delta_g]\Delta_g z}_{L^2(\M)}^2&dt\\[-8pt]
&\leq C\biggl(\int_0^T \norm{\nabla b_\omega\cdot\nabla(\Delta_g z)}_{L^2(\M)}^2dt
+\int_0^T \norm{\Delta_g b_\omega\Delta_g z}_{L^2(\M)}^2dt\biggr)\\
&\leq C\int_0^T \norm{z(t)}_{H_D^3}^2dt \leq C\norm{(z_0, z_1)}_{H_D^3\times H_D^1}^2\\
&\leq C\veps \norm{(z_0, z_1)}_{H_D^4\times H_D^2}^2
+\frac{C}{\veps}\norm{(z_0, z_1)}_{H_D^2\times L^2(\M)}^2\\
&\leq C\veps \norm{(z_0, z_1)}_{H_D^4\times H_D^2}^2
+\frac{C}{\veps}\int_0^T \norm{b_\omega \Delta_g z(t)}_{L^2(\M)}^2 dt,
\end{align*}
for $\veps>0$ to be chosen. Observe that we have used energy estimates, an interpolation inequality and the observability inequality \eqref{thm:ineq:obs-plate-1} for $s=0$. By~choosing $\veps>0$ small enough, the observability inequality \eqref{thm:ineq:obs-plate-1} for $s=2$ follows once we put all the inequalities above together.

It only remains to prove \eqref{thm:ineq:obs-plate-1bis} when there is no boundary. Decomposing $(z_0,z_1)=\pi_0(z_0,z_1)+\pi_0^{\perp}(z_0,z_1)$ where $\pi_0$ is the projection on the eigenvalues $0$ of $\Delta_g$, that is the constants, we~have obtained up to now, noticing the constant part of the initial data produce some part of the solution with zero Laplacian,
\begin{align*}
\norm{\pi_0^{\perp}(z_0,z_1)}_{H^{2+s}\times H^s(\M)}^2\leq C\int_0^T \norm{b_\omega \Delta_g z(t)}_{H^s(\M)}^2 dt.
\end{align*}
By adding the components corresponding to eigenvalue zero and noting that $[b_{\omega},\Delta_g]$ is a differential operator of order one, we~obtain
\begin{align*}
\norm{(z_0,z_1)}_{H^{2+s}\times H^s(\M)}^2\leq C\int_0^T \norm{ b_\omega z(t)}_{H^{2+s}(\M)}^2 dt+C\norm{(z_0,z_1)}_{H^{1+s}\times H^{s-1}(\M)}^2.
\end{align*}
A compactness-uniqueness argument as in \cref{prop:obs-schr-plates} allows to conclude.
\end{proof}
\subsubsection{Observability for the Schrödinger equation}\label{sssec:obs-schr} Let $\M$ be a compact Riemannian manifold with or without boundary equipped with a metric $g$ and take $\omega\subset \M$. In~what follows, we~consider an observation operator $\mathcal{C}\psi=\mathbbm{1}_\omega \psi$, unless we specify otherwise.

In any of the situations described in \cref{s:plateintro}, an observability inequality at the~$ L^2$ level holds for the linear Schrödinger equation. We~summarize this discussion in the following result.

\begin{theorem}\label{thm:schr-obs}
If we are in any of the examples described in \cref{s:plateintro}, then for every $T>0$, there exists $C>0$ such that
\begin{align*}
\norm{v_0}_{L^2(\M)}^2\leq C\int_0^T \norm{\mc{C} e^{it\Delta_g}v_0}_{L^2(\M)}^2dt,
\end{align*}
for all $v_0\in L^2(\M)$.
\end{theorem}
Note that $e^{it\Delta_g}$ is the flow with Dirichlet boundary condition in case that $\partial \M\neq \emptyset$. We~now verify that the pair $(\bA, \bC)$ satisfies \cref{assumcommu}.
\begin{proposition}\label{propcommutplate}
Let $\sigma\in [0, 1]$, $\sigma \neq 1/4$. If~$\bC$ is given by $ \bC(\phi, \psi)=(0, b_\omega\Delta_g \phi)$ with $b_{\omega}$ smooth satisfying $\partial_{\vec n}b_{\omega}=0$ on $\partial \M$, then \cref{assumcommu} is fulfilled with $s=1/2$ as long as $\veps\leq 1/2$, $\sigma+\veps<5/4$ and $\sigma+\veps\neq 1/4$. When $\partial M=\emptyset$, the same result holds with $ \bC(\phi, \psi)=(b_\omega\phi,0)$.
\end{proposition}
\begin{proof}
We compute $[(A^*A)^{1/2},\bC]=\begin{psmallmatrix}
0 & 0\\
[b_{\omega},\Delta_g] \Delta_g & 0
\end{psmallmatrix}$. Therefore, the result is true as long as $[b_{\omega},\Delta_g]\Delta_g=-2\nabla_g b_{\omega}\cdot \nabla_g \Delta_g-(\Delta_g b_{\omega})\Delta_g$ sends $H^{4+2\sigma}_D$ into $H^{2(\sigma+\veps)}_D$, that is, as long as $[b_{\omega},\Delta_g]=-2\nabla_g b_{\omega}\cdot \nabla_g -\Delta_g b_{\omega}$ sends $H^{2+2\sigma}_D$ into $H^{2(\sigma+\veps)}_D$. The assumption $\veps\leq 1/2$ ensures that the loss of derivative is correct, whereas $\sigma+\veps<5/4$ ensures that $H^{2(\sigma+\veps)}_D$ is either $H^{2(\sigma+\veps)}(\M)$ or $H^{2(\sigma+\veps)}_0(\M)$ with the Dirichlet boundary condition (provided we avoid the value $1/2$). Since $\partial_{\vec n}b_{\omega}=0$, this gives the result. This is similar in the other case.
\end{proof}

We now come to the proof of the main result of this section.

\begin{proof}[Proof of \cref{thm:analytic-nlp}]
It only remains to check that \cref{thmabstractanalyticintro} can be applied with the abstract notations in \cref{notplate}, with $\sigma=0$. The proof is very similar to the one of \cref{thm:analytic-prop}.

We have already established that $\bA$ satisfies \cref{assumAA}. Using \cref{lemma:omomtilde-smaller-subset} below, we~can construct successively $b_{\omega}$ and $\chi$ smooth, compactly supported in $\widetilde{\omega}$, with $\partial_{\vec n}\chi=\partial_{\vec n}b_{\omega}=0$ on $\partial M$ and such that $b_{\omega}=1$ on $\omega$ and $\chi=1$ on $\supp(b_{\omega})$. Take $\bC(\phi, \psi)=(0,b_\omega \Delta_g \phi)$ or $\bC(\phi, \psi)=(b_\omega \phi,0)$ if $\partial \M=\emptyset$ so that \cref{prop:obs-plate} applies. Therefore, we~get that that $\bA$ satisfies \cref{assumCC} with $\bC\in \mc{L}(X)$ and $\bC\in \mc{L}(X^\veps)$ for some $\veps>0$. Also, \cref{propcommutplate} applies so that \cref{assumcommu} is fulfilled.
By writing $u=\chi u+(1-\chi)u$, we~want to prove that $t\mto (1-\chi)u$ is analytic. Set $\Tilde{\chi}=(1-\chi)$ and consider the new variable $z=\Tilde{\chi} u$. We~then have
\begin{align*}
\partial_t^2 z+\Delta_g^2 z &=\Tilde{\chi}(\partial_t^2 u+\Delta_g^2 u)-[\Delta_g^2, \Tilde{\chi}]u=-\Tilde{\chi} f(u)+[\Delta_g^2, \chi]u\\
&=-\Tilde{\chi} f(z+\chi u)+[\Delta_g^2, \chi]u=-\Tilde{\chi} f(z+h_1)+h_2,
\end{align*}
with $h_1=\chi u$ and $h_2=[\Delta_g^2,\chi]u$. From \cref{prop:f-assumptions-pl}, we~see that $F$ satisfies \cref{assumFholom}. Since the multiplication by $\chi$ maps $H^2_D$ into itself, while $[\Delta_g^2,\chi]$ maps $H^{3+\veps}(\M)$ into $H^{\veps}_D$ for $\veps<1/2$. So, if~we choose $\sigma=0$ and $\veps<1/2$, the assumptions imply that $h_1$ and $h_2$ are analytic with value in $H^2_D$ and $H^{\veps}_D$, respectively. Now, the conclusion follows as a direct application of \cref{thmabstractanalyticintro} in the same way as \cref{thm:analytic-prop}. We~conclude that $t\mto (z(t), \partial_t z(t))$ is analytic with value in $X^0$. By~assumption, $t\mto \chi u(t)$ is analytic with value in $H^{3+\veps}\cap H_0^1$ (and so it is the same for $\chi \partial_t u$) and therefore with value in $H^2_D$. So, $t\mto (\chi u(t),\chi\partial_t u(t))$ is analytic with value in $X^0$.

Summing up, we~obtain that $t\mto ( u(t),\partial_t u(t))$ is analytic with value in $X^0$. Using the equation again as in \cref{thm:analytic-prop}, we~obtain that $t\mto \bA U(t)$ is analytic with value in $X^0$. Hence $t\mto U(t)$ is analytic with value in $H^4\cap H_0^1(\M)\times H^2\cap H_0^1(\M)$, which finishes the proof.
\end{proof}

\appendix
\setcounter{theorem}{0}
\renewcommand{\thetheorem}{A.\arabic{theorem}}
\renewcommand{\theequation}{A.\arabic{equation}}
\renewcommand{\thesubsection}{A.\arabic{subsection}}

\section*{Appendix}

\subsection{ODEs in Banach spaces} We now introduce the two different notions of ODEs in Banach spaces used in the present article. Let us consider the framework of \cref{sec:abstract-construction} and let $I\subset \R$ be a nonempty interval and take $s_0\in I$.

For any $s\in \R$, we~can easily extend $e^{sA}$ to $C^{0}([0,T],X^{\sigma})$ by the formula
\begin{align*}
\left[e^{sA}V\right](t)=e^{sA}V(t),\quad \text{for $V\in C^{0}([0,T],X^{\sigma})$.}
\end{align*}
Given $H\in L^{1}\big(I,C^{0}([0,T],X^{\sigma})\big)$, we~say that $\xi\in C^{0}\big(I,C^{0}([0,T],X^{\sigma})\big)$ satisfies
\begin{align}\label{appendix:ode:ode-banach-1}
\left\{\begin{aligned}
\dfrac{d}{ds}\xi(s)&=A \xi(s)+H(s), && s\in I, \\
\xi(s_0)&=\xi_0, &&
\end{aligned}\right.
\end{align}
with $\xi_0\in C^{0}([0,T],X^{\sigma})$, if~it satisfies
\begin{align}\label{appendix:ode:ode-banach-1-duhamel}
\xi(s)&=e^{(s-s_0)A}\xi_0+\int_{s_0}^{s}e^{(s-w)A}H(w)dw,\ \forall s\in I,
\end{align}
with equality in $C^{0}([0,T],X^{\sigma})$.
\begin{lemma}
\label{lmDuhamelCauchy}
If $H\in C^0\left(I,C^{0}([0,T],\P_{n}X^{\sigma})\right)$ and $\xi\in C^{0}\left(I,C^{0}([0,T],\P_{n}X^{\sigma})\right)$ for some $n\in\N$ and satisfies
$ \frac{d}{ds}\xi(s)=A \xi(s)+H(s)$ in the previous sense \eqref{appendix:ode:ode-banach-1-duhamel}, then $\xi\in C^{1}\left(I,C^{0}([0,T],\P_{n}X^{\sigma})\right)$ and is a classical solution.
\end{lemma}
\begin{proof}
Since $A$ is a bounded operator on $C^{0}([0,T],\P_{n}X^{\sigma})$, $\frac{d}{dt}e^{tA}\xi_0=A e^{tA}\xi_0$ in the classical sense of $C^1$ functions with value in $C^{0}([0,T],\P_{n}X^{\sigma})$. In~particular, Duhamel's formula \eqref{appendix:ode:ode-banach-1-duhamel} leads to the claimed result by usual arguments of semigroup theory.
\end{proof}

\begin{lemma}\label{lmtranslat}
For $T_1<T_2$, let us consider $T\in (0, T_2-T_1)$, $\eta\in (0, T_2-T-T_1)$ and $I:=[T_1-\eta, T_2-T-\eta]$. Let $G\in C^{0}([T_{1},T_{2}],X^{\sigma})$ and assume that $V\in C^{0}([T_{1},T_{2}],X^{\sigma})$ is a mild solution of
\begin{align*}
\left\{\begin{aligned}
\dfrac{d}{dt}V(t)&=A V(t)+G(t) &&\text{ for }\ t\in [T_1, T_2],\\
V(T_1)&=V_0. &
\end{aligned}\right.
\end{align*}
If we define $\xi,H\in C^{0}\bigl(I, C^{0}([0,T],X^{\sigma})\bigr)$ by $\xi(s)=V^s$ and $H(s)=G^s$ with
\begin{align*}
\xi(s)(t)=V^{s}(t)=V(t+s+\eta)\ \text{ and }\ H(s)(t)=G^{s}(t)=G(t+s+\eta),
\end{align*}
for all $s\in I, t\in [0,T]$, then for any $s_0\in I$, $\xi$ is solution, in the sense of \eqref{appendix:ode:ode-banach-1-duhamel}, of
\begin{align*}
\left\{\begin{aligned}
\dfrac{d}{ds}\xi(s)&=A \xi(s)+H(s), && s\in I, \\
\xi(s_0)&=\xi_0, &
\end{aligned}\right.
\end{align*}
with $\xi_0=V^{s_0}=V(\cdot+s_0+\eta)$.
\end{lemma}
\begin{proof}
By Duhamel's formula, for all $t\in [T_1, T_2]$ we have
\begin{align*}
V(t)&=e^{(t-T_1)A}V_{0}+\int_{T_1}^{t}e^{(t-\tau)A}G(\tau)d\tau,
\end{align*}
with $V(T_1)=V_0$. Pick $s_0\in I$. First observe that, for any $t\in [0, T]$,
\begin{align*}
V(t+s_0+\eta)=e^{(t+s_0+\eta-T_1)A}V_{0}+\int_{T_1}^{t+s_0+\eta}e^{(t+s_0+\eta-\tau)A}G(\tau)d\tau.
\end{align*}
Then for $s\in I$ and $t\in [0, T]$,
\begin{align*}
V^{s}(t)&=V(t+s+\eta)=e^{(t+s+\eta-T_1)A}V_{0}+\int_{T_1}^{t+s+\eta}e^{(t+s+\eta-\tau)A}G(\tau)d\tau\\
&=e^{(s-s_0)A}V(t+s_0+\eta)+\int_{s_0}^{s}e^{(s-w)A}G(t+w+\eta)dw\\
&=e^{(s-s_0)A}V^{s_0}(t)+\int_{s_0}^{s}e^{(s-w)A}G^{w}(t)dw.
\end{align*}
So, since this is true for any $t\in [0,T]$, it~gives
\begin{align*}
V^{s}=e^{(s-s_0)A}V^{s_0}+\int_{s_0}^{s}e^{(s-w)A}G^{w}dw,\quad \forall s\in I.
\end{align*}
By hypothesis, this equality holds in $C^0([0, T], X^\sigma)$, and is exactly \eqref{appendix:ode:ode-banach-1-duhamel}, as we wanted to prove.
\end{proof}

\subsection{Geometric fact}\label{s:geom}
In this section, we~describe briefly the compressed cotangent bundle and prove a geometric lemma that has been used several times in the article. We~refer to Melrose-Sjöstrand \cite{MS:78} Hörmander \cite[\S 18.3 \& 24.3]{H:07} for more precisions and \cite[\S 2.2]{BL:01} in the more specific context of the wave equation.

Let $\M$ be a smooth compact Riemannian manifold of dimension $d$ with boundary. Denote $^bT\M$ the bundle of rank $d$ whose sections are the vector fields tangent to $\partial \M$, by~$^bT^*\M$ the dual bundle (Melrose's compressed cotangent bundle), and by $\mathfrak{j}:T^*\M\rightarrow~ ^bT^*\M$ the canonical map. This is the restriction map, dual to the embedding $^bT\M\hookrightarrow T\M$. Its image can be identified with $T^* (\text{Int}\M) \sqcup T^* \partial\M$ with an appropriate topology.

We set $^bS^*\M=(^bT^*\M\setminus \M_0)/\R_+^*$ to be the cosphere bundle of $^bT^*\M$ the compressed cotangent bundle. Here $\M_0\approx \M$ is the zero section. The map $\mathfrak{j}$ can be defined on the quotient and allows to define $B:=\mathfrak{j}(S^*\M)\subset ~^bS^*\M$. It~can be identified with the image of $S^*\M$ by the continuous map $\mathfrak{j}$ and is therefore a compact space when it is equipped with the natural topology of vector bundle inherited from $~^bS^*\M$.

The bicharacteristic flow is usually defined for non-elliptic operators, but the link with generalized geodesics can, for instance, be made as follows, see Lebeau \cite[\S A.3]{L:96}. Let $P=\partial_t^2-\Delta_g$ be the wave operator defined on the manifold with boundary $\X=\R_t\times \M$ and $p=|\xi_x|_g^2-\xi_t^2$, the principal symbol of $P$, which is well-defined on $T^*\X\approx T^*\R_t\times T^*\M$. The set $p^{-1}(0)$ is conical and therefore well-defined in $T^*\X$ and $S^*\X$. We~denote $Z=\mathfrak{j}(p^{-1}(0))\subset~ ^bT^*\X$ and $SZ=(Z\setminus \X_0)/\R_+^*\subset ~ ^bS^*\X$. Note that $SZ$ has actually two connected components corresponding to $\xi_t>0$ and $\xi_t<0$, where~$\xi_t$ is the variable dual to $t$, that we denote $Z^+$ and $Z^-$. We~can see that $SZ^+$ can be identified with $(Z\setminus \X_0)\cap \{\xi_t=1/2\}$ and to $\R_t\times~ \mathfrak{j}(S^*\M)$. In~that context, since~$\xi_t$ is invariant by the flow, if~we denote $G$ the bicharacteristic flow of Melrose-Sjöstrand, we~see that it can be written $
G(s)(t,x,\xi)=(t+s,\phi_s (x,\xi))$, where $\phi_s$ is a well-defined flow on $B=\mathfrak{j}(S^*\M)\subset ~^bS^*\M$, the generalized geodesic flow.

We denote $\pi$ the natural projection from $^bS^*\M$ to $\M$. Both are continuous with the natural topologies given. For $\phi_t$, this is a consequence of the continuity of the bicharacteristic flow $G$, see \cite[Lem.\,3.31]{MS:78}.
\begin{lemma}\label{lemma:gcc-smaller-subset}
Suppose that $(\omega, T)$ satisfies \ref{assumGCC}. Then there exist $\chi\in C^{\infty}_c(\omega)$ with non negative values such that
\begin{itemize}
\item there exists an open set $\widetilde{\omega}\Subset \omega$ and a time $\widetilde{T}\in(0,T)$ such that $(\widetilde{\omega}, \widetilde{T})$ satisfies \ref{assumGCC},
\item $\chi=1$ on $\widetilde{\omega}$,
\item $\partial_{\vec n} \chi=0$ on $\partial \M$ where $\partial_{\vec n}$ is the normal derivative to the boundary.
\end{itemize}
\end{lemma}
\begin{proof}
Let $\rho=(x, \xi)\in B$. By~\ref{assumGCC}, there exists $t=t(\rho)\in (0, T)$ such that $x_t:=\pi\circ\phi_t(\rho)\in \omega$.

We first assume $x_t\in \partial \M$, the case $x_t\in \text{Int}(\M)$ being simpler.
In a sufficiently small neighborhood of $x_t$, we~can find some geodesic normal coordinates $(x_1,x')\in [0,2\eta)\times B_{\R^{d-1}}(0,2\eta)$ satisfying the following properties.
The point $x_t$ is $(0,0)$ in these coordinates, $\partial \M=\{x_1=0\}$, $\M=\{x_1\geq 0\}$ and the metric $g$ can be written in a diagonal type form $\begin{Smallpmatrix}
1 & 0 \\
0 &g'(x_1,x')
\end{Smallpmatrix}$ where $g'$ is a metric on $\R^{d-1}$ depending smoothly on $(x_1,x')$. In~particular, in~these coordinates, the normal vector at the boundary pointing inside is $\vec n=\sfrac{\partial}{\partial x_1}$. By assumption, $x_t\in \omega$, so up to diminishing $\eta$, we~can suppose that $[0,2\eta)\times B_{\R^{d-1}}(0,2\eta)\subset \omega$.

Let $\varphi\in C^{\infty}_c((-2,2),[0,1])$ such that $\varphi=1$ on $(-1,1)$. In~these coordinates, we~define $\chi_{\rho}(x_{1},x')= \varphi(x_{1}/\eta)\varphi(|x'|/\eta)$. We~verify that $\chi_{\rho}$ satisfies $\chi_{\rho}\geq 0$, $\chi_{\rho}=1$ in~$[0,\eta)\times B_{\R^{d-1}}(0,\eta)$ which contains $(0,0)$, $\partial_{x_1}\chi_{\rho}=0$ on $\{x_1=0\}$ and $\text{Supp}(\chi_{\rho})\subset \omega$. Independently of the coordinates in $\M$, the properties of $\chi_{\rho}$ can be written
\begin{itemize}
\item $\chi_{\rho}\geq 0$,
\item $\chi_{\rho}=1$ in an open set $\omega_{\rho}\subset \M$ (open for the topology of a manifold with boundary) which contains $x_t$, and we can select another open set $\widetilde{\omega}_{\rho}\Subset \omega_{\rho}$ with $x_t\in \widetilde{\omega}_{\rho}$,
\item $\partial_{\vec n}\chi_{\rho}=0$ on $\partial \M$,
\item $\chi_{\rho}\in C^{\infty}_c(\omega)$.
\end{itemize}
In the case $x_t\in \Int{\M}$, a cutoff function $\chi_{\rho}$ can be found with the same properties.

Since $x_t:=\pi\circ\phi_t(\rho)\in \widetilde{\omega}_{\rho}$ with $\widetilde{\omega}_{\rho}$ open and for fixed $t$, $\pi\circ\phi_t$ is continuous from~$B$ to $\M$, we~can find a neighborhood $\mc{V}_{\rho}\subset B$ of $\rho$ such that
\begin{align}
\label{inclusflow}
\pi\circ\phi_t(\mc{V}_\rho)\subset\widetilde{\omega}_{\rho}.
\end{align}

The open sets $\mc{V}_{\rho}$ form an open covering of $B$ by such neighborhoods. The compactness of $^bS^*\M$ allows to extract a finite subcovering of them $\mc{V}_{\rho_1},\ldots, \mc{V}_{\rho_k}$ such that
\begin{align}
\label{finitecover}B=\bigcup_{i=1}^{k}\mc{V}_{\rho_{i}}.
\end{align}

We define $\widetilde{\chi}=\sum_{i=1}^{k}\chi_{\rho_{i}}$, $\widetilde{\omega}=\bigcup_{i=1}^{k}\widetilde{\omega}_{\rho_i}$, $F=\bigcup_{i=1}^{k}\overline{\widetilde{\omega}_{\rho_i}}$ and $W=\bigcup_{i=1}^{k}\omega_{\rho_i}$. Note that we have $ \widetilde{\omega}\subset F\subset W\subset \omega$. For $\widetilde{\chi}$, we~have the following properties \begin{itemize}
\item $\chi\geq 0$ on $\M$,
\item $\widetilde{\chi}\geq 1$ in the open set $W$,
\item $\partial_{\vec n}\widetilde{\chi}=0$ on $\partial \M$,
\item $\widetilde{\chi}\in C^{\infty}_c(\omega)$.
\end{itemize}
We state the following claim that we will prove later.

\begin{claim} Let $F\subset W\subset \M$ with $F$ closed and $W$ open. Then there exists $h\in C^{\infty}(\M)$ such that $h\geq 0$ on $\M$, $h=0$ on $F$, $h\geq 1$ on $\M\setminus W$ and $\partial_{\vec n}h=0$ on $\partial \M$.
\end{claim}

We now define $\chi=\spfrac{\widetilde{\chi}}{\widetilde{\chi}+h}$. We have $\widetilde{\chi}+h\geq 1$ on $\M$, so $\chi\in C^{\infty}(\M)$. It~satisfies $\partial_{\vec n}\chi=0$ on $\partial \M$ by composition since it is the case for $\widetilde{\chi}$ and $h$. Moreover, since $h=0$ on $F$, we~have $\chi=1$ on $F$.

We now check that $(\widetilde{\omega},\widetilde{T})$ satisfies \ref{assumGCC} for some $\widetilde{T}\in (\max\{t(\rho_j),j=1,\dots, k\}, T)$.
Indeed, let $\rho \in B$. By~the finite covering property \eqref{finitecover}, there exists $j\in \llbracket 1,k \rrbracket $ such that $\rho\in \mc{V}_{\rho_j}$. \eqref{inclusflow} implies then $\pi\circ\phi_{t(\rho_{j})}(\rho)\in \widetilde{\omega}_{\rho_{j}}$ and then $\pi\circ\phi_{t(\rho_{j})}(\rho)\in\widetilde{\omega}$ as expected.

The proof is now complete except for the proof of the claim.
Let $x\in \M\setminus W$. In~particular, $x\in \M \setminus F$ which is open. By~following the same method as in the first part of the proof, we~can construct $h_{x} \in C^{\infty}_c(\M \setminus F)$ such that $h_{x}\geq 0$, $h_{x}=1$ in an open set $\omega_{x}\subset \M\setminus F$ which contains $x$. We~can also assume $\partial_{\vec n}h_{x}=0$ on~$\partial \M$. The $\omega_{x}$ form a covering of the compact set $\M\setminus W$. So, we~can select a finite covering $\M\setminus W=\bigcup_{i=1}^{k} \omega_{x_{i}}$ and define $h=\sum_{i=1}^{k}h_{i}$ which satisfies the expected properties.
\end{proof}
With similar arguments, we~can also prove the following result.
\begin{lemma}\label{lemma:omomtilde-smaller-subset}
Suppose that $\omega$ and $\widetilde{\omega}$ are two open subsets of $\M$ such that $\omega \Subset \widetilde{\omega}$. Then there exists $\chi\in C^{\infty}_c(\widetilde{\omega})$ with non negative values such that
\begin{itemize}
\item $\chi=1$ on $\omega$,
\item $\partial_{\vec n} \chi=0$ on $\partial \M$.
\end{itemize}
\end{lemma}
\begin{proof}
Since $\omega \Subset \widetilde{\omega}$ where both sets are open, we~can find $\omega_1$ open such that $\omega \Subset \omega_1 \Subset \widetilde{\omega}$. We~apply the claim with $F=\M\setminus \widetilde{\omega}$, $W=\M\setminus \overline{\omega_1}$ to get $h\in C^{\infty}_c(\widetilde{\omega})$ such that $h\geq 1$ on $\overline{\omega_1}$. If we apply again the claim with $F=\overline{\omega}$, $W=\omega_1$ to get $m\in C^{\infty}(\M)$ such that $m=0$ on $\overline{\omega}$ and $m\geq 1$ on $\M\setminus\omega_1$. In~both cases, the functions are nonnegative and we have $\partial_{\vec n} h=\partial_{\vec n} m=0$ on $\partial \M$. The function $\chi=\spfrac{h}{m+h}$ is regular and has the desired properties.
\end{proof}

\backmatter
\bibliographystyle{jepplain+eid}
\bibliography{laurent-loyola}
\end{document}